\chapter{文本生成与转换}
\label{chap:nlp-text-generation}

前一章已经能够判断通知的类别、找到相关材料，并把重复报道组织到一起。现在需要交付的不再是标签或材料列表，而是一段可以直接阅读的文字：把通知译给外国来宾，压缩成日历说明，从活动记录表生成提醒，或改写成容易理解的参加指南。这些输出都可能写得流畅，却犯下不同的错误。译文不能把“校外来宾”变成所有参加者；摘要是否必须保留报名截止条件，取决于它的用途；改写若删掉“须”，就改变了通知的要求。本章研究输出为文本、且输出须对给定条件负责的任务：怎样依据材料形成新文本，以及怎样控制和检查其中的信息变化。

这些任务的输入、输出和允许的改动各不相同，却反复遇到同一组挑战：从条件中选出该说的内容；准确复制名称与数字；处理语序重排与输入输出之间的对齐；不补造无依据的内容；避免重复与退化；在没有成对数据时学习转换；处理超出上下文范围的长输入。方法的演进大致沿一条主线：模板与规则流水线保证可控但难以迁移；统计模型（短语翻译、混淆集打分）从数据中估计转换；注意力编码器—解码器加复制机制把各任务统一为条件序列生成；去噪预训练的编码器—解码器改善了低资源下的表达；大语言模型则以提示、微调与人类反馈完成大部分任务。各任务节沿这条主线说明每种方法补上了前一方法的哪个缺口、又带来了什么代价，并在节末给出“方法／解决的问题／核心机制／新代价”对照表；七项挑战与各节的对应关系见表~\ref{tab:nlp-generation-challenges}。

本章建立在前面几章之上。第~\ref{chap:nlp-large-language-models}章给出了模型结构、交叉注意力和T5、BART的去噪预训练；第~\ref{chap:nlp-supervised-finetuning}章给出了教师强制与掩码损失；第~\ref{chap:nlp-preference-alignment}章给出了策略梯度与基于人类反馈的强化学习；第~\ref{chap:nlp-model-inference}章给出了束搜索、截断采样、重排、最小贝叶斯风险解码和文法约束解码；第~\ref{chap:nlp-text-analysis}章在字符级任务中引入了复制机制。本章不再重复这些机制，只写它们接入具体任务后需要补充的结构与目标。第~\ref{sec:nlp-generation-interface}节先集中处理各任务共用的问题：条件怎样读取、罕见词怎样复制、训练与生成为什么不一致，以及怎样核验结果。随后按内容从哪里来、什么必须保留划分任务：数据到文本必须忠实于给定记录，开放写作允许虚构但要连贯且可控，文本编辑的输入与输出大部分相同，摘要要按用途选择并压缩内容，翻译要把全部意义带到另一种语言。最后两节处理跨任务的硬约束与评价：约束解码保证必须出现的术语和格式，评价指标则各自检查质量的一个方面。

以下通知、资料、分数和轨迹均为教学构造，不是模型实测。共同通知于2025年4月15日（北京时间）发布，活动编号为\code{ML-0418}：
\begin{quote}
机器学习读书会原定于4月17日举行，现改至4月18日（周五）14:00，地点为海棠楼201。
校内师生可直接参加，校外来宾须在4月16日18:00前提交报名信息。
\end{quote}
材料没有说明主讲人、费用、人数或时长，生成器不能把这些空白自动补成事实。

\section{条件生成的共同问题}
\label{sec:nlp-generation-interface}

翻译、摘要、数据描述和改写的输入各不相同，但都要求模型在给定条件下写出一段文本，并为其中的每条信息负责。它们因此共用几项基本问题：条件和目标怎样定义，模型怎样在每一步读取任意长度的条件，名称和数字这类罕见词元怎样原样保留，按参考前缀训练的模型在自己生成的前缀上为什么会出错，以及哪些困难反复出现在所有任务中。本节依次回答这些问题，并把全章共同挑战整理成一张索引表，最后给出各任务共用的结果核验流程；后续各任务节回指这里，只补任务特有的结构与目标。

\subsection{输入内容、要求与目标文本}
\label{sec:nlp-generation-training}

把输入条件记为$c=(x,u,b)$：$x$是待处理材料，$u$是任务要求，$b$是允许使用的背景。例如$x$为章首通知，$u$可以是“译成英语，保留全部参加条件”，也可以是“为校外来宾生成两句参加指南”；$b$可以给出年份、时区和经确认的名称译法，但不能包含测试参考答案。目标文本$y_{1:T}$采用固定分词器，其中$y_T=\mathrm{EOS}$是结束符；起始符$\mathrm{BOS}$只用于启动生成，不是要交付的文字。

同一材料在不同要求下有不同的合法目标。英文译文“External guests must submit their registration information before 18:00 on April 16”需要保留截止条件；只供日历显示的“读书会改至4月18日14:00，海棠楼201”可以不承担参加指南的职责；面向来宾的易读改写则应写明“校外来宾请在4月16日18:00前提交报名信息”。多参考文本表示同一要求下的表达差异，不同用途的答案不能混在一起充当参考。

在当前材料与要求下，第~\ref{chap:nlp-large-language-models}章的序列概率写成链式分解
\begin{equation}
  p_{\bm{\theta}}(y_{1:T}\mid c)
  =\prod_{t=1}^{T}p_{\bm{\theta}}(y_t\mid c,y_{<t}),
  \label{eq:nlp-generation-factorization}
\end{equation}
其中$\bm{\theta}$是可学习参数，$y_{<t}$是位置$t$之前的目标词元。对$n$个独立训练样本的对数似然取负，得到\term{条件最大似然}（conditional maximum likelihood）的最小化目标
\begin{equation}
  L(\bm{\theta})
  =-\sum_{i=1}^{n}\sum_{t=1}^{T_i}
  \log p_{\bm{\theta}}(y_{i,t}\mid c_i,y_{i,<t}).
  \label{eq:nlp-generation-nll}
\end{equation}
它让正确词元在同一位置的候选中获得更大概率。先把每条样本的损失除以自身长度再平均，会让各句等权；把全部有效目标词元的损失一起平均，则长句贡献更多。训练记录应写清采用哪一层归一化。

训练时每个位置都以参考前缀为条件，这就是第~\ref{sec:nlp-sft-masked-loss}小节的\term{教师强制}（teacher forcing）。设简化的目标词元为“周五／举行／EOS”，解码输入是“BOS／周五／举行”，预测标签是“周五／举行／EOS”。若三个标签的条件概率为$0.6,0.5,0.8$（教学数值），样本概率为$0.24$，负对数似然为$-\log0.24\approx1.4271$，每个位置平均$0.4757$。结束符参与监督，模型才能学会何时停止。

条件怎样交给模型有两种组织方式。编码器—解码器把$c$交给编码器，解码器读取移位后的目标，只在目标位置计算损失。仅解码器把“条件／分隔符／目标”拼成一条因果序列，条件位置和补齐位置的损失掩码为零，目标及其EOS位置为一；掩掉条件的损失并不掩掉对条件的注意，目标仍须读取条件。评价时模型只能看到$x,u,b$、指定的词典与资料库以及自身生成历史；选择提示、检索资料或阈值时用过的测试参考，不能再充当独立测试。

\subsection{编码器—解码器与交叉注意力}
\label{sec:nlp-generation-encoder-decoder}

条件生成的第一个结构问题是：解码器每写一个词元，怎样读取长度任意的条件。最早的神经序列到序列模型用循环编码器读完源句，把最后的隐状态作为一个定长向量交给循环解码器\scite{sutskever2014seq2seq,cho2014}。所有源端信息都要挤进这一个向量，句子越长，解码器越难从中取回某个具体的日期或名称。Bahdanau等的比较显示，这类定长编码模型的译文质量随句长增加而明显下降\scite{bahdanau2015}。

注意力改变的是读取方式：解码器在每一步以当前状态为查询，对全部编码器状态打分并加权求和，各步可以读取不同的源位置。设编码器状态为$\bm h_1,\ldots,\bm h_S\in\mathbb R^{d}$，解码器在第$t$步的状态为$\bm s_t$，加性注意力（additive attention）计算
\begin{equation}
  \begin{gathered}
  e_{t,i}=\bm v^{\top}\tanh\bigl(\bm W_h\bm h_i+\bm W_s\bm s_t+\bm b\bigr),\\
  a_{t,i}=\frac{\exp e_{t,i}}{\sum_{j=1}^{S}\exp e_{t,j}},\qquad
  \bm h^{*}_t=\sum_{i=1}^{S}a_{t,i}\bm h_i .
  \end{gathered}
  \label{eq:nlp-generation-bahdanau}
\end{equation}
$e_{t,i}$是第$t$步对源位置$i$的打分，$a_{t,i}$是归一化后的注意力权重，上下文向量$\bm h^{*}_t$随后与$\bm s_t$一起决定输出分布。本章后文只用到这一机制的一点：每一步按当前查询重新加权全部源状态，权重可以作为“正在读哪里”的分布直接使用。循环网络、门控单元以及注意力循环网络的完整结构可参考\BookChapterRef{chap:recurrent-networks}{卷五的循环神经网络章}。

Transformer把同一思想写成交叉注意力，其单头公式、掩码规则和键值缓存已在第~\ref{sec:nlp-llm-objective-cross-attention}小节式~\eqref{eq:nlp-llm-cross-attention}给出：查询来自解码器，键和值来自编码器输出。这里补上条件生成训练时一个批次的张量组织与解码层的计算次序。设批大小为$B$，源长为$S$，目标长为$T$，宽度为$d$，头数为$H$，每头维度$d_h=d/H$，词表大小为$V$。
\begin{itemize}
  \item 编码器输入为$B\times S$的词元编号与同形状的填充掩码，嵌入加位置后为$B\times S\times d$，经$L_{\mathrm e}$层双向自注意力与前馈网络得到$\bm H_{\mathrm{enc}}\in\mathbb R^{B\times S\times d}$。
  \item 解码器输入是右移一位的目标$B\times T$，嵌入后为$\bm H^{(0)}\in\mathbb R^{B\times T\times d}$；标签是未移位的目标$B\times T$。
  \item 每个解码层的自注意力权重为$B\times H\times T\times T$并带因果掩码；交叉注意力权重为$B\times H\times T\times S$，只对源端填充列加掩码。
  \item 最后一层输出$B\times T\times d$乘以输出矩阵得到$B\times T\times V$的对数几率，在标签位置计算式~\eqref{eq:nlp-generation-nll}。
\end{itemize}

按预归一化写法，第$\ell$个解码层依次执行三个子层：
\begin{equation}
  \begin{aligned}
    \bm A&=\bm H^{(\ell-1)}+\operatorname{MaskedSelfAttn}\bigl(\operatorname{Norm}(\bm H^{(\ell-1)})\bigr),\\
    \bm C&=\bm A+\operatorname{CrossAttn}\bigl(\operatorname{Norm}(\bm A),\,\bm H_{\mathrm{enc}}\bigr),\\
    \bm H^{(\ell)}&=\bm C+\operatorname{FFN}\bigl(\operatorname{Norm}(\bm C)\bigr).
  \end{aligned}
  \label{eq:nlp-generation-decoder-layer}
\end{equation}
第一行让目标位置读取已写出的前缀，第二行读取条件，第三行逐位置变换。三个子层的输出都保持$B\times T\times d$，残差连接要求形状不变。

用一个两头的小例子看交叉注意力读到了什么。设源端三个位置为“改至／4月18日／海棠楼201”，解码器正在写下一个词元，两头的缩放后得分（即$\bm q\bm k^{\top}/\sqrt{d_h}$）分别为$(0,2,0)$与$(0,0,1)$（教学数值）。softmax后，第一头权重为
\[
  \Bigl(\tfrac{1}{e^{2}+2},\tfrac{e^{2}}{e^{2}+2},\tfrac{1}{e^{2}+2}\Bigr)\approx(0.107,0.787,0.107),
\]
第二头为$(1,1,e)/(2+e)\approx(0.212,0.212,0.576)$。若第一头的值向量依次为$(1,0),(0,1),(0,0)$，第二头为$(0,0),(1,0),(0,1)$，两头输出分别为$(0.107,0.787)$和$(0.212,0.576)$，拼接成四维向量$(0.107,0.787,0.212,0.576)$后再乘输出投影。第一头主要读取日期，第二头主要读取地点，同一步的两个头可以从不同位置取回信息；权重全部为正，所以每头的输出仍混入少量其他位置的值。

推断时，每个样本的交叉注意力键值共$2L_{\mathrm d}\times S\times d$个数，在第一步算好后整个生成过程都可复用（第~\ref{sec:nlp-llm-objective-cross-attention}小节）。仅解码器把条件拼成前缀时，条件的键值同样只算一次，但条件位置只能看左侧上下文，条件与目标共用同一组参数；两种结构的取舍见第~\ref{sec:nlp-llm-objective-architecture}小节。对本章任务而言，编码器—解码器的优点是源端双向编码、输入与输出分工明确，在翻译、摘要等输入远长于输出的任务中仍常用；代价是条件更新后需要重新编码，难以像因果前缀那样追加多轮内容。

\subsection{复制与指针机制}
\label{sec:nlp-generation-copy-pointer}

注意力让解码器读得到“海棠楼201”，却不保证它写得出。输出层只能在固定词表上分配概率，房间号、活动编号、人名和罕见数字往往不在词表中，或在训练中出现次数太少，模型倾向于用常见词替换它们；而这些恰是条件生成最需要准确保留的内容。子词切分（第~\ref{sec:nlp-llm-bpe}小节）让任何字符串都能拼出来，却没有让模型更倾向于照抄输入。复制机制直接给“从输入中取一个位置”分配概率。

最直接的做法是\term{指针网络}（pointer network）：输出不是词表中的词，而是输入的位置\scite{nlp-generation-vinyals2015pointer}。第~\ref{sec:nlp-analysis-pointer-ordering}小节已用它逐步选出下一个句子来排序段落；一般形式是第$t$步对源位置$i$打分后在位置上归一化，
\begin{equation}
  u_{t,i}=\bm v^{\top}\tanh\bigl(\bm W_1\bm h_i+\bm W_2\bm s_t\bigr),\qquad
  p(z_t=i\mid z_{<t},x)=\frac{\exp u_{t,i}}{\sum_{j=1}^{S}\exp u_{t,j}} .
  \label{eq:nlp-generation-pointer-network}
\end{equation}
它与式~\eqref{eq:nlp-generation-bahdanau}的打分相同，区别在于注意力权重不再用于汇聚上下文，而是直接作为输出分布；输出空间的大小随输入长度变化。指针网络适合输出必为输入元素的问题，例如排序、选择记录或抽取片段；第~\ref{sec:nlp-generation-content-planning}小节的内容规划就用它在记录集合上逐步选择。它不能输出输入中没有的词。

条件生成的输出大多是新写的句子夹着照抄的名称，需要在生成与复制之间切换。\term{指针生成网络}（pointer-generator network）用一个生成门把两种分布混合起来\scite{nlp-intro-see2017}。设注意力权重$a_{t,i}$与上下文向量$\bm h^{*}_t$按式~\eqref{eq:nlp-generation-bahdanau}计算，词表分布$P_{\mathrm{vocab}}=\operatorname{softmax}(\bm W_o[\bm s_t;\bm h^{*}_t]+\bm b_o)$，则
\begin{equation}
  \begin{aligned}
    p_{\mathrm{gen},t}&=\sigma\bigl(\bm w_h^{\top}\bm h^{*}_t+\bm w_s^{\top}\bm s_t+\bm w_x^{\top}\bm x_t+b\bigr),\\
    p_t(w)&=p_{\mathrm{gen},t}\,P_{\mathrm{vocab}}(w)+\bigl(1-p_{\mathrm{gen},t}\bigr)\sum_{i:\,x_i=w}a_{t,i}.
  \end{aligned}
  \label{eq:nlp-generation-pointer-generator}
\end{equation}
按原文记号，$\bm x_t$是第$t$步解码器输入（即上一目标词元）的嵌入，不要与源端词元$x_i$混淆；$\sigma$是逻辑函数，$p_{\mathrm{gen},t}\in(0,1)$是这一步选择“生成”的概率。第二行的复制项把注意力权重当作复制分布，对源端所有等于$w$的位置求和。第~\ref{chap:nlp-text-analysis}章式~\eqref{eq:nlp-analysis-copy}已在字符级词形还原中给出同一形式及重复字符求和的算例，这里关注它在词级生成中的两处实现细节。

第一处是扩展词表。每个样本把源端出现而固定词表中没有的词临时加入，得到$V_{\mathrm{ext}}=V\cup\{x_i:x_i\notin V\}$，输出分布定义在$V_{\mathrm{ext}}$上。$w\notin V$时$P_{\mathrm{vocab}}(w)=0$，只能靠复制；$w$不在源端时复制项为$0$，只能靠生成。训练标签若是源端未登录词，就用它在扩展词表中的临时编号计算损失；若既不在词表也不在源端，仍记为未知词。第二处是门的输入。$p_{\mathrm{gen},t}$读取上下文、解码器状态和上一词元，模型因此能学到“刚写完‘地点为’之后更可能复制”这类规律，门值本身没有额外监督，随式~\eqref{eq:nlp-generation-nll}一起训练。

图~\ref{fig:nlp-generation-pointer-generator}画出第$t$步的计算。设上一词元为“海棠楼”，对源端“读书会／改至／海棠楼／201”的注意力为$(0.05,0.10,0.25,0.60)$，$p_{\mathrm{gen},t}=0.3$，而“201”不在固定词表中（教学数值）。则$p_t(\text{201})=0.7\times0.60=0.42$；固定词表中最可能的候选“教室”若有$P_{\mathrm{vocab}}=0.5$，混合后只有$0.3\times0.5=0.15$。门值偏向复制时，注意力集中的位置直接成为输出；同一词在源端多处出现时，各处权重相加。

\begin{figure*}[htbp]
  \centering
  % 指针生成网络在解码第t步的计算：编码器状态、注意力分布、上下文向量、生成门与混合输出分布。
% 注意力数值为教学示例，与正文算例一致。
\begin{tikzpicture}[
  x=1mm,y=1mm,figure picture,
  cell/.style={figure node,minimum width=13mm,minimum height=6mm,
    inner sep=1.5pt,font=\figurefont\scriptsize},
  tok/.style={font=\figurefont\scriptsize},
  flow/.style={figure flow},
  note/.style={font=\figurefont\scriptsize,text=FigureMuted},
  row label/.style={font=\figurefont\scriptsize,text=FigureMuted,anchor=east},
  rnn/.style={draw=FigureStroke,line width=.6pt,
    {Stealth[length=1.1mm]}-{Stealth[length=1.1mm]}},
  bar/.style={fill=FigureInputFill,draw=FigureInput,line width=.5pt},
  halo/.style={preaction={draw=white,line width=2.6pt}}
]
  % 源端词元、编码器状态与注意力分布
  \foreach \i/\w/\a in {1/读书会/0.05,2/改至/0.10,3/海棠楼/0.25,4/201/0.60}{
    \pgfmathsetmacro{\xx}{(\i-1)*18}
    \node[tok] (x\i) at (\xx,0) {\w};
    \node[cell,draw=FigureInput,fill=FigureInputFill] (h\i) at (\xx,10) {$\bm h_{\i}$};
    \draw[flow] (x\i) -- (h\i);
    \pgfmathsetmacro{\hh}{\a*20}
    \draw[bar] (\xx-3,20) rectangle (\xx+3,20+\hh);
    \node[note,anchor=south,inner sep=1pt] at (\xx,20+\hh) {\a};
  }
  \draw[rnn] (h1) -- (h2);
  \draw[rnn] (h2) -- (h3);
  \draw[rnn] (h3) -- (h4);
  \draw[figure link] (-4,20) -- (58,20);
  \node[row label] at (-6,0) {源词元};
  \node[row label] at (-6,10) {编码器状态};
  \node[row label] at (-6,25) {注意力$\bm a_{t}$};

  % 上下文向量
  \node[cell,draw=FigureModel,minimum width=34mm] (ctx) at (27,46)
    {上下文$\bm h^{*}_{t}=\sum_i a_{t,i}\bm h_i$};
  \draw[flow] (27,36) -- (ctx.south);

  % 解码器
  \node[tok] (yprev) at (88,0) {上一词元“海棠楼”};
  \node[cell,draw=FigureModel] (s) at (88,10) {$\bm s_t$};
  \draw[flow] (yprev) -- (s);
  \draw[figure link,densely dashed,-{Stealth[length=1.2mm]}] (s.west)
    -- node[note,above,sloped,pos=.45] {查询} (61,23);

  % 词表分布
  \node[cell,draw=FigureModel,minimum width=28mm] (pv) at (88,46)
    {词表分布$P_{\mathrm{vocab}}$};
  \draw[flow] (s.north) -- (pv.south);
  \draw[flow] (ctx.east) -- (pv.west);

  % 生成门
  \node[circle,draw=FigureLoss,fill=FigureLossFill,line width=.75pt,
    minimum size=7mm,inner sep=0pt,font=\figurefont\scriptsize] (g) at (58,62) {$\sigma$};
  \node[note,anchor=west] at (g.east) {生成门$p_{\mathrm{gen}}$};
  \draw[flow] (ctx.north east) -- (g);
  \draw[flow,halo] (s.north west) -- (g.south east);

  % 混合输出分布
  \node[cell,draw=FigureOutput,fill=FigureOutputFill,minimum width=72mm] (out) at (58,80)
    {$p(w)=p_{\mathrm{gen}}P_{\mathrm{vocab}}(w)+(1-p_{\mathrm{gen}})\sum_{i:x_i=w}a_{t,i}$};
  \draw[flow] (g.north) -- (g.north |- out.south);
  \draw[flow] (pv.north) -- node[note,right] {生成} (pv.north |- out.south);
  \draw[flow] (-4,22) -- node[note,left,pos=.55] {复制} (-4,80) -- (out.west);
\end{tikzpicture}

  \caption{指针生成网络在解码第$t$步的计算。下方是源端词元与双向编码器状态$\bm h_i$，条形为解码器状态$\bm s_t$作查询得到的注意力$\bm a_t$；注意力一方面汇聚成上下文$\bm h^{*}_t$，与$\bm s_t$一起给出固定词表上的分布，另一方面直接作为复制分布（左侧路径）。生成门$p_{\mathrm{gen}}$读取$\bm h^{*}_t$、$\bm s_t$与解码器输入嵌入（图中略去后者的连线），按式~\eqref{eq:nlp-generation-pointer-generator}混合两路分布。注意力数值为教学示例。}
  \label{fig:nlp-generation-pointer-generator}
\end{figure*}

复制机制还有另一种归一化方式。CopyNet不设显式的门，而是给生成模式和复制模式各算一组得分，再在“词表词与源端位置”的并集上统一做softmax\scite{nlp-analysis-gu2016copynet}。两种写法都让输出空间成为词表与源端词的并集；显式门便于观察每步偏向复制还是生成，共享归一化则让两种模式的得分直接竞争。指针生成网络在摘要中还配有覆盖机制，用累计注意力惩罚重复关注同一位置，见第~\ref{sec:nlp-generation-summarization}节。

复制解决的是“写得出”，不解决“该不该写”。模型可以准确复制“海棠楼201”，却把它绑定到另一场活动，或在不需要地点的日历摘要中照抄整句。复制概率高也不说明输出忠实，所以复制机制之后仍需要内容选择（第~\ref{sec:nlp-generation-data-to-text}节）和结果核验（第~\ref{sec:nlp-generation-recovery}小节）。后文的数据到文本、摘要和对话状态追踪都沿用式~\eqref{eq:nlp-generation-pointer-generator}，只说明各自复制的单位（记录值、源句词元或对话历史词元）和门的输入。

\subsection{训练与生成的不一致}
\label{sec:nlp-generation-exposure-bias}

式~\eqref{eq:nlp-generation-nll}在每个位置都以参考前缀$y^{\star}_{<t}$为条件，生成时模型只能以自己写出的前缀$\hat y_{<t}$为条件。两者在训练初期几乎处处相同，一旦生成早期出错，例如把“4月18日”写成“4月17日”，后续每一步的条件就是训练中从未出现过的前缀，模型在这些状态下的行为没有受过监督，错误可能继续扩散。这就是第~\ref{sec:nlp-sft-on-policy-distillation}小节定义的暴露偏差（exposure bias）：模型训练时只“暴露”在参考前缀上\scite{nlp-generation-ranzato2016mixer}。另一处不一致在目标上：逐词元交叉熵对每个位置独立计分，而任务关心整句的BLEU、ROUGE或事实正确性；一个词的同义替换在交叉熵下与严重错误同样受罚，在序列指标下却几乎不扣分。

第~\ref{chap:nlp-text-analysis}章的动态预言（第~\ref{sec:nlp-analysis-paradigm-transition}小节）和第~\ref{chap:nlp-supervised-finetuning}章在学生分布上的蒸馏处理的是同一类问题：让训练发生在模型自己会到达的状态上。条件生成中的对策分三类，依次放宽对参考前缀的依赖，并逐步把逐词元目标换成整句指标。

\term{计划采样}（scheduled sampling）在训练中逐步把参考前缀换成模型预测\scite{nlp-analysis-bengio2015scheduled}。第$k$个训练步中，解码第$t$个位置时以概率$\epsilon_k$输入参考词元$y^{\star}_{t-1}$，以概率$1-\epsilon_k$输入模型自己的预测$\hat y_{t-1}$（从当前分布采样或取最大值），标签仍是$y^{\star}_t$。$\epsilon_k$从$1$开始随训练步数下降，原文给出三种日程：
\begin{equation}
  \begin{gathered}
  \epsilon_k=\max(\epsilon_{\min},\,\kappa-\gamma k),\qquad
  \epsilon_k=\kappa^{k}\ (\kappa<1),\\
  \epsilon_k=\frac{\kappa}{\kappa+\exp(k/\kappa)}\ (\kappa\geq1),
  \end{gathered}
  \label{eq:nlp-generation-scheduled-sampling}
\end{equation}
分别是线性衰减、指数衰减和逆S形衰减。训练前期接近教师强制，保证模型先学会基本分布；后期更多地以自己的预测为条件，学会从偏离的前缀继续。

计划采样改动很小，但有两个代价。其一，模型偏离参考后，“下一个参考词元”未必仍是合理的标签：若模型漏写了“改至”，参考的下一个词“4月18日”接在漏写之后，标签与前缀已经错位。Huszár指出，这样得到的目标并不对应一个一致的概率估计准则\scite{nlp-generation-huszar2015scheduled}。其二，Transformer的教师强制能在一次前向中计算所有位置，计划采样需要先前向一次得到预测，再用混合后的前缀前向第二次，计算量约翻倍。

\term{序列级训练}（sequence-level training）直接优化整句奖励。设$r(y,y^{\star})$是候选$y$相对参考的句子级指标，例如句子BLEU或ROUGE，目标是最大化模型分布下的期望奖励
\begin{equation}
  J(\bm\theta)=\mathbb E_{y\sim p_{\bm\theta}(\cdot\mid c)}\bigl[r(y,y^{\star})\bigr],\qquad
  \nabla_{\bm\theta}J=\mathbb E_{y\sim p_{\bm\theta}}\Bigl[\bigl(r(y,y^{\star})-\beta\bigr)\nabla_{\bm\theta}\log p_{\bm\theta}(y\mid c)\Bigr].
  \label{eq:nlp-generation-sequence-reinforce}
\end{equation}
指标对词元不可导，梯度用REINFORCE估计：采样一条候选，按奖励减基线$\beta$加权它的对数概率梯度。这就是第~\ref{sec:nlp-align-mdp}小节陈述的策略梯度，把整句看作一个动作；基线不改变梯度期望，只降低方差。因为采样来自模型自身，训练前缀与生成前缀的分布一致，暴露偏差与指标不一致同时得到处理。

纯随机初始化的模型几乎采不到有奖励的句子，梯度方差也很大。MIXER先用交叉熵训练，再从句尾开始逐步扩大用REINFORCE训练的位置数：前$T-\Delta$个位置仍用教师强制，后$\Delta$个位置从模型采样并按奖励更新，$\Delta$随训练增大；基线由一个读取解码器状态的线性回归器估计\scite{nlp-generation-ranzato2016mixer}。自我批判序列训练（self-critical sequence training, SCST）换用更简单的基线：同一输入下贪心解码结果的奖励，$\beta=r(\hat y_{\mathrm{greedy}},y^{\star})$，采样句比贪心句好才被加强\scite{nlp-generation-rennie2017scst}。第~\ref{chap:nlp-preference-alignment}章的RLHF沿用同一梯度，只是把参考指标换成学习到的奖励模型，并加入KL约束。

序列级训练直接针对任务指标，代价是方差高、每步需要完整采样与打分，并且模型会利用指标的漏洞：只看$n$元重合的奖励会鼓励重复高频短语或删去难译部分。实际使用时通常与交叉熵损失加权混合。

REINFORCE每次只用一两个样本估计梯度。\term{最小风险训练}（minimum risk training, MRT）改为在一组候选上精确计算期望代价\scite{nlp-generation-shen2016mrt}。对每个训练样本，把参考译文与从当前模型采样的候选去重后合成候选集合$\mathcal S(c)$（原文每句采样100条），定义锐化后的候选分布与风险
\begin{equation}
  \begin{gathered}
  Q_\alpha(y\mid c)=\frac{p_{\bm\theta}(y\mid c)^{\alpha}}{\sum_{y'\in\mathcal S(c)}p_{\bm\theta}(y'\mid c)^{\alpha}},\\
  R(\bm\theta)=\sum_{(c,y^{\star})}\ \sum_{y\in\mathcal S(c)}Q_\alpha(y\mid c)\,\Delta(y,y^{\star}),
  \end{gathered}
  \label{eq:nlp-generation-mrt}
\end{equation}
其中$\Delta(y,y^{\star})$是代价，例如$1-\text{句子BLEU}$，$\alpha>0$调节分布的平滑程度。整句概率之间常相差几个数量级，$\alpha<1$把它们拉近，使集合中多个候选都能分到权重；$\alpha$越大，$Q_\alpha$越集中在概率最高的候选上。

对$\log Q_\alpha$求导再整理，梯度可以写成
\begin{equation}
  \nabla_{\bm\theta}R=\alpha\sum_{y\in\mathcal S(c)}Q_\alpha(y\mid c)\Bigl(\Delta(y,y^{\star})-\mathbb E_{Q_\alpha}[\Delta]\Bigr)\nabla_{\bm\theta}\log p_{\bm\theta}(y\mid c).
  \label{eq:nlp-generation-mrt-gradient}
\end{equation}
推导只用到$\nabla\log Q_\alpha(y)=\alpha\bigl(\nabla\log p_{\bm\theta}(y)-\mathbb E_{Q_\alpha}[\nabla\log p_{\bm\theta}]\bigr)$以及$\sum_yQ_\alpha(y)=1$。代价低于集合平均的候选概率被提高，高于平均的被压低。与式~\eqref{eq:nlp-generation-sequence-reinforce}相比，期望代价$\mathbb E_{Q_\alpha}[\Delta]$自动充当基线，且对整组候选求和而不是单样本估计，方差更低。

用三条候选核对式~\eqref{eq:nlp-generation-mrt-gradient}的方向。设$\alpha=1$，集合内归一化后的概率为$(0.5,0.3,0.2)$，代价为$(0.6,0.1,0.9)$（教学数值）。风险为$0.5\times0.6+0.3\times0.1+0.2\times0.9=0.51$，三项梯度权重$Q(\Delta-0.51)$为$(0.045,-0.123,0.078)$。梯度下降压低第一、三条候选，提高代价最低的第二条，尽管它不是当前概率最高的。若取$\alpha=0.5$，$Q$变为$(0.416,0.322,0.263)$，风险升为$0.518$，权重分布更平，三条候选都参与更新。

\Needspace{12\baselineskip}
\begin{algorithm}[最小风险训练的一步更新]
  \label{alg:nlp-generation-mrt}
  \begin{algorithmic}[1]
    \Require 用条件最大似然预训练的模型$p_{\bm\theta}$；小批次$\{(c_b,y^{\star}_b)\}_{b=1}^{B}$；候选数$N$；平滑系数$\alpha$；代价函数$\Delta$
    \Ensure 更新后的参数$\bm\theta$
    \For{$b=1,\ldots,B$}
      \State 从$p_{\bm\theta}(\cdot\mid c_b)$采样$N$条候选，去重后并入参考$y^{\star}_b$，得到$\mathcal S(c_b)$
      \State 计算每条候选的$\log p_{\bm\theta}(y\mid c_b)$与代价$\Delta(y,y^{\star}_b)$
      \State 按式~\eqref{eq:nlp-generation-mrt}在$\mathcal S(c_b)$内归一化得到$Q_\alpha$，求$R_b=\sum_yQ_\alpha(y\mid c_b)\Delta(y,y^{\star}_b)$
    \EndFor
    \State 对$\frac1B\sum_bR_b$反向传播（梯度即式~\eqref{eq:nlp-generation-mrt-gradient}），更新$\bm\theta$
  \end{algorithmic}
\end{algorithm}

算法~\ref{alg:nlp-generation-mrt}中，候选的采样过程不参与求导，梯度只经过各候选的对数概率。MRT与第~\ref{sec:nlp-inference-mbr}小节的最小贝叶斯风险解码形式相同，都是对候选集合上的期望代价求最小：MBR在推断时从候选中选择输出，MRT在训练时调整参数，使低代价的候选获得更多概率。MRT的代价是每个样本要生成$N$条候选并逐一计算概率与代价，训练开销远高于交叉熵；结果也依赖代价函数与$\alpha$的选择。

表~\ref{tab:nlp-generation-exposure-methods}比较三类对策。它们都需要先用条件最大似然训练出可用的模型，再在此基础上修正；大语言模型时代，后两类思路以偏好对齐和基于结果奖励的强化学习的形式继续使用。

\begin{table*}[htbp]
  \centering
  \small
  \begin{tabularx}{\linewidth}{@{}>{\raggedright\arraybackslash}p{22mm}
    >{\raggedright\arraybackslash}X>{\raggedright\arraybackslash}X
    >{\raggedright\arraybackslash}X@{}}
    \toprule
    方法 & 解决的问题 & 核心机制 & 新代价或局限 \\
    \midrule
    教师强制 & 并行、稳定地学习条件分布 & 每个位置以参考前缀为条件，逐词元交叉熵 & 训练前缀与生成前缀不同，目标与序列指标不一致 \\
    计划采样 & 模型从未见过自己的错误前缀 & 按日程以模型预测替换参考前缀 & 标签可能与偏离的前缀错位；Transformer需两次前向 \\
    序列级训练 & 交叉熵与整句指标不一致 & 采样整句，按奖励减基线的REINFORCE梯度更新 & 方差高，需预训练；模型会利用指标漏洞 \\
    最小风险训练 & 单样本梯度方差大 & 在候选集合上用锐化分布精确计算期望代价 & 每样本生成并打分$N$条候选，训练慢；依赖$\alpha$与代价函数 \\
    \bottomrule
  \end{tabularx}
  \caption{训练与生成不一致的对策。自上而下，对参考前缀的依赖逐步减少；后三种方法都从教师强制训练的模型出发。}
  \label{tab:nlp-generation-exposure-methods}
\end{table*}

\subsection{条件生成的共同挑战}
\label{sec:nlp-generation-challenges}

有了读取条件、复制和训练的共同机制，还要看清各任务反复遇到哪些困难。以下七项挑战贯穿本章，后续各节按它们组织方法谱系。

\begin{itemize}
  \item \emph{内容选择}：条件中的信息多于输出能容纳的信息，必须决定说什么。五条活动记录写成一句日历提示时，报名截止条件可以省略；写成校外来宾指南时却必须保留。摘要从长文中选句，数据到文本从记录表中选记录，都属于这一挑战。
  \item \emph{罕见词与复制}：名称、编号、数字和术语在词表中稀少，却必须原样保留。“海棠楼201”“ML-0418”被写成相近的常见词，句子仍然通顺，信息已经错了。
  \item \emph{重排与对齐}：输出的顺序和粒度与输入不同。中文“4月18日（周五）14:00”译成英文时，星期、日期和时间的语序都要变化；数据到文本要把无序的记录排成合理的叙述顺序；评价和编辑还需要知道输出的哪一段对应输入的哪一段。
  \item \emph{忠实性与无依据内容}：输出中出现条件不支持的内容，常称为\term{幻觉}（hallucination）。Maynez等区分两类：与来源矛盾的内在错误，例如把“校外来宾须报名”写成“无需报名”；来源未提及的外在错误，例如补写“主讲人为张老师”\scite{nlp-generation-maynez2020faithfulness}。开放写作允许虚构情节，但仍不得改动题目给定的事实。
  \item \emph{重复与退化}：模型重复同一短语、在长输出中兜圈子，或写出平淡空泛的句子。第~\ref{sec:nlp-inference-degeneration}小节已从解码角度分析了这一现象及采样对策；本章讨论任务侧的覆盖惩罚和内容规划。
  \item \emph{无平行数据}：许多转换任务没有成对的输入输出，例如同一内容的正式与口语版本、低资源语言之间的译文。方法需要从单语语料、回译或可分离的属性标记中构造训练信号。
  \item \emph{长输入}：多文档摘要、篇章翻译和长篇写作的条件超出模型的上下文窗口，或者虽然放得下却难以有效利用。需要分层编码、先选后写或递归规划。
\end{itemize}

表~\ref{tab:nlp-generation-challenges}把每项挑战对应到负责的方法和所在位置，可作为全章的索引。一种方法常同时回应几项挑战，例如显式规划既做内容选择也缓解长篇中的离题；表中只列主要对应关系。

\begin{table*}[htbp]
  \centering
  \small
  \begin{tabularx}{\linewidth}{@{}>{\raggedright\arraybackslash}p{24mm}
    >{\raggedright\arraybackslash}X>{\raggedright\arraybackslash}p{40mm}@{}}
    \toprule
    挑战 & 负责的方法 & 所在位置 \\
    \midrule
    内容选择 & 经典流水线的内容确定；内容选择门与指针规划；抽取式摘要与句子打分 & 第~\ref{sec:nlp-generation-data-to-text}节，第~\ref{sec:nlp-generation-summarization}节 \\
    罕见词与复制 & 复制与指针机制；子词词表；术语与词汇约束解码 & 第~\ref{sec:nlp-generation-copy-pointer}小节，第~\ref{sec:nlp-generation-constraints}节 \\
    重排与对齐 & 交叉注意力；短语翻译的扭曲模型与词对齐；编辑距离对齐 & 第~\ref{sec:nlp-generation-encoder-decoder}小节，第~\ref{sec:nlp-generation-translation}节，第~\ref{sec:nlp-generation-editing}节 \\
    忠实性与无依据内容 & 复制与记录回查；规划后生成；蕴含与问答式核验 & 第~\ref{sec:nlp-generation-data-to-text}节，第~\ref{sec:nlp-generation-faithfulness}小节 \\
    重复与退化 & 覆盖机制；分层与规划后写作；截断采样与候选重排 & 第~\ref{sec:nlp-generation-summarization}节，第~\ref{sec:nlp-generation-open-writing}节，第~\ref{chap:nlp-model-inference}章 \\
    无平行数据 & 删除—检索—生成与回译式风格转换；回译与无监督翻译；挖掘复述 & 第~\ref{sec:nlp-generation-editing}节，第~\ref{sec:nlp-generation-translation}节 \\
    长输入 & 长文档与分层摘要；篇章翻译；递归规划与修订 & 第~\ref{sec:nlp-generation-summarization}节，第~\ref{sec:nlp-generation-translation}节，第~\ref{sec:nlp-generation-open-writing}节 \\
    \bottomrule
  \end{tabularx}
  \caption{条件生成的七项共同挑战及其在本章中的对应方法。训练与生成的不一致贯穿所有挑战，其对策集中在第~\ref{sec:nlp-generation-exposure-bias}小节；各项挑战的评价方法见第~\ref{sec:nlp-generation-evaluation}节。}
  \label{tab:nlp-generation-challenges}
\end{table*}

\subsection{结果解析、核验与修订}
\label{sec:nlp-generation-recovery}

模型输出的词元不是最后的交付格式，也不能直接视为正确。交付前要完成三步：把词元恢复成可检查的对象，逐项核对事实，只对出错部分做局部修订并重新核验。

恢复按分词器规则合并子词、处理空格和特殊标记、保留段落边界；结构化输出还要解析字段名、类型和必填项。生成的日期字符串可以解析为2025年4月18日，随后仍要核对它属于\code{ML-0418}的活动时间，而不是报名截止时间：可解析与值正确是两项检查。编辑类任务还需要知道改了哪里，这要用编辑距离对齐源与输出，其递推与回溯见第~\ref{sec:nlp-generation-editing}节。

核验先做硬筛，再在合格候选之间比较。假设任务要求两句校外参加指南，候选甲为“读书会4月18日14:00在海棠楼201举行。校外来宾须在4月16日18:00前报名。”，候选乙为“读书会4月18日14:00在海棠楼201举行。校外来宾无需报名。”即使教学评分分别为$0.78,0.91$，乙因与来源矛盾先被淘汰，分数高不能补偿事实错误。甲进入合格集合后才与其他合格候选比较简洁程度；“提交报名信息”若在业务上不等于报名成功，甲还应修订为原通知的准确动作。

修订要给出可执行的反馈。设另有一份教学记录“资料袋共12份”，而候选写成“已备21份资料袋”。反馈应保存候选位置、错误值21、记录值12及允许的局部替换，而不只是“请检查”。替换后重新核对整段的实体绑定、日期、数量、句数和格式，因为局部修订也可能引入新错误。原始材料本身冲突时，核验器不能擅自选值，应返回待澄清状态。算法~\ref{alg:nlp-generation-verify-revise}给出这一流程。

\Needspace{12\baselineskip}
\begin{algorithm}[结果解析、核验与局部修订]
  \label{alg:nlp-generation-verify-revise}
  \begin{algorithmic}[1]
    \Require 条件$c$；候选$y$；字段解析器；核验器集合（实体绑定、否定、日期、数量、格式）；修订预算$K$
    \Ensure 通过核验的文本，或失败原因及证据
    \For{$k=0,1,\ldots,K$}
      \State 把$y$恢复为文本与字段；解析失败则记录位置并转到修订
      \State 逐项核验，收集错误集合$E$，每项含位置、错误值、来源值与来源位置
      \If{$E=\varnothing$}
        \State \Return $y$
      \EndIf
      \If{$E$中某项的来源值本身冲突}
        \State \Return 待澄清状态与冲突来源
      \EndIf
      \State 只对$E$涉及的片段做局部替换或重新生成，得到新的$y$
    \EndFor
    \State \Return 失败：最后一轮的$E$与候选
  \end{algorithmic}
\end{algorithm}

核验器自身也会出错，需要用独立标注检查其精度和召回，不能只让生成器认可自己的输出。基于蕴含和问答的忠实性检测器见第~\ref{sec:nlp-generation-faithfulness}小节；需要先检索证据再核验的问题留到第~\ref{chap:nlp-question-answering-dialogue}章。

\section{数据到文本}
\label{sec:nlp-generation-data-to-text}

\term{数据到文本生成}（data-to-text generation）以表格、记录或三元组为输入，生成描述它们的自然语言。输入已经是结构化的事实，模型不必再识别信息，难点转到表达上：在许多记录中选出该说的几条，安排叙述顺序，把数值和名称准确写出，并且不补造记录中没有的事实。这四项分别对应第~\ref{sec:nlp-generation-challenges}小节的内容选择、重排、复制和忠实性挑战。

方法谱系沿这四项展开。人工流水线把生成拆成若干阶段，每一阶段由规则完成，可控、可检查，但规则难以迁移到新领域。端到端神经生成去掉人工规则，输出流畅，却常写出记录不支持的数字和陈述，内容选择也不稳定。显式内容选择与规划把流水线的前两个阶段改成可学习的模块，在写句子之前先产生可检查的计划。结构线性化后交给预训练编码器—解码器或大语言模型，则借用预训练学到的语言能力改善流畅性和低资源表现。节末表~\ref{tab:nlp-generation-d2t-methods}比较四者。

本节使用五条规范化记录，每条写成（实体，类型，值）三元组，实体均为活动编号\code{ML-0418}：$r_1$的类型为名称，值为“机器学习读书会”；$r_2$为时间“2025年4月18日14:00”；$r_3$为地点“海棠楼201”；$r_4$为校内规则“校内师生可直接参加”；$r_5$为校外规则“校外来宾须在2025年4月16日18:00前提交报名信息”。这是从原通知抽出的当前安排视图，没有原日期字段，因而不能单凭这张表生成“从4月17日改期”的历史陈述。

\subsection{经典生成流水线}
\label{sec:nlp-generation-classic-pipeline}

在神经方法之前，实用的生成系统多采用Reiter与Dale总结的流水线\scite{nlp-generation-reiter2000building}。它把“从数据到句子”拆成三个模块：\term{文档规划}（document planning）决定说什么、按什么结构说，包括内容确定与文档结构化两项任务；\term{微规划}（microplanning）决定用什么词、怎样指代、哪些事实合并成一句；\term{表层实现}（surface realization）按语法把句子计划变成字符串，处理词序、虚词和标点。每一阶段的输出都是显式的中间表示，可以单独检查和修改，这是流水线的主要优点；代价是规则和模板需要按领域人工编写，换领域或换语言时重写成本高，表达也较单一。

\subsubsection{内容确定、文档规划与微规划}
\label{sec:nlp-generation-nlg-pipeline-stages}

\term{内容确定}（content determination）从全部记录中选出要报道的条目，通常由领域规则完成。设任务为“给校外来宾的两句通知”，规则可以写成：名称、时间、地点必选；校内规则与读者无关，不选；校外规则必选。五条记录因此选出$r_1,r_2,r_3,r_5$，得到四条待表达的消息。

文档结构化把消息组织成有序的文档计划，常用修辞关系连接。这里第一句介绍活动本身（$r_1,r_2,r_3$），第二句说明参加条件（$r_5$），两句之间是“说明”关系：后者补充前者的参加方式。文档计划可以写成树：根节点为“说明”，左子树是由$r_1,r_2,r_3$组成的事件描述，右子树是$r_5$。

微规划把文档计划变成句子计划。词汇化为每条消息选词，例如时间类型选“于……举行”而不是“定在”；指代表达生成决定第二次提到活动时用全名、简称还是省略，第二句以“校外来宾”开头，不必再提活动名；句子聚合把同一事件的名称、时间和地点合并成一个小句，而不是写成三句“活动名为……。时间为……。地点为……。”。最后，表层实现按中文语序输出
\begin{quote}
机器学习读书会于4月18日14:00在海棠楼201举行。校外来宾须在4月16日18:00前提交报名信息。
\end{quote}
日期按规则从完整格式缩写为“4月18日”，这一步也要记录在实现规则中，否则读者无法确认年份是否被有意省略。

\subsubsection{模板与槽位填充}
\label{sec:nlp-generation-template}

字段和句式稳定时，可以把微规划与表层实现合并成模板：
“\{名称\}于\{时间\}在\{地点\}举行。\{校外规则\}。”
模板不需要训练，参数来自经确认的模式、日期格式和语法规则。执行时先按活动编号连接字段，统一时区与日期，再填槽；不能把两条不同活动的同名字段按表格行号直接拼接。

填入上述记录后，第一句的三个槽位分别回指$r_1,r_2,r_3$，第二句回指$r_5$，每个陈述都能追溯到一条记录。若地点缺失，模板应选择明确的缺值分支，如“地点尚未提供”，或者停止生成并报告缺字段，不能填入默认教室。两条完全相同的地点记录只表达一次，但日期相同而编号不同的活动不能去重：去重按对象身份进行，不按文字进行。

模板的事实覆盖可控，弱点在组合。每增加一种字段组合、一种缺值情形或一种语气，模板数就要成倍增加；大量记录写成长文时，固定句式反复出现，读起来机械。把同样的记录直接交给条件生成器，表达更灵活，却可能写出“读书会周五下午举行，欢迎参加”：语句通顺，却没有表达精确时间、地点和报名规则，以五条记录为单位核对，它至多完整覆盖$r_1$。下一小节的端到端方法正是从这种取舍出发的。

\subsection{端到端神经生成与无依据内容}
\label{sec:nlp-generation-end-to-end-d2t}

端到端方法的动机是去掉人工规则：把全部记录交给编码器—解码器，由数据自动学出选择、排序和措辞。Wiseman等在体育比赛数据上系统地检验了这一路线\scite{nlp-generation-wiseman2017}。他们构建的RotoWire数据集把NBA比赛的球员与球队统计表与人写的赛后报道配对，每场比赛平均约有六百条记录，而一篇报道只提到其中一小部分，因此内容选择是主要难点。

模型的输入是记录集合$R=\{r_1,\ldots,r_J\}$，每条记录有实体$e_j$、类型$t_j$和值$m_j$（原数据还有主客队标记）。记录编码把各字段嵌入拼接后经一层前馈网络：
\[
  \bm u_j=\bigl[\bm E_e(e_j);\bm E_t(t_j);\bm E_m(m_j)\bigr],\qquad
  \bm r_j=\operatorname{ReLU}(\bm W_r\bm u_j+\bm b_r)\in\mathbb R^{d}.
\]
解码器是带注意力的LSTM，第$t$步按式~\eqref{eq:nlp-generation-bahdanau}对全部$\bm r_j$求注意力$a_{t,j}$。复制单位是记录的值：门控复制按
\[
  p_t(w)=p(z_t=1\mid\bm s_t)\sum_{j:\,m_j=w}a_{t,j}+p(z_t=0\mid\bm s_t)\,P_{\mathrm{vocab}}(w)
\]
混合，与式~\eqref{eq:nlp-generation-pointer-generator}相同，只是复制分布定义在记录值而不是源端词元上；他们也比较了CopyNet式的共享归一化复制。为了让解码器状态覆盖更多输入记录，他们还试用了只在训练时使用的重建损失：把解码器隐状态切成长度不超过$B$的连续块，要求从每一块预测出表中若干条记录的实体、类型与值，以表中最匹配记录的负对数似然计损失，作为辅助目标加到式~\eqref{eq:nlp-generation-nll}上。

只看BLEU无法判断报道说对了什么，Wiseman等因此提出\term{抽取式评价}（extractive evaluation）：先在参考报道上训练一个关系抽取器，再用它从生成文本$\hat y$中抽出（实体，值，类型）关系序列$E(\hat y)$，从参考文本中抽出$E(y^{\star})$，计算三项指标。关系生成（RG）统计$E(\hat y)$中有多少条能在记录$R$中找到，衡量忠实性；内容选择（CS）以$E(y^{\star})$为参照计算$E(\hat y)$的精确率与召回率，衡量是否说了人会说的事；内容排序（CO）用两条关系序列之间的归一化Damerau--Levenshtein距离衡量顺序是否相近。例如生成句“读书会于4月18日14:00在海棠楼301举行，主讲人为张老师”抽出三条关系，只有时间一条在记录中，RG精确率为$1/3$；地点值写错，主讲人是记录没有的补造内容（教学示例）。

评价结果揭示了端到端路线的典型问题。神经模型的输出比模板流畅得多，但RG精确率明显低于模板：模型会写出记录中不存在的数字、把一名球员的数据安到另一名球员身上，或补写记录无法支持的评论；CS也不稳定，同一模型时而遗漏关键记录、时而罗列次要统计。原因在于所有决定都隐含在逐词生成里：注意力每步只局部地决定读哪条记录，没有一个全局计划约束“这篇报道要讲哪几条、按什么顺序”。抽取式评价本身也依赖抽取器的精度，抽取错误会同时影响三项指标。下一小节把内容选择与规划重新显式化。

\subsection{内容选择与规划后生成}
\label{sec:nlp-generation-content-planning}

Puduppully等的方法把流水线中的内容确定与文档规划变成可学习的模块，并保留端到端训练\scite{nlp-generation-puduppully2019}。\term{内容选择与规划}（content selection and planning）先从记录集合$R$产生记录编号序列$z=(z_1,\ldots,z_K)$作为内容计划，再以计划为条件写文本。图~\ref{fig:nlp-generation-content-planning}给出五个计算阶段。

\begin{figure*}[tbp]
  \centering
  % 内容选择与规划后生成（Puduppully等，2019）的五个计算阶段及两项训练损失。
% 下方示例行使用正文的五条活动记录，为教学示例。
\begin{tikzpicture}[
  x=1mm,y=1mm,figure picture,
  stage/.style={figure node,minimum width=24mm,text width=22mm,minimum height=17mm,
    inner sep=2pt,font=\figurefont\scriptsize,align=flush center},
  ex/.style={font=\figurefont\scriptsize,text=FigureMuted,text width=26mm,align=flush center},
  loss/.style={figure node,draw=FigureLoss,fill=FigureLossFill,minimum height=6mm,
    inner sep=1.5pt,font=\figurefont\scriptsize},
  flow/.style={figure flow}
]
  \node[stage,draw=FigureInput,fill=FigureInputFill] (enc) at (0,0)
    {记录编码\\字段嵌入拼接后\\经前馈层得$\bm r_j$};
  \node[stage,draw=FigureModel] (gate) at (30,0)
    {内容选择门\\$\bm r^{\mathrm{cs}}_j=\bm g_j\odot\bm r_j$};
  \node[stage,draw=FigureModel] (plan) at (60,0)
    {指针规划\\$p(z_k=j\mid z_{<k},R)$};
  \node[stage,draw=FigureModel] (penc) at (90,0)
    {计划编码\\双向LSTM得$\bm e_k$};
  \node[stage,draw=FigureOutput,fill=FigureOutputFill] (dec) at (120,0)
    {文本解码\\注意力$\bm e_k$与复制};
  \foreach \a/\b in {enc/gate,gate/plan,plan/penc,penc/dec}{
    \draw[flow] (\a.east) -- (\b.west);
  }
  \node[ex,anchor=north] at (enc.south) {$r_1,\ldots,r_5$各为\\（实体，类型，值）};
  \node[ex,anchor=north] at (gate.south) {每条记录读取\\其他记录后得门值};
  \node[ex,anchor=north] at (plan.south) {$z=(r_1,r_2,r_3,r_5,$\\$\mathrm{EOS})$};
  \node[ex,anchor=north] at (penc.south) {只编码计划中的\\四条记录};
  \node[ex,anchor=north] at (dec.south) {“读书会于4月18日\\14:00在海棠楼201……”};
  \node[loss] (lp) at (60,14) {$-\log p_{\bm\phi}(z\mid R)$};
  \node[loss] (ly) at (120,14) {$-\log p_{\bm\theta}(y\mid R,z)$};
  \draw[figure link,densely dashed] (lp) -- (plan.north);
  \draw[figure link,densely dashed] (ly) -- (dec.north);
\end{tikzpicture}

  \caption{内容选择与规划后生成的计算阶段。记录编码与端到端方法相同；内容选择门让每条记录参照其他记录调节自身表示；指针网络在记录集合上逐步选出内容计划，并以结束符终止；计划经双向LSTM编码后，由带注意力与复制的解码器写出文本。虚线框为两项训练损失，分别监督计划与文本。下方示例行使用本节的五条活动记录，为教学示例。}
  \label{fig:nlp-generation-content-planning}
\end{figure*}

第一步是内容选择门。一条记录是否值得报道，取决于其他记录：一名球员得分很高才值得提，平常的数值则不必。门控让每条记录先注意其他记录，再据此缩放自身表示：
\begin{equation}
  \begin{gathered}
  \alpha_{j,k}=\frac{\exp(\bm r_j^{\top}\bm W_a\bm r_k)}{\sum_{k'\neq j}\exp(\bm r_j^{\top}\bm W_a\bm r_{k'})},\qquad
  \bm c_j=\sum_{k\neq j}\alpha_{j,k}\bm r_k,\\
  \bm g_j=\sigma\bigl(\bm W_g[\bm r_j;\bm c_j]\bigr),\qquad
  \bm r^{\mathrm{cs}}_j=\bm g_j\odot\bm r_j .
  \end{gathered}
  \label{eq:nlp-generation-content-gate}
\end{equation}
$\bm c_j$汇总与第$j$条记录相关的其他记录，门向量$\bm g_j\in(0,1)^d$逐维决定保留多少，$\odot$是逐元素乘法。门没有单独的标签，随后续两项损失一起训练。

第二步是指针规划。一个LSTM解码器在记录集合上运行指针网络：初始状态取全部$\bm r^{\mathrm{cs}}_j$的平均，第$k$步以上一选中记录的$\bm r^{\mathrm{cs}}$为输入得到状态$\bm h_k$，按$p(z_k=j\mid z_{<k},R)\propto\exp(\bm h_k^{\top}\bm W_c\bm r^{\mathrm{cs}}_j)$在记录上归一化。它与式~\eqref{eq:nlp-generation-pointer-network}的区别只在打分用双线性形式而非加性形式，输出同样是输入位置。计划长度不固定，实现时在记录集合中加一个特殊的结束记录，选中它即停止。第三步，把选中的记录$\bm r^{\mathrm{cs}}_{z_1},\ldots,\bm r^{\mathrm{cs}}_{z_K}$依次送入双向LSTM，得到计划编码$\bm e_1,\ldots,\bm e_K$。第四步，文本解码器只对$\bm e_k$做注意力，并从计划中的记录复制值。解码器不再面对数百条记录，只需把一条短计划依次写出来。

训练需要记录、计划和文本三种监督，目标是两项负对数似然之和：
\begin{equation}
  L_{\mathrm{plan}}
  =-\log p_{\bm{\phi}}(z\mid R)
   -\log p_{\bm{\theta}}(y\mid R,z).
  \label{eq:nlp-generation-plan-loss}
\end{equation}
第一项监督选择与次序，第二项监督如何把选定记录说出来。数据集只有记录与文本，计划标签要另行构造：用上一小节的关系抽取器从参考文本中抽出关系，按出现顺序匹配到记录，得到弱监督的计划。推断时两个阶段各用一次束搜索：先求近似最优的计划$\hat z$，再以$\hat z$为条件生成文本。

以给校外来宾的两句通知为例，计划为$z=(r_1,r_2,r_3,r_5,\mathrm{EOS})$。若规划器给这五步的概率依次为$0.7,0.6,0.8,0.9,0.95$，第一项为$-\log(0.7\times0.6\times0.8\times0.9\times0.95)=-\log0.28728\approx1.2473$；文本部分沿用第~\ref{sec:nlp-generation-training}小节的三词元算例，第二项为$1.4271$，总损失约$2.6744$（教学数值）。两项可以分别监控：计划损失下降而文本损失不降，说明问题在表达；反之说明选错了内容。

显式计划带来两项好处。其一，内容选择可以在写句子之前单独检查：计划中有$r_5$而无$r_4$，正符合来宾通知的用途；若任务改成完整通知，同样的计划就不合格。其二，解码器只读计划中的少数记录，补造其他记录内容的机会减少。在RotoWire测试集上，Puduppully等报告的RG精确率从Wiseman等端到端基线的约75\%提高到约87\%。代价同样来自计划：弱对齐的计划标签含有抽取错误；推断要多运行一次规划解码；计划一旦漏选某条记录，文本解码器只读计划，无法再补回。复制能写出“海棠楼201”，但不阻止模型把它绑定到另一场活动，生成后仍要把每个陈述对齐回$r_j$核验。

\subsection{结构线性化与预训练模型}
\label{sec:nlp-generation-linearization}

前两小节的模型从零训练，数据到文本的平行语料又通常只有数千到数万对，模型要同时学会语言和数据，表达能力受限。预训练编码器—解码器已经掌握了语言表达，剩下的问题是把结构送进一个只接受词元序列的模型。最常用的做法是\term{线性化}（linearization）：按固定顺序把结构展开成带分隔标记的文本。三元组可以写成
\begin{center}
  \code{<H> ML-0418 <R> 时间 <T> 2025-04-18 14:00 <H> ML-0418 <R> 地点 <T> 海棠楼201}
\end{center}
其中\code{<H>}、\code{<R>}、\code{<T>}分别标出主语、关系和宾语；表格可以按行展开，每个单元格前加列名与行名。线性化后的序列作为编码器输入，参考文本作为目标，按式~\eqref{eq:nlp-generation-nll}微调T5或BART；两者的去噪预训练目标见第~\ref{sec:nlp-llm-objective-t5-bart}小节。Kale与Rastogi用这一做法微调T5，在表格到文本、三元组到文本和对话行为到文本等任务上都超过了专门设计的结构编码器，在训练中未见过的类别上提高尤其明显\scite{nlp-generation-kale2020t5d2t}。

两个数据集体现了这一阶段的任务设置。WebNLG以DBpedia中的一到七条RDF三元组为输入，测试集包含训练时未出现的类别，用来检验对新领域关系的泛化\scite{nlp-intro-gardent2017}。ToTTo以维基百科表格加一组高亮单元格为输入，目标是一句只描述高亮内容的句子\scite{nlp-generation-parikh2020totto}。高亮单元格把内容选择从模型手中拿走，任务集中在忠实表达上；参考句由标注者从原维基句子逐步修改而来，删去表格不支持的内容，因此适合衡量忠实性。

线性化把结构信息变成了顺序信息，带来三个实现问题。第一，顺序本身没有含义，但模型会对顺序敏感：同一组三元组换一种排列，输出可能不同。训练时随机打乱三元组顺序可以减少这种依赖。第二，长表格超出编码器长度时只能截断或先选子表，ToTTo的高亮单元格正是一种预先选择。第三，数值和日期按子词切分后，“2025-04-18”可能被拆成若干片段，格式不统一时模型容易抄错；应统一数值格式，必要时让模型输出记录编号或占位符，再由程序填值。

大语言模型可以直接据提示完成数据到文本：把记录或表格放进提示，说明读者与篇幅，必要时给几个示例。这省去了微调数据，却把内容选择重新变成隐式的：模型自己决定说哪些记录，也可能用常识补写记录没有的内容。对大模型输出，应把每个数值和实体回查到记录字段，按第~\ref{sec:nlp-generation-recovery}小节的流程核验；需要可审计的内容选择时，可以先让模型输出记录编号形式的计划，经检查后再写文本，这正是上一小节显式规划在提示层面的对应做法。

表~\ref{tab:nlp-generation-d2t-methods}比较本节四类方法。预训练路线显著改善了流畅性和低资源表现，显式规划仍在可审计性上有优势，两者可以结合：规划部分输出记录编号，表达部分由预训练模型完成。

\begin{table*}[htbp]
  \centering
  \small
  \begin{tabularx}{\linewidth}{@{}>{\raggedright\arraybackslash}p{24mm}
    >{\raggedright\arraybackslash}X>{\raggedright\arraybackslash}X
    >{\raggedright\arraybackslash}X@{}}
    \toprule
    方法 & 解决的问题 & 核心机制 & 新代价或局限 \\
    \midrule
    流水线与模板 & 需要可控、可逐步检查的生成 & 内容确定、文档规划、微规划、表层实现由规则或模板完成 & 规则按领域人工编写，组合随字段增多而膨胀，表达单一 \\
    端到端编码器—解码器 & 去掉人工规则，从数据学出选择与措辞 & 记录嵌入、注意力与门控复制，抽取式评价衡量忠实性 & 无依据内容多，内容选择不稳定，所有决定隐含在逐词生成中 \\
    内容选择与规划 & 选择与次序缺少全局约束 & 内容选择门、指针网络规划、计划编码后再生成 & 计划标签靠弱对齐；多一次规划解码；漏选无法补回 \\
    线性化与预训练模型 & 平行数据少，表达能力受限 & 结构展开为带分隔符的序列，微调T5或BART，或提示大语言模型 & 对顺序敏感，长表需截断，内容选择重新变为隐式 \\
    \bottomrule
  \end{tabularx}
  \caption{数据到文本的方法谱系。自上而下，前一行的局限大致是后一行的动机；后一行并不在所有维度上都优于前一行，例如模板的事实覆盖最容易保证。}
  \label{tab:nlp-generation-d2t-methods}
\end{table*}

\section{开放写作}
\label{sec:nlp-generation-open-writing}

开放写作包括故事、续写和带风格要求的自由文本。它允许虚构，内容不再由输入记录逐条决定，挑战随之改变：长文本要前后连贯、情节要有组织；采样需要在多样性与第~\ref{sec:nlp-inference-degeneration}小节的重复退化之间取得平衡；常常还要控制情感、话题或正式程度等属性；长篇中的人物设定和已写事实必须保持一致。

方法谱系按这些挑战排列。分层故事生成先产生计划或提示，再据此写正文，处理长程连贯。属性控制有三条路线：训练时把属性作为条件，解码时用属性模型的梯度修改隐状态，解码时用属性分类器对候选重加权；第一条要训练整个生成器，后两条冻结生成器、只外接小模型，第三条还省去了解码时的反向传播。数千词的长篇超出单次生成的连贯范围，需要递归规划与修订。大语言模型直接续写再对候选重排，是所有路线的对照基线。节末表~\ref{tab:nlp-generation-writing-methods}汇总比较。

本节固定一个写作任务。题目为“失而复得的借书卡”，人物只有林岚，故事发生在同一天的图书馆，分三段，每段一句，文体平实；允许虚构找卡经过，但不得新增人物替林岚找到卡，也不得把图书馆业务规则写成真实规定。为小说构造情节不属于事实错误，前提是任务明确允许虚构；评价分别检查情节多样性、阅读偏好与约束满足，任何一个维度都不能代表全部创作质量。

\subsection{分层故事生成}
\label{sec:nlp-generation-hierarchical-story}

单层语言模型从标题直接写长故事，常在几句之后偏离主题：每一步只依据前文局部地决定下一个词，开头的线索没有在结尾回收，远处的句子之间缺少共同的约束。\term{分层生成}（hierarchical generation）先产生一个高层计划，再以计划为条件写正文，使相距很远的句子共享同一条件。两种代表路线的区别在于计划从哪里来：Plan-and-Write从故事中自动抽取关键词作为计划监督，Fan等的分层生成以人写的写作提示作为上层。

\subsubsection{Plan-and-Write情节规划}
\label{sec:nlp-generation-plan-write}

Plan-and-Write先用有序的关键词表示故事线，再依据故事线生成文本\scite{nlp-generation-yao2019}。它在五句短故事语料上训练，用RAKE关键词抽取算法从每句抽出一个关键词，形成与句子一一对应的故事线，作为规划监督。故事线只是关键词序列，不是完备的事件逻辑，也不是人物状态记录。

原方法比较两种规划方式。静态规划由题目$u$一次生成完整故事线$z_{1:k}$，再生成故事$y$，两个模型分别最大化$p(z\mid u)$与$p(y\mid u,z)$。动态规划交替进行：以题目、此前句子和上一关键词预测下一关键词，再据此写下一句，规划条件是$(u,y_{<j},z_{j-1})$，写作条件还包含新关键词$z_j$。两个模型可以分别训练、运行时串接。静态方式在写正文之前就决定了全部情节，便于检查首尾呼应；动态方式能响应已写内容，但计划和正文交替推进，前面的线索可能在后面被遗忘。

静态教学计划为“遗失／回找／夹页”，据此形成三段：
\begin{quote}
林岚在图书馆借书时发现借书卡不见了，记起自己刚翻过一本厚书。

她返回阅览桌，把那本厚书从原处取下来，一页页查看。

借书卡正夹在书页之间；她收好卡，回到借书处。
\end{quote}
另设一份教学检查记录：第一段林岚在借书处，卡的位置未知；第二段她到阅览桌，取书后开始找；第三段卡从书页转到她手中，她返回借书处。开头的厚书线索在结尾回收，人物移动和事件先后可以逐项检查。这份记录是本书的核验手段，不是原模型维护的状态。

动态方式写完第一段后，也可能选择“衣袋”作为第二个关键词，形成“她先检查外套衣袋，找到夹在票据中的借书卡”，再以“收好”结尾。这时第一段的“厚书”线索失去作用；若第三段仍生成“从书中取出刚找到的卡”，同一张卡就出现在两个位置。静态规划也会失败：计划写“衣袋”而正文仍沿用“夹页”时，只有核对最终文本才能发现计划与实现脱节。

\subsubsection{提示到故事的分层生成}
\label{sec:nlp-generation-prompt-story}

Fan等把计划换成一句人写的写作提示，建立了提示—故事配对的WritingPrompts数据集，数据来自网络写作社区，故事平均长达数百词\scite{nlp-inference-fan2018story}。生成分两层：先用卷积语言模型生成一句提示，再以提示为条件生成故事。提示比关键词序列信息更丰富，也更接近人类写作时的起点。

故事模型是卷积编码器—解码器。卷积只能看到有限窗口，故事又长达数百词，他们在解码器中加入门控多尺度自注意力，让每个位置能读取全部已写内容：查询、键和值不用线性投影，而由带门控线性单元的小网络给出；各头在不同时间尺度上读取历史，第一个头看全部位置，第二个头每隔一步取一个位置，第三个头每隔两步取一个，依此类推，使不同的头关注不同范围的信息。更关键的问题是，在提示—故事数据上直接训练的编码器—解码器常常忽视提示，只学会写通顺的故事：词与词之间的局部依赖比故事与提示之间的依赖容易学得多。他们的对策是\term{模型融合}（model fusion）：先训练一个编码器—解码器并固定参数，再训练第二个编码器—解码器，第二个模型每步读取第一个模型的隐状态。设两者在第$t$步的隐状态为$\bm h^{\mathrm{new}}_t$与$\bm h^{\mathrm{pre}}_t$，融合层先拼接再逐维门控，
\[
  \bm g_t=\sigma\bigl(\bm W[\bm h^{\mathrm{new}}_t;\bm h^{\mathrm{pre}}_t]+\bm b\bigr),\qquad
  \bm h_t=\bm g_t\odot[\bm h^{\mathrm{new}}_t;\bm h^{\mathrm{pre}}_t],
\]
再经若干全连接层进入输出层。Fan等把这看作一种残差学习：第一个模型已经承担了大部分语言建模，第二个模型可以集中学习第一个模型没有学到的部分，例如故事对提示的依赖。

生成时，Fan等使用top-k采样（$k=10$），每步只在概率最高的十个词元中重新归一化抽样；束搜索在开放写作中产生重复、平淡的文本，原因及截断采样的一般形式见第~\ref{sec:nlp-generation-sampling}小节。与Plan-and-Write相比，提示—故事路线的上层更自然，但提示本身不规定情节结构，长故事中段仍可能离题；两者都只做一层规划，第~\ref{sec:nlp-generation-long-form}小节把规划扩展到多层。

\subsection{属性控制生成}
\label{sec:nlp-generation-attribute-control}

许多写作任务还要求文本具有指定属性$a$：情感为正面、话题为科技、措辞正式。目标是从条件分布$p(x\mid a)$生成。由贝叶斯公式，$p(x\mid a)\propto p(x)\,p(a\mid x)$；逐词元写开，第$t$步应满足
\begin{equation}
  p(x_t\mid x_{<t},a)\propto p(x_t\mid x_{<t})\,p(a\mid x_{\le t}),
  \label{eq:nlp-generation-attribute-bayes}
\end{equation}
其中$p(a\mid x_{\le t})$是“当前前缀写完后满足属性”的概率。三条路线以不同方式实现这个关系。CTRL直接训练条件模型$p(x\mid a)$，控制强但新增属性需重训；PPLM冻结语言模型，用属性模型的梯度修改隐状态，近似地把分布推向高$p(a\mid x)$的区域，不需重训，但每步要多次反向传播；FUDGE训练一个前缀分类器估计$p(a\mid x_{\le t})$，按式~\eqref{eq:nlp-generation-attribute-bayes}对候选重加权，只需前向计算。表~\ref{tab:nlp-generation-attribute-control}在本小节末比较三者。

\subsubsection{控制码条件生成：CTRL}
\label{sec:nlp-generation-ctrl}

CTRL是一个约16亿参数的Transformer语言模型\scite{nlp-inference-keskar2019ctrl}。训练时在每条序列开头放一个\term{控制码}（control code），即标明这条文本来源或类型的特殊前缀，学习条件语言模型
\[
  p(x\mid a)=\prod_{t}p(x_t\mid x_{<t},a).
\]
控制码来自语料的来源信息：维基百科、图书、评论等领域，评论的评分，网页的网址，以及问答、翻译等任务标记。生成时写下“Reviews Rating: 5.0”再接开头，模型就按五星评论的风格续写。

控制码把属性控制变成普通的条件生成，推断不增加任何开销，控制效果也最直接。它的局限来自训练：控制码集合在训练时就固定了，只能控制语料中天然带有标签的属性；想要“正式而简短”这类新属性或任意组合，只能重新训练或微调。CTRL还提出了第~\ref{chap:nlp-model-inference}章介绍的重复惩罚。

\subsubsection{梯度引导控制：PPLM}
\label{sec:nlp-generation-pplm}

\term{即插即用语言模型}（plug and play language model, PPLM）的出发点是不改动已经训练好的语言模型，只在解码时外接一个小的属性模型\scite{nlp-generation-dathathri2020pplm}。属性模型给出$p(a\mid x)$，有两种：词袋模型取一组属性词，以下一步分布落在这组词上的总概率作为$p(a\mid x)$；判别器模型在语言模型隐状态的平均上接一个线性分类器。属性模型的参数很少，可以按需更换。

控制通过修改语言模型的历史状态实现。设$\bm H_t$是第$t$步之前各层的键值缓存，PPLM给它加一个扰动$\Delta\bm H_t$，沿属性对数概率的梯度方向更新若干次：
\begin{equation}
  \Delta\bm H_t\leftarrow\Delta\bm H_t+\eta\,
  \frac{\nabla_{\Delta\bm H_t}\log p(a\mid\bm H_t+\Delta\bm H_t)}
  {\bigl\lVert\nabla_{\Delta\bm H_t}\log p(a\mid\bm H_t+\Delta\bm H_t)\bigr\rVert^{\gamma}},
  \label{eq:nlp-generation-pplm}
\end{equation}
$\eta$是步长，分母按梯度范数的$\gamma$次方归一化，控制每次修改的幅度。只追求属性会损害流畅性，因此更新目标中还加了一项KL惩罚，使修改后的下一词元分布$\widetilde p_{t+1}$不远离原分布$p_{t+1}$；采样时再取二者的几何加权平均$\widetilde p_{t+1}^{\,\gamma_{\mathrm{gm}}}p_{t+1}^{\,1-\gamma_{\mathrm{gm}}}$并重新归一化，$\gamma_{\mathrm{gm}}$接近1时更偏向控制后的分布。

\Needspace{12\baselineskip}
\begin{algorithm}[PPLM的单步生成]
  \label{alg:nlp-generation-pplm}
  \begin{algorithmic}[1]
    \Require 冻结的语言模型；历史键值$\bm H_t$；属性模型$p(a\mid\cdot)$；步数$m$，步长$\eta$，系数$\gamma,\lambda_{\mathrm{KL}},\gamma_{\mathrm{gm}}$
    \Ensure 下一词元$x_{t+1}$与更新后的历史
    \State 用$\bm H_t$前向一次，得到原分布$p_{t+1}$；令$\Delta\bm H_t\gets\bm 0$
    \For{$i=1,\ldots,m$}
      \State 用$\bm H_t+\Delta\bm H_t$前向，得到$\widetilde p_{t+1}$与$\log p(a\mid\bm H_t+\Delta\bm H_t)$
      \State 目标取$\log p(a\mid\cdot)-\lambda_{\mathrm{KL}}\operatorname{KL}(\widetilde p_{t+1}\Vert p_{t+1})$，反向传播到$\Delta\bm H_t$
      \State 按式~\eqref{eq:nlp-generation-pplm}更新$\Delta\bm H_t$
    \EndFor
    \State 以$\widetilde p_{t+1}^{\,\gamma_{\mathrm{gm}}}p_{t+1}^{\,1-\gamma_{\mathrm{gm}}}$归一化后采样$x_{t+1}$
    \State 用修改后的历史前向$x_{t+1}$，追加其键值，进入下一步
  \end{algorithmic}
\end{algorithm}

算法~\ref{alg:nlp-generation-pplm}的每一步都要做$m$次前向与反向传播，生成速度比普通采样慢数倍；步长过大时文本会变得不通顺甚至重复属性词，需要调节$\eta,\lambda_{\mathrm{KL}},\gamma_{\mathrm{gm}}$。原文每步更新$m=3$到$10$次，$\lambda_{\mathrm{KL}}=0.01$，$\gamma_{\mathrm{gm}}$取$0.8$到$0.95$，梯度范数的归一化对Transformer每层分别计算。与CTRL相比，PPLM不需重训且属性模型可任意更换；与下文的FUDGE相比，它修改的是隐状态而不是输出概率，代价是每步的反向传播。

\subsubsection{FUDGE软属性控制}
\label{sec:nlp-generation-fudge}

正式、礼貌等属性难以写成一个词袋。FUDGE直接实现式~\eqref{eq:nlp-generation-attribute-bayes}：训练一个前缀分类器$q_{\bm\phi}(a\mid x_{\le t})$，预测“当前前缀继续写完后，完整文本会不会满足属性”，再用它对下一词元重加权\scite{yang-klein-2021-fudge}。训练样本是带完整文本属性标签的序列，把同一个完整文本的标签附给它的每个前缀，用二元交叉熵训练；分类器判断的是前缀的未来走向，不是前缀本身是否已经正式。生成器保持不变。

记生成器的前缀为$h=x_{<t}$，候选下一词元为$v$。重加权后的分布为
\begin{equation}
  \widetilde p(v\mid h,a)=
  \frac{p_{\bm{\theta}}(v\mid h)\,q_{\bm{\phi}}(a\mid hv)}
  {\sum_w p_{\bm{\theta}}(w\mid h)\,q_{\bm{\phi}}(a\mid hw)}.
  \label{eq:nlp-generation-fudge}
\end{equation}
输入材料也属于条件，这里为简洁略去。分母要对每个候选调用一次分类器，实际实现只对基础分布概率最高的若干词元计算，这也改变了可搜索的集合。分类器只是近似，不一定与生成器的真实未来分布校准一致，因此公式实现的是控制的近似，不是属性满足的保证。

给定开头“校外来宾”，三个教学候选片段为“须提交”“记得交”“无需提交”，基础概率为$(0.40,0.50,0.10)$，预测最终正式的概率为$(0.90,0.20,0.95)$。乘积为$(0.36,0.10,0.095)$，总和$0.555$，归一化为$(0.649,0.180,0.171)$。束宽为一时，选择从口语的“记得交”变成“须提交”。“无需提交”同样很正式，控制后它的相对概率反而从$0.10$升到$0.171$，却与通知要求矛盾：软属性控制只管措辞，内容是否正确要单独核验。

FUDGE只需前向计算，分类器可以比生成器小得多，多个属性的分类器还可以相乘组合；后者只有在相应条件独立假设下才有概率解释，“极短”与“完整保留长条件”这类属性还可能直接冲突。与PPLM相比，FUDGE不修改隐状态，控制强度受限于基础分布前若干候选中有没有满足属性的词元。

\begin{table*}[htbp]
  \centering
  \small
  \begin{tabularx}{\linewidth}{@{}>{\raggedright\arraybackslash}p{18mm}
    >{\raggedright\arraybackslash}X>{\raggedright\arraybackslash}X
    >{\raggedright\arraybackslash}X>{\raggedright\arraybackslash}X@{}}
    \toprule
    方法 & 控制的实现 & 训练代价 & 推断代价 & 控制强度与流畅性 \\
    \midrule
    CTRL & 训练条件模型$p(x\mid a)$，控制码放在序列开头 & 整个模型在带来源标签的语料上训练；新属性需重训 & 与普通生成相同 & 控制直接；属性集合固定，难以组合 \\
    PPLM & 沿$\nabla\log p(a\mid x)$修改键值历史，KL项与几何平均保持流畅 & 只训练小属性模型 & 每步$m$次前向与反向 & 控制较强；步长过大损害流畅性 \\
    FUDGE & 前缀分类器估计$p(a\mid x_{\le t})$，对候选按式~\eqref{eq:nlp-generation-fudge}重加权 & 训练前缀分类器 & 每步对前若干候选各做一次分类器前向 & 不改动生成器，流畅性好；强度受候选集合限制 \\
    \bottomrule
  \end{tabularx}
  \caption{三种属性控制方法的比较。三者都实现式~\eqref{eq:nlp-generation-attribute-bayes}的某种近似：CTRL在训练时学习条件分布，PPLM和FUDGE在解码时引入外部属性模型。属性控制只改变措辞或话题倾向，不检查内容是否与条件一致。}
  \label{tab:nlp-generation-attribute-control}
\end{table*}

\subsection{长篇规划与递归修订}
\label{sec:nlp-generation-long-form}

前两小节的方法处理的是几百词的故事。数千词的长篇超出了单次生成能保持连贯的范围：上下文窗口装不下全部已写内容，即使装得下，模型也会逐渐偏离最初的设定，前文定下的人物属性在后文被改写。分层故事生成只有一层计划，不足以约束这么长的文本。

Re3（recursive reprompting and revision）用大语言模型的提示把长篇写作拆成“计划—草拟—重排—编辑”的循环\scite{nlp-generation-yang2022re3}。计划阶段由前提生成背景、人物及其描述和分条大纲。草拟阶段每次写一段，把与当前大纲条目相关的计划内容、前文摘要和最近写出的几段重新组织进提示，生成多个候选续写；这种每段重建提示的做法称为\term{递归重提示}（recursive reprompting），使远处的设定以摘要形式始终留在上下文中。重排阶段用两个训练好的打分模型选出与大纲最相关、与前文最连贯的候选。编辑阶段从已写内容中抽取人物属性，检测新段落与之矛盾的地方，再调用编辑模型修改。

\Needspace{12\baselineskip}
\begin{algorithm}[递归重提示与修订的长篇生成]
  \label{alg:nlp-generation-re3}
  \begin{algorithmic}[1]
    \Require 前提$u$；大语言模型；相关性与连贯性打分模型；每段候选数$N$；段数上限$M$
    \Ensure 长篇文本$y$
    \State 由$u$生成背景、人物描述与大纲条目$o_1,\ldots,o_K$；令$y\gets\varnothing$
    \For{$j=1,\ldots,M$}
      \State 选出当前大纲条目$o_k$，组织提示：相关计划、前文摘要、最近若干段
      \State 生成$N$条候选段落
      \State 按与$o_k$的相关性和与$y$的连贯性重排，取最高者$y_j$
      \State 抽取人物属性，检测$y_j$与已写内容的矛盾，必要时编辑$y_j$
      \State 把$y_j$接到$y$后，更新前文摘要；当前条目写完则转到$o_{k+1}$
    \EndFor
    \State \Return $y$
  \end{algorithmic}
\end{algorithm}

DOC在Re3的基础上加深计划\scite{nlp-generation-yang2023doc}。它把大纲逐层细化成多级条目，每个条目附有场景和出场人物，细化时对候选子条目过滤和重排；生成每一段时，用一个按FUDGE方式实现的控制器（基于OPT-350m）逐词元调整候选的概率，使正文贴合当前的细化条目。控制器用段落与其摘要的配对做对比训练，负例包括流畅但与摘要无关的续写，使长段落不只是开头切题。这里第~\ref{sec:nlp-generation-fudge}小节的属性从“正式”换成了“与某个大纲条目相关”。

长篇规划路线的代价有三项。其一，调用次数随段数和候选数成倍增加，一篇长篇需要数百次模型调用和多个辅助打分模型。其二，错误沿层级传播：大纲层的矛盾会被忠实地写进每一段，编辑阶段只能修补局部，无法改正结构问题。其三，评价长篇的连贯和有趣主要依赖人工判断，自动指标难以反映，第~\ref{sec:nlp-generation-evaluation}节将讨论这一问题。

\subsection{语言模型直接续写与候选重排}
\label{sec:nlp-generation-continuation}

作为所有路线的对照基线，可以把题目、人物要求和已写段落作为前缀，让大语言模型直接续写，生成方式与采样规则沿用第~\ref{sec:nlp-generation-sampling}小节。运行时要写清模型能否看到此前所有段落、用何种段落标记，以及第三段结束后的EOS；没有EOS而撞上长度上限的文本仍是截断候选。多轮共同修改题目或人物设定的交互过程见第~\ref{chap:nlp-question-answering-dialogue}章。

直接续写的结果需要筛选和重排，方法见第~\ref{sec:nlp-inference-reranking}小节。固定上一小节故事的第一段，教学候选甲续写“返回阅览桌找书，卡夹在书中”，候选乙续写“管理员把卡送来，林岚感谢他”。乙新增人物代为找卡，违反已定条件，即使语言偏好评分更高也先被淘汰。在合法候选之间，再按线索回收、重复和约定风格重排；重排器可以是人工量表、规则或有标注的偏好模型，来源不同需要分别验证。

比较规划路线与直接续写时，要把规划消耗的词元与调用计入预算。例如设两条完整候选、每条最多90个新词元：规划路线每条分配12个规划词元和78个正文词元，直接路线每条90个正文词元，两者再接受相同次数的核验，记录实际词元、前向调用和延迟（教学协议）。规划在写正文之前决定部分情节，直接续写在正文中逐步决定；完整上下文可用时二者都能读取已写内容，发生截断时都可能丢失人物条件。偏好最高的一条仍须检查位置、时间线、题目要求和允许虚构的边界。

表~\ref{tab:nlp-generation-writing-methods}汇总本节方法。规划、控制和修订各自回应开放写作的一项挑战，可以组合使用；DOC就是分层规划与FUDGE式控制的组合。

\begin{table*}[htbp]
  \centering
  \small
  \begin{tabularx}{\linewidth}{@{}>{\raggedright\arraybackslash}p{24mm}
    >{\raggedright\arraybackslash}X>{\raggedright\arraybackslash}X
    >{\raggedright\arraybackslash}X@{}}
    \toprule
    方法 & 解决的问题 & 核心机制 & 新代价或局限 \\
    \midrule
    Plan-and-Write & 单层生成离题，首尾不呼应 & 自动抽取的关键词故事线，静态或动态规划后写作 & 关键词不表达事件逻辑；计划与正文可能脱节 \\
    提示到故事 & 计划信息太少；模型忽视条件 & 先生成提示，卷积编码器—解码器加多尺度自注意力，模型融合 & 只有一层规划，长故事中段仍会离题 \\
    CTRL & 需要指定风格或领域 & 训练时以控制码为条件 & 属性集合固定，新属性需重训 \\
    PPLM & 不重训而外接属性 & 属性梯度修改键值历史 & 每步多次反向传播，流畅性受步长影响 \\
    FUDGE & 降低外接控制的代价 & 前缀分类器对候选重加权 & 近似未校准；强度受候选集合限制 \\
    Re3与DOC & 数千词长篇的连贯与人物一致 & 多层大纲，递归重提示，候选重排与一致性编辑 & 调用次数多；大纲错误沿层级传播；评价依赖人工 \\
    直接续写与重排 & 作为对照基线，实现最简单 & 前缀续写，截断采样，规则或偏好模型重排 & 无显式计划；长文本中设定易被遗忘 \\
    \bottomrule
  \end{tabularx}
  \caption{开放写作的方法谱系。前两行处理长程连贯，中间三行处理属性控制，Re3与DOC把规划扩展到长篇；直接续写是各路线的共同基线。比较时须计入规划、控制与重排消耗的调用和词元。}
  \label{tab:nlp-generation-writing-methods}
\end{table*}

\section{文本编辑}
\label{sec:nlp-generation-editing}

前几节的输出几乎都要从头写出；\term{文本编辑}（text editing）则从一段已经存在的文字出发，
输出与输入的大部分词元相同。改写调整表达，风格转换改变指定属性，纠错修复错误，
简化为特定读者降低理解负担。四类任务共有四个挑战：保留不该改的内容，
多数位置的正确动作是原样复制；找到该改的位置，修改往往稀疏而分散；
缺少同内容不同表达的平行句对，风格转换几乎没有；区分漏改、正确修改、过改和意义改变，
把“须提交”改成“不必提交”只动了几个字，意义变化却大于一次保持原意的整句重写。

方法谱系沿两条线展开。一条是输出的表示方式：规则与统计替换只在少数位置提出候选；
整句序列生成用第~\ref{sec:nlp-generation-encoder-decoder}小节的编码器—解码器重写全句，
借第~\ref{sec:nlp-generation-copy-pointer}小节的复制保留原词；编辑标签方法只为每个源词元
预测一个动作；大语言模型以提示完成编辑。另一条是监督来源：有平行句对时直接训练，
没有时用回译、属性语料、合成错误或挖掘复述构造训练信号。四个小节按修改理由划分任务，
在每个任务内沿这两条线比较方法，节末表~\ref{tab:nlp-generation-editing-methods}汇总对照。

\paragraph{从源与输出的对齐恢复修改}
上面两条线都要回答同一个问题：一对句子之间究竟改了什么。编辑标签方法要据此把平行句对
转成逐词元的训练标签，整句生成要据此从输出中找出实际修改，评价也要在修改集合上计数。
\term{编辑距离}（edit distance）是把源序列$x_{1:n}$变成目标序列$y_{1:m}$所需的
最少插入、删除和替换次数；用动态规划求它并回溯路径，就得到两者的
\term{Levenshtein对齐}。设$d_{i,j}$为前缀$x_{1:i}$与$y_{1:j}$之间的编辑距离，
以$d_{i,0}=i$、$d_{0,j}=j$初始化，按行递推
\begin{equation}
  d_{i,j}=\min\bigl\{\,d_{i-1,j}+1,\;\; d_{i,j-1}+1,\;\;
  d_{i-1,j-1}+\mathbb{I}[x_i\neq y_j]\,\bigr\}.
  \label{eq:nlp-generation-edit-distance}
\end{equation}
三项依次对应删除$x_i$、插入$y_j$、替换或保留当前位置。递推只依赖左、上和左上三个格子，
时间与空间都是$O(nm)$；记录每格取到最小值的来源，从$(n,m)$回溯到$(0,0)$，
就得到操作序列及每个操作的源坐标。几个来源并列最小时，按预先固定的优先次序取一个。

以$x=$“go／to／meeting”、$y=$“goes／to／the／meeting”为例，
表~\ref{tab:nlp-generation-levenshtein}给出全部$d_{i,j}$。右下角$d_{3,4}=2$；
回溯时，$(3,4)$来自左上方的匹配meeting，$(2,3)$来自左方，即插入the，
$(2,2)$来自左上方的匹配to，$(1,1)$来自左上方的替换go为goes。
两次编辑恰好是一次词形修正和一次冠词插入。

\begin{table}[htbp]
  \centering
  \small
  \begin{tabular}{@{}lccccc@{}}
    \toprule
    & $\varnothing$ & goes & to & the & meeting\\
    \midrule
    $\varnothing$ & 0 & 1 & 2 & 3 & 4\\
    go & 1 & \textbf{1} & 2 & 3 & 4\\
    to & 2 & 2 & \textbf{1} & \textbf{2} & 3\\
    meeting & 3 & 3 & 2 & 2 & \textbf{2}\\
    \bottomrule
  \end{tabular}
  \caption{“go to meeting”到“goes to the meeting”的编辑距离表。行是源前缀，列是目标前缀，
  粗体格是回溯路径：从右下角经匹配、插入、匹配回到替换，共两次编辑。}
  \label{tab:nlp-generation-levenshtein}
\end{table}

若去掉替换操作、只允许插入和删除，编辑距离的回溯路径就是最长公共子序列对齐，
后文LaserTagger构造标签用的正是这一特例；GECToR用Levenshtein对齐的一个变体，
第~\ref{sec:nlp-generation-edit-metrics}小节的编辑指标也在对齐恢复的编辑集合上计数。
对齐只说明改了哪里，不说明改得对不对：删除“不得参加”中的“不”同样只需一次编辑。

\subsection{改写与复述}
\label{sec:nlp-generation-paraphrase}

\term{改写}（paraphrasing）要求意义不变而表达变化。以通知为例，把“须在4月16日18:00前提交报名信息”
改成“必须在4月16日18:00前提交报名信息”可以接受；改成“可在4月16日提交信息”
同时丢掉了义务、精确时点和“前”。改写有两个核心困难：训练需要大量意义相同的句对，
即\term{复述}（paraphrase），而人工写复述代价高；句子高度重合时，整句生成仍要逐词重写，
既浪费计算，也给补造内容留下机会。回译批量构造复述对，回应前一点；
局部编辑标注只预测改动，回应后一点。

\subsubsection{回译复述}
\label{sec:nlp-generation-backtranslation-paraphrase}

两句话若能译成同一句外语，它们多半表达相同的意思。
Bannard与Callison-Burch据此从双语平行语料中抽取短语复述：
设$e_1,e_2$为同一语言的两个短语，$f$为与它们对齐的外语短语，
复述概率按枢轴语言上的全概率展开
\begin{equation}
  p(e_2\mid e_1)=\sum_{f}p(e_2\mid f)\,p(f\mid e_1),
  \label{eq:nlp-generation-pivot-paraphrase}
\end{equation}
两个条件概率都由短语表的对齐计数估计\scite{bannard2005paraphrasing}
式~\eqref{eq:nlp-generation-pivot-paraphrase}把“意义相同”换成“共享译文”，
其中求和覆盖了一个短语的多种译法，单个错误对齐的影响因此被其他译法稀释。

神经翻译把这一思想推广到整句：把句子$x$译成枢轴语言$z$，再把$z$译回原语言，
得到复述$\tilde x$。Mallinson等用多条枢轴译文及多个枢轴语言对回译分布加权平均，
近似式~\eqref{eq:nlp-generation-pivot-paraphrase}的求和\scite{mallinson2017paraphrasing}
Wieting与Gimpel走得更远：他们取大规模捷克语—英语平行语料，把捷克语一侧机器翻译成英语，
与原有英语句配成五千万量级的复述对ParaNMT，再按句向量相似度和词重合度过滤，
用来训练复述句向量\scite{wieting2018paranmt}
这些做法只需现成的翻译系统或双语语料，不需要人工复述。

回译的多样性取决于译回时怎样选择候选。束搜索倾向于给出最常见的译法，
结果常与原句几乎相同；从模型分布中采样或在束内引入随机性，
表达更多样，也更容易偏离原意。因此要过滤：相似度太低的句对可能已经改变意义，
太高的又几乎没有改写。设$\operatorname{sim}$为句向量余弦相似度，可只保留
$\tau_{\mathrm{low}}\leq\operatorname{sim}(x,\tilde x)\leq\tau_{\mathrm{high}}$且编辑距离大于零的句对。
整句相似度对少数关键词不敏感：通知经英语回译成“来宾可在4月16日提交信息”时，
相似度可能仍然不低，却丢掉了人群、义务和时点。日期、实体和否定词这类保护片段
应在过滤阶段逐项比对。

回译复述的长处是规模大、成本低、领域可以跟随单语材料；
代价是意义漂移只能事后过滤，而且复述的多样性受限于翻译系统本身的表达习惯。

\subsubsection{局部编辑标注：LaserTagger}
\label{sec:nlp-generation-lasertagger}

有了复述对，可以用编码器—解码器整句生成输出。但在句子融合、拆句和局部改写里，
输出往往只在几个位置与输入不同，整句生成需要逐词重写那些本该照抄的部分。
LaserTagger把输出表示为对源词元的一串编辑动作，只需预测这些动作，
再由确定的程序重建输出\scite{malmi2019lasertagger}

\paragraph{标签与短语词表}
LaserTagger为每个源词元$x_i$分配一个标签$t_i=(b_i,p_i)$：基础动作
$b_i\in\{\code{KEEP},\code{DELETE}\}$决定是否保留$x_i$，可选的插入短语$p_i$
放在$x_i$之前，取值于\term{短语词表}（phrase vocabulary）$V$或为空。
标签集的大小是$2(|V|+1)$。替换由“删除原词并在它之前插入新短语”表示。
以句子融合为例，源文本
“Turing／was／born／in／1912／.／Turing／died／in／1954／.”的目标是
“Turing was born in 1912 and died in 1954.”，只需给第一个句号打\code{DELETE}，
给第二个Turing打\code{DELETE}，再给died打$(\code{KEEP},\text{and})$，即在died之前插入and，
其余全部\code{KEEP}。
中文通知中“须”改为“必须”，同样是在“须”上打$(\code{DELETE},\text{必须})$。

短语词表决定了模型能表达哪些输出。把训练集每对句子用最长公共子序列对齐，
未被对齐的目标片段就是需要插入的短语。理想的$V$应在大小上限$k$内
使尽可能多的训练样本可以完整重建，这个选择问题是NP困难的；
LaserTagger按短语在训练集中被需要的次数排序，取前$k$个。
例如三对训练句分别需要插入$\{\text{and}\}$、$\{\text{and}\}$和$\{\text{, which}\}$，
取$k=1$时选and，覆盖三对中的两对；第三对无法用标签表示，训练时被丢弃。

\paragraph{标签转换与模型}
给定平行句对和$V$，下面的贪心程序把目标转成标签；转换失败的样本不参与训练。

\begin{algorithm}[LaserTagger的标签转换]
  \label{alg:nlp-generation-lasertagger-convert}
  \begin{algorithmic}[1]
    \Require 源词元$x_{1:n}$，目标词元$y_{1:m}$，短语词表$V$
    \Ensure 标签序列$t_{1:n}$，或“无法表示”
    \State $j\gets1$\Comment{目标中下一个待匹配的位置}
    \For{$i=1,\ldots,n$}
      \If{存在$j'\geq j$使$y_{j'}=x_i$，且$y_{j:j'-1}$为空或属于$V$}
        \State 取满足条件的最小$j'$；$t_i\gets(\code{KEEP},y_{j:j'-1})$；$j\gets j'+1$
      \Else
        \State $t_i\gets(\code{DELETE},\varnothing)$
      \EndIf
    \EndFor
    \If{$j\leq m$}\Comment{目标末尾仍有未匹配词元}
      \State 若$y_{j:m}\in V$，把它作为插入短语附到末尾的特殊位置；否则\Return 无法表示
    \EndIf
  \end{algorithmic}
\end{algorithm}

对上面的句子融合，算法~\ref{alg:nlp-generation-lasertagger-convert}处理第一个句号时，
目标中下一个句号在句末，两者之间的“and died in 1954”不在词表中，打\code{DELETE}；第二个Turing在目标剩余部分中不出现，也打\code{DELETE}；
died与目标匹配，其前的未匹配片段and在词表中，于是得到$(\code{KEEP},\text{and})$。
模型用BERT编码源句，得到$\bm H\in\mathbb{R}^{n\times768}$。
前馈版本LaserTagger\textsubscript{FF}对每个位置独立分类：
$p(t_i\mid x)=\operatorname{softmax}(\bm W\bm h_i+\bm b)$，$\bm W\in\mathbb{R}^{|T|\times768}$，
$|T|$为标签数；自回归版本LaserTagger\textsubscript{AR}在编码器上接一层Transformer解码器，
逐位置预测标签，并以已预测标签为条件。两者都最小化逐位置交叉熵
$-\sum_{i}\log p(t_i\mid x,t_{<i})$，前馈版本中条件$t_{<i}$不出现。
推断时把标签依次应用到源句即可得到输出。

\paragraph{关联、推广与边界}
LaserTagger是第~\ref{chap:nlp-text-analysis}章序列标注在生成任务中的用法：
输出空间从“任意词元序列”收缩为“每个源词元一个标签”，前馈版本可以并行预测全部位置，
推断比同规模的序列生成快得多。它的约束也来自这里：插入只能取词表内的短语，
大幅重排无法表达，训练集中需要罕见插入的样本只能丢弃。
Seq2Edits把标签推广为带跨度的编辑序列，每步输出“编辑类型、跨度终点、替换短语”三元组，
替换短语可从完整词表生成，从而能表达跨多个词的修改\scite{stahlberg2020seq2edits}
在源与目标高度重合、需要插入的短语种类有限时，编辑标签方法通常比整句生成更快，
也更少补造内容；修改幅度越大，整句生成越有优势。

\subsection{无平行数据的风格转换}
\label{sec:nlp-generation-style-transfer}

\term{风格转换}（style transfer）要改变句子的一个属性，例如情感、正式程度或礼貌程度，
同时保留其余内容。形式上，给定按属性分开的单语语料$X_1,\ldots,X_K$，
没有同内容不同属性的句对，要学习映射$x\mapsto\tilde x$，使$\tilde x$具有目标属性$a'$，
并与$x$讲同一件事。缺少平行数据是这项任务的核心挑战，各方法的差别在于用什么替代平行监督：
DRG假定属性集中在少数可删除的短语，用检索与重建代替；
隐表示方法试图把内容与属性分开，或改用回译直接构造伪平行对；
大语言模型借预训练中学到的风格知识，以提示直接转换。
评价也因此分三项：目标属性是否达成，通常用属性分类器判断；
内容是否保持，用与源句的相似度或人工判断；输出是否流畅。
三项相互牵制：原样复制源句的内容保持最好，属性却没有改变。

\subsubsection{Delete--Retrieve--Generate属性转换}
\label{sec:nlp-generation-delete-retrieve}

情感等属性常集中体现在少数短语上：“令人失望”“十分周到”带有明显倾向，
“这家书店的服务”则不带。Delete--Retrieve--Generate（DRG）据此把句子拆成
属性标记与内容两部分：删除原属性的标记，从目标属性语料检索适合当前内容的标记，
再生成新句\scite{li-etal-2018-delete}

\paragraph{属性标记的识别}
设$n_a(u)$为短语$u$在属性$a$语料中的出现次数。DRG用平滑后的相对频率衡量$u$
对属性$a$的\term{属性标记}（attribute marker）程度：
\begin{equation}
  s(u,a)=\frac{n_a(u)+\lambda}{\sum_{a'\neq a}n_{a'}(u)+\lambda},
  \label{eq:nlp-generation-drg-salience}
\end{equation}
$\lambda>0$防止分母为零，并压低只出现一两次的短语。$s(u,a)$超过阈值$\gamma$的
$n$元短语（原文取到四元）被视为属性$a$的标记。
例如“令人失望”在负面语料中出现40次、在正面语料中出现2次，取$\lambda=1$得
$s=41/3\approx13.7$；“服务”在两类语料中分别出现30次和32次，$s=31/33\approx0.94$。
取$\gamma=5$时，只有前者被删除。删除句$x$中属性$a$的全部标记后，
得到内容序列$c(x,a)$，被删的标记序列记为$m(x,a)$。

\paragraph{检索与生成}
检索在目标属性$a'$的语料中寻找内容最接近的句子$x'$：
$x'=\argmin_{x''\in X_{a'}}\operatorname{dist}\bigl(c(x,a),c(x'',a')\bigr)$，
距离可取TF-IDF加权的词重合，或句向量之间的欧氏距离。$x'$的标记$m(x',a')$
就是候选的目标属性表达。由此有四种输出方式：
RetrieveOnly直接返回$x'$；TemplateBased把$m(x',a')$填回$x$中被删除的位置；
DeleteOnly用LSTM编码$c(x,a)$，把编码与目标属性的可学习嵌入拼接作为解码器初态，生成新句；
DeleteAndRetrieve再用另一个编码器读取$m(x',a')$，与内容编码一起交给解码器。

生成器没有跨属性的平行目标，只能用重建训练：以句子自身的$c(x,a)$和属性$a$为输入，
最大化原句$x$的概率。DeleteAndRetrieve的训练输入还包括自身的标记$m(x,a)$，
若原样提供，模型只需学会照抄；因此训练时把它替换为同属性语料中相近的另一组标记，
迫使模型学会把检索来的标记接到内容上。运行时把$a$换成目标属性$a'$、
把$m(x,a)$换成检索到的$m(x',a')$。

以负面到正面的情感转换为例，原句“这家书店的服务令人失望”删除“令人失望”后
留下“这家书店的服务”，正面语料中内容最近的句子是“这家商场的服务十分周到”。
RetrieveOnly直接返回它，把书店换成了商场；TemplateBased把“十分周到”填回空位，
得到“这家书店的服务十分周到”；DeleteOnly只凭内容和属性生成，可能输出“这家书店服务很好”；
DeleteAndRetrieve把检索来的“十分周到”接到内容上，并可调整衔接。
这些是为说明四种方式构造的候选，不是实测输出。

DRG的优点是每一步都可检查：被删的标记、检索到的句子和生成依据都能看到。
它的假设是属性可以由局部短语承载；讽刺、隐含评价或分布在整句结构中的风格不满足这一假设。
按频率识别的标记也可能误伤内容：“须”“不得”这类义务词即使在正式语料中更常见，
也不能当作文体标记删除。

\subsubsection{隐表示解纠缠与回译转换}
\label{sec:nlp-generation-latent-style}

DRG把内容与风格的分离放在词面上。另一条路线把分离放到隐表示里：
若编码器能把句子映射为只含内容的向量$\bm z$，解码器再以目标风格为条件生成，
就不必假设风格集中在可删除的短语。

\paragraph{交叉对齐自编码器}
Shen等考虑两种风格的语料$X_1,X_2$\scite{shen2017crossalign}
编码器$E(x,y)$读取句子$x$及其风格$y$，输出内容表示$\bm z$；生成器$G(y,\bm z)$
是以风格$y$和$\bm z$为初态的GRU。自编码重建损失为
\begin{equation}
  L_{\mathrm{rec}}=\sum_{k=1}^{2}\mathbb{E}_{x\sim X_k}
  \bigl[-\log p_G\bigl(x\mid y_k,E(x,y_k)\bigr)\bigr].
  \label{eq:nlp-generation-shen-rec}
\end{equation}
只有重建并不能阻止$\bm z$携带风格。他们先尝试对齐两种风格下$\bm z$的分布，
用一个判别器区分$\bm z$来自哪种语料，再让编码器骗过它；
进一步的\term{交叉对齐}（cross-alignment）改为对齐生成器的隐状态序列：
判别器$D_1$区分两类序列，一类是生成器以$y_1$为条件、在教师强制下重建真实$x_1$时的隐状态
$\bm h_1$，另一类是把$x_2$的内容$E(x_2,y_2)$配上风格$y_1$、自由生成时的隐状态
$\tilde{\bm h}_2$。若生成器转换得好，后者应与真实风格$y_1$的句子不可区分。
自由生成的离散采样不可求导，他们用带温度的softmax输出作为下一步输入，使梯度可以回传。
$D_2$对称地处理另一方向，总目标为
\begin{equation}
  \begin{aligned}
    &\min_{E,G}\max_{D_1,D_2}\;L_{\mathrm{rec}}
    -\rho\,(L_{\mathrm{adv},1}+L_{\mathrm{adv},2}),\\
    &L_{\mathrm{adv},1}=-\mathbb{E}\bigl[\log D_1(\bm h_1)\bigr]
    -\mathbb{E}\bigl[\log\bigl(1-D_1(\tilde{\bm h}_2)\bigr)\bigr].
  \end{aligned}
  \label{eq:nlp-generation-shen-adv}
\end{equation}
这里$\rho>0$平衡重建与对抗两项。判别器最大化自己的分类对数似然，
编码器和生成器则让转换后的隐状态骗过判别器。这一对抗训练与生成对抗网络相同：
一方生成样本，另一方区分真伪，二者交替更新；其训练不稳定和需要调节更新步数的问题，
可参考\BookChapterRef{chap:neural-generative-models}{卷五的神经生成模型章}。

\paragraph{对抗解纠缠}
Fu等把对抗直接放在内容表示上：一个分类器试图从$\bm z$预测风格，
编码器则被训练成让这个分类器无法预测，再用风格专属的解码器或风格嵌入生成
\scite{fu2018styletransfer}
这类方法统称为解纠缠：希望$\bm z$与风格统计独立，从而“换风格不换内容”。
困难在于，对抗训练使某个分类器失败，并不说明$\bm z$已经不含风格；
训练后另取一个分类器，往往仍能从$\bm z$中读出风格。

\paragraph{去噪与回译}
Lample等的观察正是如此：解纠缠后的表示仍可读出属性，转换却照样有效，
说明解纠缠既难做到，也未必必要\scite{lample2019multiattr}
他们借鉴无监督机器翻译，用两项损失训练一个以属性为条件的编码器—解码器。
去噪自编码项把句子加噪（随机删词、局部打乱）后重建原句；
回译项先用当前模型把$x$转成另一属性$\tilde a$下的$\hat x=G(E(x),\tilde a)$，
不对这一步求导，再要求模型从$\hat x$还原$x$：
\begin{equation}
  L=\lambda_{\mathrm{ae}}\,\mathbb{E}_{x}
    \bigl[-\log p\bigl(x\mid E(\operatorname{noise}(x)),a\bigr)\bigr]
   +\lambda_{\mathrm{bt}}\,\mathbb{E}_{x,\tilde a}
    \bigl[-\log p\bigl(x\mid E(\hat x),a\bigr)\bigr].
  \label{eq:nlp-generation-lample-bt}
\end{equation}
回译项把模型自己生成的$(\hat x,x)$当作伪平行对：输入是属性$\tilde a$的句子，
输出是属性$a$的原句。随着模型变好，伪平行对的质量也提高。
Prabhumoye等用另一种回译：先把英语句子译成法语，再用法语到英语翻译模型的编码器
得到表示，假定翻译过程已削弱了作者风格，最后由各风格专属的解码器生成，
并用风格分类器的反馈训练解码器\scite{prabhumoye2018style}

交叉对齐与对抗解纠缠要求一个难以验证的中间状态，即不含风格的$\bm z$；
去噪加回译放弃这一要求，转而依赖迭代产生的伪平行对，
代价是早期的转换错误会写进伪数据，内容可能逐轮漂移。
两条路线的共同限制是内容保持只受重建损失间接约束，需要单独检查实体、数字与否定。

\subsubsection{大语言模型提示改写}
\label{sec:nlp-generation-llm-style}

前面几种方法都要为每种属性准备语料并单独训练。大语言模型在预训练中见过大量不同风格的文本，
可以直接按指令改写。Reif等提出增强零样本提示：提示中给出几个“原句—某种风格的改写”示例，
示例覆盖多种不同风格而不必包含目标风格，再要求模型把输入改写成目标风格
\scite{reif2022recipe}
示例的作用是让模型理解“只改风格、保留内容”这一任务格式，目标风格本身由一句自然语言描述，
因此“更有诗意”“像新闻标题”这类没有标注语料的风格也能请求。
上下文示例的选择与数目对结果的影响，见第~\ref{sec:nlp-inference-icl-shots}小节。

提示只描述了要求，模型是否做到要靠筛选。具体做法是在提示中列出必须保留的内容，
如“保持人群、日期、时刻与义务不变”；再对同一输入采样$k$个候选，
用风格分类器给出目标属性概率$p_{\mathrm{sty}}(a'\mid\tilde x)$，用句向量给出与源句的相似度$\operatorname{sim}(x,\tilde x)$，
剔除改动了保护片段的候选后，在相似度不低于阈值的候选中取$p_{\mathrm{sty}}$最大者，
或按两者的几何平均排序。这一步就是第~\ref{sec:nlp-inference-reranking}小节的候选重排。
提示改写省去了训练，适用的风格范围也最广；代价是每次调用的推断成本高，
过改与意义漂移需要筛选步骤控制，转换程度也难以像控制码那样精确调节。

\subsection{语法纠错}
\label{sec:nlp-generation-correction}

\term{语法纠错}（grammatical error correction, GEC）的依据是原文存在错误，
而不是另一种表达更好。“请与周五到场”中的同音字、“提高报名流程”中的搭配，
以及“She go to meeting”中的动词词形和冠词，都有明确的修订；
“请于周五到场”已经正确，改成“烦请周五莅临”属于文体选择，不能算作纠错收益。
英文流畅性纠错数据集JFLEG允许比最小语法修改更宽的改写，使用时应沿用其任务定义
\scite{nlp-intro-napoles2017}

纠错有四个挑战。错误稀疏，一句中多数词应原样保留，模型很容易学成“什么都不改”；
真实的错句—正句平行数据有限；一个错句常有多种合法修订；模型又容易过改，
把正确的部分也改掉。方法谱系为：混淆集加语言模型打分；统计翻译式纠错；
神经序列生成；编辑标签模型GECToR；大语言模型。合成错误数据贯穿后三者，
用来弥补平行数据的不足。评价时区分漏改、正确修改、过改与意义改变，
修改级指标见第~\ref{sec:nlp-generation-edit-metrics}小节。

\subsubsection{混淆集与语言分数筛选}
\label{sec:nlp-generation-confusion}

若错误集中在同音、形近或拼写混淆，可以先限制每个位置的替换范围。
\term{混淆集}（confusion set）是某个位置允许尝试的候选字词集合，它来自词典、输入法混淆统计、
人工规则或有标注错误。集合必须包含原词，否则系统遇到正确输入也被迫改写。
给完整候选句计分：
\begin{equation}
  s(y;x)=\log p_{\mathrm{LM}}(y)-\lambda\,d(x,y),
  \label{eq:nlp-generation-confusion-score}
\end{equation}
其中语言模型提供上下文分数，$d$为式~\eqref{eq:nlp-generation-edit-distance}的编辑距离，
$\lambda\geq0$惩罚不必要的改动。只有相对原句的得分增量超过开发集确定的阈值，
且规则检查通过，才执行替换。

对“请／与／周五／到场”，候选位置取$\{\text{与},\text{于},\text{予}\}$。
设三条完整句子的教学对数分数分别为$-8,-4,-10$，$\lambda=1$，
则保留原文得$-8$，替换为“于”得$-5$，替换为“予”得$-11$。
阈值为2时，增量3使“于”通过，再由时间介词规则核验，得到“请于周五到场”。

混淆集方法的覆盖面受限于集合本身：词典没有收录真实误字时，搜索不可能找到正确替换；
逐位置替换也无法补冠词、删重复词或调整远处的搭配。领域术语还须保护，
“束搜索”中的“束”虽不如“数”常见，却不能被改成“数搜索”。
突破逐位置替换的第一步是把纠错看作翻译：错句是源语言，正句是目标语言。
Junczys-Dowmunt与Grundkiewicz用第~\ref{sec:nlp-generation-smt}小节的短语统计翻译实现这一点，
在对数线性模型中加入编辑操作计数等稠密与稀疏特征，并直接按纠错评价指标调节特征权重，
得到了当时的强基线\scite{junczys2016pbmt}
短语表从平行纠错语料中学到“go to meeting→goes to the meeting”这类短语级修改，
不再需要人工列出候选；但它仍受短语表覆盖与局部重排的限制。
后文的神经方法沿用“纠错即翻译”的视角，并用合成数据弥补平行语料的不足。

\subsubsection{合成错误数据}
\label{sec:nlp-generation-synthetic-errors}

序列模型需要大量错句—正句对，真实标注语料却只有几十万到几百万句。
干净文本则几乎无限，若能从正句自动生成逼真的错句，就能构造任意规模的
\term{合成错误数据}（synthetic error data）。

\paragraph{规则加噪}
最直接的做法是对正句随机施加编辑。Grundkiewicz等对每个词元按概率执行替换、删除、
插入或与相邻词交换，其中替换优先从拼写检查器给出的混淆集中取词，
使生成的错误更接近真实拼写与用词错误\scite{grundkiewicz2019synthetic}
算法~\ref{alg:nlp-generation-noise}给出一般形式。

\begin{algorithm}[规则加噪生成错句]
  \label{alg:nlp-generation-noise}
  \begin{algorithmic}[1]
    \Require 正句$y_{1:m}$；错误率分布的均值$\mu$与标准差$\sigma$；
      四种操作的概率$(\pi_{\mathrm{sub}},\pi_{\mathrm{del}},\pi_{\mathrm{ins}},\pi_{\mathrm{swap}})$；
      混淆集$C(\cdot)$与词表$V$
    \Ensure 合成错句$x$
    \State 采样本句错误率$\rho\sim\mathcal{N}(\mu,\sigma^2)$，截断到$[0,1]$；$x\gets y$
    \For{$x$中的每个位置$i$}
      \State 以概率$\rho$决定是否在此处加噪；若不加噪，继续下一个位置
      \State 按$\pi$选择操作：替换为$C(x_i)$中的随机词；删除$x_i$；
        在$x_i$后插入$V$中的随机词；或交换$x_i$与$x_{i+1}$
    \EndFor
    \State \Return $(x,y)$作为训练对
  \end{algorithmic}
\end{algorithm}

$\mu$通常取开发集真实错句相对正句的词级错误率，使合成数据的错误密度与真实数据相近。
例如取四种操作概率为$0.7,0.1,0.1,0.1$、错误率$\rho=0.25$，
对“She goes to the meeting”的期望加噪次数为$5\times0.25=1.25$，
一次可能的结果是把goes换成混淆集中的go、再删除the，得到“She go to meeting”。

\paragraph{反向翻译生成错误}
规则噪声的分布与学习者的真实错误并不相同：真实错误集中在冠词、介词、词形和搭配，
随机插入的词则往往一眼可辨。另一种做法借用机器翻译中的回译：
在真实平行语料上反向训练一个“正句→错句”的模型，再用它把大量干净句子译成错句
\scite{xie2018noising}
若用束搜索解码，反向模型倾向于输出最常见的少数错误，甚至原样复制正句；
Xie等在束搜索的每个候选分数上加入随机噪声或排名惩罚，使生成的错误更多样。
Kiyono等在同一模型结构下比较了直接加噪与回译式合成，以及“先在合成数据上预训练、
再在真实数据上微调”与二者混合训练两种用法，结论是用带噪束搜索的回译数据预训练、
再用真实数据微调最为有效\scite{kiyono2019pseudo}

合成数据的风险在于分布偏移：模型在合成数据上学到的错误类型和错误密度若与真实输入不符，
就会漏改真实错误，或在正确句上“找错”。因此合成数据通常只用于预训练，
再用真实平行数据微调；训练数据中保留一定比例的无错句，模型才会学到“不改”也是正常输出。

\subsubsection{条件序列生成纠错}
\label{sec:nlp-generation-seq-correction}

把纠错当作从错句到正句的翻译，就可以直接使用编码器—解码器整句生成正句。
这条路线能一次写出任意修改，包括长短语插入和重排；问题在于正句的绝大多数词元与源句相同，
逐词元交叉熵大部分花在“照抄”上，真正需要学的修改只占很小比例，模型因此偏向不改。
Junczys-Dowmunt等把纠错看作低资源翻译，提出\term{编辑加权最大似然}
（edit-weighted maximum likelihood），提高被修改词元的损失权重\scite{junczys2018approaching}
\begin{equation}
  L(\bm\theta)=-\sum_{t=1}^{T}\lambda_t\log p_{\bm\theta}(y_t\mid x,y_{<t}),\qquad
  \lambda_t=\begin{cases}\Lambda,& y_t\in\mathcal{M},\\
  1,&y_t\notin\mathcal{M},\end{cases}
  \label{eq:nlp-generation-edit-weighted}
\end{equation}
其中$\mathcal{M}$是目标中被修改的词元，即按式~\eqref{eq:nlp-generation-edit-distance}
与源句对齐后不是原样保留的位置，$\Lambda>1$。
它是式~\eqref{eq:nlp-generation-nll}的加权版本：$\Lambda=1$时退化为普通最大似然，
$\Lambda$越大，模型越敢于修改，精确率也随之下降。
另一种做法从结构上处理“照抄”：Zhao等在Transformer上加入第~\ref{sec:nlp-generation-copy-pointer}小节
的复制分布，并用去噪方式在无标注文本上预训练，使模型把复制当作默认动作，
把生成留给真正需要修改的位置\scite{zhao2019improving}
预训练的BART、T5类编码器—解码器本身就在去噪中学会了复制未损坏部分，
微调到纠错时也继承了这一偏置；其去噪目标见第~\ref{sec:nlp-llm-objective-t5-bart}小节。

整句生成的输出要重新与源句对齐，才能知道模型改了哪里。固定输入
“She go to meeting on Friday”，参考修改为go到goes，以及在to后插入the。
生成候选“She goes to the meeting on Monday”与源句对齐后包含三项编辑：
前两项正确，Friday到Monday是过改，并改变了事实。正确编辑2项、系统编辑3项、
参考编辑2项，精确率为$2/3$，召回率为1。整句生成可以写出任何修改，
也就要为输出中每一处非必要的变化负责。

\subsubsection{GECToR编辑标签纠错}
\label{sec:nlp-generation-gector}

整句生成要逐词重写，推断慢，也给过改留下空间。GECToR沿LaserTagger的思路，
把纠错写成对每个源词元的标签预测，并针对纠错扩充了标签集和推断控制
\scite{omelianchuk-etal-2020-gector}

\paragraph{编码器与两个输出头}
输入是在句首加入特殊位置\code{\$START}的源词元序列，句首位置使得在第一个词之前插入成为可能。
预训练编码器（BERT、RoBERTa或XLNet）给出各位置的向量；一个词被切成多个子词时，
取第一个子词的向量$\bm h_i\in\mathbb{R}^{d}$，基础规模模型$d=768$。
编码器上接两个线性层加softmax：检测头输出该词是否有错的二分类概率
$p_{\mathrm{det}}(\cdot\mid i)=\operatorname{softmax}(\bm W_{\mathrm{d}}\bm h_i+\bm b_{\mathrm{d}})$；
标签头输出变换标签的分布
$p_{\mathrm{tag}}(\cdot\mid i)=\operatorname{softmax}(\bm W_{\mathrm{t}}\bm h_i+\bm b_{\mathrm{t}})$，
$\bm W_{\mathrm{t}}\in\mathbb{R}^{|T|\times d}$。默认标签集共5000个：
4971个基本变换，即\code{\$KEEP}、\code{\$DELETE}和带具体词的\code{\$APPEND\_w}、\code{\$REPLACE\_w}；
另有29个与具体词无关的\term{g变换}（g-transformation），如大小写变换、相邻词合并与拆分、
名词单复数和动词形式变换（\code{\$VERB\_FORM\_VB\_VBZ}把go变成goes）。
g变换让一个标签覆盖一整类词形修改，不必为每个词各设一个替换标签。

训练标签由Levenshtein对齐的一个变体从平行句对导出；一个源词元对应多个修改时，
每轮只保留第一个非\code{\$KEEP}标签，其余修改留给下一轮。检测头的标签由此确定：
标签不是\code{\$KEEP}的位置为“错”。两个头的交叉熵相加作为损失：
\begin{equation}
  L=-\sum_{i=0}^{n}\Bigl[\log p_{\mathrm{tag}}(t_i\mid i)
  +\log p_{\mathrm{det}}(e_i\mid i)\Bigr],\qquad e_i=\mathbb{I}[t_i\neq\code{\$KEEP}].
  \label{eq:nlp-generation-gector-loss}
\end{equation}
检测头与标签头共享编码器，检测头给出的是“这里是否需要动”的整体判断，
推断时可以用来决定整句是否值得修改。

\begin{figure*}[htbp]
  \centering
  % GECToR一轮推断：源词元（含句首位置）经编码器得到逐词向量，检测头与标签头逐位置输出，
% 加置信度偏置后取标签并同时应用，结果进入下一轮。概率为教学示例，与正文算例一致。
\begin{tikzpicture}[
  x=1mm,y=1mm,figure picture,
  tok/.style={figure node,minimum width=17mm,minimum height=6mm,
    inner sep=1pt,draw=FigureInput,fill=FigureInputFill,font=\figurefont\scriptsize},
  enc/.style={figure node,minimum height=7mm,inner sep=1.5pt,
    draw=FigureModel,fill=FigureModelFill,font=\figurefont\scriptsize},
  vec/.style={font=\figurefont\scriptsize},
  outc/.style={figure node,minimum width=19.5mm,minimum height=10mm,align=center,
    inner sep=1pt,draw=FigureOutput,fill=FigureOutputFill,font=\figurefont\scriptsize},
  res/.style={figure node,minimum height=6.5mm,inner sep=1.5pt,
    draw=FigureOutput,fill=FigureOutputFill,font=\figurefont\scriptsize},
  note/.style={font=\figurefont\scriptsize,text=FigureMuted,align=center},
  row label/.style={font=\figurefont\scriptsize,text=FigureMuted,anchor=east,align=right},
  flow/.style={figure flow}
]
  \foreach \i/\w/\tag/\pe in {0/{\$START}/KEEP/0.02,1/She/KEEP/0.03,2/go/{VB$\to$VBZ}/0.85,
      3/to/{APPEND\_the}/0.70,4/meeting/KEEP/0.45}{
    \node[tok] (x\i) at (\i*22,0) {\w};
    \node[vec] (h\i) at (\i*22,22) {$\bm h_{\i}$};
    \node[outc] (t\i) at (\i*22,36) {标签：\tag\\有错概率\pe};
    \draw[flow] (x\i) -- (x\i |- 0,8.5);
    \draw[flow] (\i*22,15.5) -- (h\i);
    \draw[flow] (h\i) -- (t\i);
  }
  \node[enc,minimum width=108mm] (enc) at (44,12)
    {预训练编码器（每个词取首个子词的向量）};
  \node[row label] at (-12,0) {本轮输入};
  \node[row label] at (-12,22) {逐词向量};
  \node[row label] at (-12,36) {标签头\\检测头};

  \node[res,minimum width=60mm] (y) at (44,54) {She goes to the meeting};
  \draw[flow] (t2.north) -- (t2.north |- y.south);
  \draw[flow] (t3.north) -- (t3.north |- y.south);
  \node[note,anchor=west] at (76,54) {同时应用，\\作为下一轮输入};
  \node[row label] at (-12,54) {重建结果};
  \node[note] at (88,44.5) {加偏置$\beta$后KEEP胜出};
\end{tikzpicture}

  \caption{GECToR的一轮推断。每个源词元（含句首位置）取编码器向量，检测头给出有错概率，
  标签头给出变换标签；标签按本轮的原坐标同时应用，得到的新句作为下一轮输入，
  直到不再产生非KEEP标签或达到轮数上限。图中概率为教学示例：meeting处的
  复数变换概率略高于KEEP，加入置信度偏置后保持原词。}
  \label{fig:nlp-generation-gector}
\end{figure*}

\paragraph{推断控制与迭代}
标签头在每个位置给出分布后，GECToR用两个超参数控制修改的积极程度。
\term{置信度偏置}（confidence bias）$\beta\geq0$加到\code{\$KEEP}的概率上，
使模型只在足够确信时才修改；\term{最小错误概率}（minimum error probability）$\delta$
作用于整句：若检测头在所有位置给出的最大有错概率$\max_i p_{\mathrm{det}}(\text{错}\mid i)$
低于$\delta$，整句不改。两者都在开发集上调节，用以权衡精确率与召回率。
如图~\ref{fig:nlp-generation-gector}，设go处
$p(\code{\$VERB\_FORM\_VB\_VBZ})=0.70$、$p(\code{\$KEEP})=0.20$，
meeting处$p(\code{\$NOUN\_NUMBER\_PLURAL})=0.48$、$p(\code{\$KEEP})=0.42$。
取$\beta=0.10$后，go处\code{\$KEEP}升到$0.30$，仍执行词形变换；
meeting处\code{\$KEEP}升到$0.52$，超过复数变换，避免了“meetings”这一过改。

修改之间可能相互依赖，一轮未必能完成。例如“She go to meeting”的目标若是
“She goes to the annual meeting”，每个位置一轮只能追加一个词，
第一轮在to之后追加the，第二轮在新序列的the之后追加annual。
算法~\ref{alg:nlp-generation-gector}给出迭代过程。

\begin{algorithm}[GECToR的迭代纠错]
  \label{alg:nlp-generation-gector}
  \begin{algorithmic}[1]
    \Require 源句$x$；置信度偏置$\beta$，最小错误概率$\delta$，最大轮数$K$
    \Ensure 修订句
    \For{$k=1,\ldots,K$}
      \State 在$x$前加\code{\$START}，编码得到$\bm h_0,\ldots,\bm h_n$
      \State 计算各位置的$p_{\mathrm{det}}$与$p_{\mathrm{tag}}$；
        令$p_{\mathrm{tag}}(\code{\$KEEP}\mid i)\gets p_{\mathrm{tag}}(\code{\$KEEP}\mid i)+\beta$
      \If{$\max_i p_{\mathrm{det}}(\text{错}\mid i)<\delta$}
        \State \Return $x$
      \EndIf
      \State $t_i\gets\argmax_t p_{\mathrm{tag}}(t\mid i)$，$i=0,\ldots,n$
      \If{所有$t_i$都是\code{\$KEEP}}
        \State \Return $x$
      \EndIf
      \State 以本轮$x$的坐标同时应用全部$t_i$，得到新的$x$\Comment{下一轮重新编号}
    \EndFor
    \State \Return $x$
  \end{algorithmic}
\end{algorithm}

一轮中的所有标签都以同一份原句为坐标，从左到右重建，
不能一边插入一边仍按已移动的位置执行后面的动作。原文报告多数修改在前两轮完成，
限制轮数能以少量精度换取速度。

\paragraph{三阶段训练}
GECToR按数据质量分三阶段训练：第一阶段在约九百万句的合成错误数据上预训练；
第二阶段只用真实标注语料中含错误的句子微调；第三阶段在规模较小、质量较高的语料上
用含错句和无错句的混合数据微调。第二阶段集中学习修改，第三阶段把无错句加回来，
让模型学会对正确句不动手，这与上一小节保留无错句的理由相同。
前两个阶段的最初两轮还冻结编码器，只训练两个输出头，避免随机初始化的头把编码器带偏。

\paragraph{与整句生成的比较}
与条件序列生成相比，GECToR的每个位置独立预测，一轮只需一次编码器前向；
原文在批量推断下测得其速度最高可达束宽12的Transformer纠错模型\cite{zhao2019improving}的十倍。
编辑可见，可以逐项接受或拒绝。代价同LaserTagger：标签集决定了可表达的修改，
未收录的替换、长短语插入和大幅重排需要多轮，甚至无法完成。
它的g变换是为英语形态设计的；用于中文时，需要重新规定按字还是按词对齐、
同音字替换动作和词语边界，而不能把单复数、时态标签原样搬过去。

\subsubsection{大语言模型纠错}
\label{sec:nlp-generation-llm-gec}

大语言模型可以直接按指令纠错：给出错句和“修正其中的语法错误”的要求，
或在纠错平行语料上微调。它的优势是无需专门的标签集，能处理搭配、语序与上下文相关的错误，
输出也更流畅。Fang等对ChatGPT的系统评测发现，它在允许流畅改写的JFLEG上表现突出，
在以最小编辑为参考的数据集上自动分数却明显低于专门的纠错系统；
人工分析显示其主要问题是\term{过度修改}（overcorrection）：
把语法正确但不够地道的部分也改写，与最小编辑参考对不上\scite{fang2023chatgpt}

过度修改既是评价口径问题，也是实际风险：改写可能改动实体、数字或语气。
控制方法有两类：一是在提示中明确要求“只修正语法错误，尽量少改，不改动专名和数字”，
并给出最小编辑的示例；二是事后过滤：用式~\eqref{eq:nlp-generation-edit-distance}
把输出与源句对齐，恢复编辑集合，删去落在保护片段上的编辑和不属于目标错误类型的编辑，
再把保留的编辑应用到源句。后一类把大语言模型当作候选编辑的提供者，
最终修改仍可逐项检查。评价时同时报告最小编辑与流畅性两类参考上的结果，
具体指标见第~\ref{sec:nlp-generation-edit-metrics}小节。

\subsection{文本简化}
\label{sec:nlp-generation-simplification}

\term{文本简化}（text simplification）要为指定读者降低理解负担，
先确定必须保留的信息和可删的背景，再通过词汇替换、拆句、删除和重排实现。
它的挑战是：平行简化语料规模有限，常见的维基百科与简明维基百科句对对齐还有噪声；
“更短”“用更常见的词”只是代理属性，删得越多越短，却可能删掉关键条件；
不同读者需要的简化程度不同，一个模型只学到训练语料的平均程度。
ASSET的多参考简化同时包含词汇替换、拆句和删减，提醒评价不能只奖励缩短
\scite{nlp-intro-alvamanchego2020}
方法谱系为：规则与统计翻译式简化；以简化指标为奖励的强化学习；
用控制属性显式调节简化程度；从无标注语料挖掘复述实现无监督简化。
自动指标SARI见第~\ref{sec:nlp-generation-sari}小节，事实保持与读者理解则需要另行检查。

\subsubsection{词典替换、拆句与规则重排}
\label{sec:nlp-generation-rule-simplification}

规则路线从适合读者的词典中取常用表达替换难词，再利用句法结构确定拆句位置和省略成分。
例如把“须”解释为“必须”，而不替换成义务更弱的“可以”；拆开条件句时，把适用人群带到新句中；
代词只有在指代唯一时才换成明确名词。句法与指代分析见第~\ref{chap:nlp-text-analysis}章。

以带否定与条件的长句为例：
“若您是校外来宾，须在4月16日18:00前提交报名信息；这项要求不适用于校内师生。”
分析得到第一分句的条件是“校外”，主干是“须提交”；第二分句的“这项要求”指提前提交报名信息。
替换“须”、把条件改为显式主语并拆句，得到
“校外来宾必须在4月16日18:00前提交报名信息。校内师生不受这项要求限制。”
若改成“请在4月16日18:00前提交报名信息。校内师生不受限制”，
第一句丢掉了人群范围，第二句把一项要求扩大成所有限制。拆句减少了句内依赖，
却增加了跨句指代；保留原指代比补错指代更安全。

规则可审查，但覆盖有限。统计翻译式简化把复杂句到简单句看作单语翻译：
Wubben等在维基百科—简明维基百科句对上训练短语统计翻译模型PBMT-R
\scite{wubben2012sentence}
简化语料中大量句对几乎相同，模型学到的最优输出常常就是原句；
他们因此在模型分数最高的$N$个候选$\mathcal{N}_N(x)$中，
选与原句差异最大的那一个：
\begin{equation}
  \hat y=\argmax_{y\in\mathcal{N}_N(x)}d(x,y),
  \label{eq:nlp-generation-pbmtr}
\end{equation}
$d$为词级编辑距离。这一重排强迫模型做出修改，但“与原句差异大”并不等于“更简单”，
可能选中删错或改错的候选。下面的DRESS把“更简单”本身写进训练目标。

\subsubsection{DRESS强化学习简化}
\label{sec:nlp-generation-dress}

最大似然训练让模型模仿参考简化句，而参考中大量是原样复制，
模型学到的往往是保守的复制。Zhang与Lapata提出DRESS，
直接以简化质量为奖励训练生成模型\scite{nlp-generation-zhang2017dress}
策略是一个带注意力的LSTM编码器—解码器$p_{\bm\theta}(y\mid x)$，
先用最大似然预训练，再从中采样完整输出$\hat y$，按三项奖励的加权和打分：
\begin{equation}
  r(\hat y)=\lambda_{\mathrm{S}}\operatorname{SARI}(x,\hat y,y^{\ast})
  +\lambda_{\mathrm{R}}\cos(\bm q_x,\bm q_{\hat y})
  +\lambda_{\mathrm{F}}\exp\Bigl(\frac{1}{|\hat y|}\sum_{t}\log p_{\mathrm{LM}}(\hat y_t\mid\hat y_{<t})\Bigr).
  \label{eq:nlp-generation-dress-reward}
\end{equation}
第一项是简化指标SARI，它同时比较源句、输出和参考，奖励正确的增加、保留与删除；
第二项是意义相关性，$\bm q_x,\bm q_{\hat y}$为一个预训练序列自编码器给出的句向量；
第三项是流畅度，即简单语料上训练的语言模型给出的逐词元概率的几何平均，
也就是困惑度的倒数。三项分别对应简单、保意和通顺，单独优化任一项都会退化：
只看SARI可能删掉关键信息，只看相关性会原样复制，只看流畅度会生成短而空洞的句子。

参数用REINFORCE加基线更新，梯度为
\[
  \nabla_{\bm\theta}J\approx\bigl(r(\hat y)-b(x)\bigr)\nabla_{\bm\theta}\log p_{\bm\theta}(\hat y\mid x),
\]
基线$b(x)$由解码器隐状态的线性回归估计，用于降低梯度方差。
这正是第~\ref{sec:nlp-generation-exposure-bias}小节序列级训练的一个实例：
训练时模型以自己的采样为前缀，奖励直接对应任务目标。
设教学权重$\lambda_{\mathrm{S}}=1$、$\lambda_{\mathrm{R}}=\lambda_{\mathrm{F}}=0.5$，
某个采样的SARI为0.40、相关性0.80、流畅度0.50，则$r=0.40+0.40+0.25=1.05$；
基线为0.95时优势为0.10，这条采样的对数概率被小幅提高。
另一条采样删掉了截止时间，SARI升到0.45，相关性却降到0.50，
$r=0.45+0.25+0.25=0.95$，优势为零，不受鼓励。

DRESS直接对准简化目标，但三项权重需要在开发集上调节，奖励函数本身也只是代理：
SARI依赖参考，相关性与流畅度来自各自的模型，强化学习可能找到骗过它们的输出。

\subsubsection{ACCESS受控简化}
\label{sec:nlp-generation-access}

DRESS学到的简化程度是训练数据与奖励权重决定的平均值，用户无法指定“保留信息但多拆句”。
ACCESS把目标属性编码为输入前的控制标记，让同一个模型按指定程度简化
\scite{martin-etal-2020-controllable}
它使用四类可由句对直接计算的属性：字符长度比、字符级编辑相似度、
词汇复杂度和依存树深度。训练时对每个句对测出这四个值，离散化后作为前缀
拼到原句前，模型学习“在这些属性下，原句应怎样变成简化句”。

记原句与目标的字符数为$n_x,n_y$，长度比为$n_y/n_x$；编辑相似度为
$1-d_{\mathrm{char}}(x,y)/\max(n_x,n_y)$，值越大表示改动越少；
词汇复杂度先取各词词频排名的对数，用第三四分位数汇总，再取目标与源之比；
句法复杂度取依存分析得到的最大依存深度，本节统一为目标深度与源深度之比。
例如一对教学句经固定预处理后有40与36个字符、字符编辑距离12，
对数词频排名的第三四分位数为5与4，最大依存深度为5与3，
四个控制值便是$0.90,0.70,0.80,0.60$。按0.05为间隔离散并把比值上限设为2，
把对应的四个控制标记接在原句前，以真实简化句为目标训练条件最大似然。

推断时目标句未知，控制值由用户或开发集选定，例如要求长度比0.8、编辑相似度0.6，
模型便倾向于输出更短、改动更多的句子。控制标记是条件而不是硬约束：
生成后应重新计算实际属性，检查是否接近请求值；属性达标也不说明截止时间、
义务和人群都已保留，这些仍要逐项核对。

\subsubsection{MUSS挖掘复述的无监督简化}
\label{sec:nlp-generation-muss}

ACCESS仍需平行简化语料，而这类语料在英语之外很少。MUSS的观察是：
在大量复述对中，有的方向是变简单，有的是变复杂，长度和词频都在变化；
若用ACCESS的控制标记把每对的属性变化写进输入，模型就能从普通复述中学会按控制值改写，
推断时把控制值设到“更简单”一侧即可\scite{martin2022muss}
这样，简化所需的监督换成了更容易获得的复述。

复述对从网络语料中挖掘：把大规模单语网页文本切成句子或短句序列，
用多语句向量模型LASER编码，以近似最近邻检索为每个序列找最相近的若干序列，
保留距离低于阈值的对；再过滤掉几乎相同的对和一方包含另一方的对，
这些对没有改写信息。挖掘结果按ACCESS的方式计算控制值，在其上微调预训练的BART
（多语时用mBART）。推断时控制值在开发集上选择。
这一做法把第~\ref{sec:nlp-generation-backtranslation-paraphrase}小节的复述构造从翻译换成检索，
也是第~\ref{chap:nlp-classification-matching-organization}章向量检索在数据构造中的用法。

MUSS不需要简化标注，可用于标注稀缺的语言；代价是挖掘与过滤的计算量、
仍需开发集或人工设定的控制值，以及挖掘到的复述本身可能改变意义，这类错误会传给模型。

\paragraph{文本编辑方法的对照}
表~\ref{tab:nlp-generation-editing-methods}按四个任务汇总本节方法。
纵向看，输出表示从局部替换、整句生成到编辑标签，再到提示，
越靠近整句生成，可表达的修改越多，过改的风险也越大；
编辑标签和选择式方法则以有限的表达能力换取速度和可检查性。
横向看，平行数据不足是共同问题，回译、合成错误、属性语料和挖掘复述
是四种替代监督，各自带来分布偏移或意义漂移，需要过滤与核验配合。

\begin{table*}[htbp]
  \centering
  \small
  \begin{tabularx}{\linewidth}{@{}>{\raggedright\arraybackslash}p{24mm}
    >{\raggedright\arraybackslash}p{30mm}>{\raggedright\arraybackslash}p{38mm}
    >{\raggedright\arraybackslash}p{30mm}>{\raggedright\arraybackslash}X@{}}
    \toprule
    方法 & 解决的问题 & 核心机制 & 数据需求 & 新代价与风险\\
    \midrule
    回译复述 & 复述对稀缺 & 经枢轴语言译回，按相似度过滤 & 翻译系统或双语语料 & 意义漂移只能事后过滤\\
    LaserTagger & 整句重写慢、易补造 & KEEP/DELETE加词表内插入短语 & 平行句对 & 词表外插入与大幅重排无法表达\\
    DRG & 无平行风格数据 & 删属性标记、检索、重建 & 按属性分开的语料 & 依赖属性可由局部短语承载\\
    交叉对齐、对抗解纠缠 & 风格不局限于短语 & 对抗对齐隐表示或隐状态 & 属性语料 & 解纠缠难以验证，训练不稳\\
    去噪加回译 & 不必解纠缠 & 迭代生成伪平行对 & 属性语料 & 伪数据错误逐轮累积\\
    提示改写 & 每种风格都要训练 & 示例加风格描述，候选筛选 & 无需训练 & 推断成本高，过改需筛选\\
    混淆集打分 & 同音形近错误 & 局部候选加语言分数减编辑代价 & 混淆集与语言模型 & 覆盖受限于候选集合\\
    合成错误数据 & 平行纠错数据不足 & 规则加噪或反向模型生成错句 & 干净文本 & 合成错误分布偏移\\
    序列生成纠错 & 任意形式的修改 & 编码器—解码器，编辑加权或复制 & 平行数据加合成数据 & 推断慢，需对齐恢复编辑\\
    GECToR & 推断慢，过改 & 检测头加标签头，偏置与阈值，迭代 & 三阶段数据 & 标签集限制表达，英语专用变换\\
    大模型纠错 & 专门标签与训练成本 & 提示或微调，事后过滤编辑 & 少量或无 & 过度修改\\
    PBMT-R & 规则覆盖有限 & 单语短语翻译，按差异重排 & 平行简化语料 & 差异大不等于更简单\\
    DRESS & 最大似然偏向复制 & 以SARI、相关性、流畅度为奖励 & 平行简化语料 & 奖励权重需调，代理可被利用\\
    ACCESS & 简化程度不可控 & 四类属性的离散控制标记 & 平行简化语料 & 控制值是条件而非约束\\
    MUSS & 简化语料稀缺 & 挖掘复述加控制标记 & 无标注单语语料 & 挖掘成本，复述本身可能漂移\\
    \bottomrule
  \end{tabularx}
  \caption{文本编辑方法对照。“数据需求”指训练该方法所需的主要监督来源；
  各方法的输出都需要按任务另行核验保护片段与意义是否保持。}
  \label{tab:nlp-generation-editing-methods}
\end{table*}
\FloatBarrier

\section{摘要与信息综合}
\label{sec:nlp-generation-summarization}

\term{文本摘要}（text summarization）要在长度预算内保留与用途相关的信息。
日历条目主要保留时间和地点，参加指南必须保留行动条件，会议纪要还须分清建议、
已作决定和待办事项。摘要面临的挑战依次是：从大量内容中选出显著信息；
在句内压缩与跨句融合时改写表达；去除重复；保持对来源的忠实，不补造事实；
处理超出模型长度上限的输入；多来源时还要处理转载、互补与冲突。

方法谱系大致沿时间展开，每一步回应前一步留下的问题。
位置基线与图中心性从原文选句，事实来自原句，却不能压缩句内信息；
有监督句子抽取学习哪些句子重要，需要先从人工摘要构造句子标签；
注意力编码器—解码器能写出原文没有的句子，却难以复制名称和数字，也容易重复；
复制与覆盖机制补上这两点；去噪预训练先在无标注文本上学会重建，再微调到摘要；
长文档结构与稀疏注意力把输入推到数万词元；人类反馈把训练目标从模仿参考换成人类偏好；
多来源综合在选择、融合与预训练三个层面处理多份材料。节末
表~\ref{tab:nlp-generation-summarization-methods}比较各法的忠实性、压缩能力、输入长度与训练数据。

\subsection{抽取式摘要}
\label{sec:nlp-generation-extractive-summary}

\term{抽取式摘要}（extractive summarization）从原文选句组成摘要。
设文档有$n$个句子，第$i$句的长度为$c_i$，抽取式摘要可以写成在预算$B$下选择句子子集
\begin{equation}
  S^{\ast}=\argmax_{S\subseteq\{1,\ldots,n\},\;\sum_{i\in S}c_i\leq B}\;
  Q(S),
  \label{eq:nlp-generation-extractive-objective}
\end{equation}
其中$Q$衡量所选句子的整体质量。无监督方法用位置或图中心性近似$Q$；
有监督方法把问题简化为对每句做二分类，再按预测概率在预算内选句。
被选的句子都是原句，事实不会被改写错，但句内的冗余无法删去，被选句也可能因指代脱离上下文
而误导：单独摘出“这项要求不适用于校内师生”却不带上“这项要求”所指的内容，正确原句也成了
误导性的摘要。

\subsubsection{Lead与TextRank抽取摘要}
\label{sec:nlp-generation-extractive}

Lead直接按原文顺序取前若干句，直到句数或词元预算用尽。它不需要训练样本，
在新闻这类重要信息位于开头的文体上是很强的基线，也是其他方法必须超过的参照。
TextRank改用句子之间的联系估计重要性：把各句作为节点，用词重合或其他相似度为边赋权，
在图上迭代传播分数，与其他句子联系越多、越强的句子得分越高\scite{nlp-mihalcea2004textrank}

设$w_{i,j}\geq0$为不同句子间的对称边权，自环权为零，分数归一化为总和1：
\begin{equation}
  r_i^{(k+1)}
  =\frac{1-d}{n}
   +d\sum_{j:\,w_{j,i}>0}
   \frac{w_{j,i}}{\sum_\ell w_{j,\ell}}r_j^{(k)}.
  \label{eq:nlp-generation-textrank}
\end{equation}
$d$是沿图传播的权重，$1-d$的部分均匀跳转到任一句，保证迭代收敛；
没有出边的句子按均匀分布分配其质量。从均匀分数开始迭代，
到相邻两轮差小于阈值为止。按分数选完句子后，通常恢复原文次序，以免打乱条件和指代。

将通知整理为三句教学输入：$s_1$“读书会改至4月18日14:00举行”，
$s_2$“读书会4月18日在海棠楼201举行”，
$s_3$“校外来宾参加读书会须在4月16日18:00前提交报名信息”。
只保留检索词集合$K_1=\{\text{读书会},\text{4月18日},\text{14:00}\}$、
$K_2=\{\text{读书会},\text{4月18日},\text{海棠楼201}\}$、
$K_3=\{\text{读书会},\text{校外来宾},\text{报名信息}\}$，
取$w_{i,j}=|K_i\cap K_j|$，得$w_{1,2}=2$、$w_{1,3}=w_{2,3}=1$。
令$d=0.85$，从$(1/3,1/3,1/3)$开始，第一轮得$(0.38056,0.38056,0.23889)$，
收敛到$(0.37013,0.37013,0.25974)$。预算为两句、平局取靠前句时，
Lead和TextRank都选择$s_1,s_2$，重复了日期却遗漏校外要求。
中心性奖励与其他句子相似的句子，而最独特的条件信息恰恰与其他句子相似度低；
若查询是“校外来宾如何办理”，就需要把查询相关性写进分数，
或用第~\ref{sec:nlp-generation-mmr}小节的方法显式去冗余。

\subsubsection{有监督句子抽取：SummaRuNNer与BertSum}
\label{sec:nlp-generation-supervised-extractive}

Lead和TextRank的打分规则是人为指定的。有了大规模新闻摘要语料，就可以让模型学习
什么样的句子应当入选。困难在于这些语料提供的是人工写的生成式摘要，而不是句子标签。

\paragraph{贪心构造oracle标签}
常用做法是用ROUGE（第~\ref{sec:nlp-generation-rouge}小节）衡量句子组合与人工摘要的重合，
贪心地构造一个近似最优的抽取结果，作为\term{oracle摘要}（oracle summary），
其中的句子标为1、其余标为0\scite{nallapati2017summarunner}

\begin{algorithm}[贪心构造抽取式oracle标签]
  \label{alg:nlp-generation-greedy-oracle}
  \begin{algorithmic}[1]
    \Require 文档句子$s_1,\ldots,s_n$，参考摘要$y^{\ast}$，评分函数$\operatorname{R}$（如ROUGE-1与ROUGE-2的F值之和）
    \Ensure 句子标签$z_1,\ldots,z_n\in\{0,1\}$
    \State $S\gets\varnothing$，$\mathrm{best}\gets0$
    \Loop
      \State $i^{\ast}\gets\argmax_{i\notin S}\operatorname{R}(S\cup\{i\},y^{\ast})$
      \If{$\operatorname{R}(S\cup\{i^{\ast}\},y^{\ast})\leq\mathrm{best}$}
        \State \textbf{break}\Comment{再加任何一句都不能提高分数}
      \EndIf
      \State $S\gets S\cup\{i^{\ast}\}$；$\mathrm{best}\gets\operatorname{R}(S,y^{\ast})$
    \EndLoop
    \State $z_i\gets\mathbb{I}[i\in S]$，$i=1,\ldots,n$
  \end{algorithmic}
\end{algorithm}

用词级ROUGE-1的F值手算一例。教学参考摘要与四个句子的词如下：
\begin{center}
  \small
  \begin{tabular}{@{}ll@{}}
    \toprule
    对象 & 词\\
    \midrule
    参考$y^{\ast}$ & 读书会／4月18日／14:00／海棠楼201／校外来宾／报名\\
    $s_1$ & 读书会／改期／4月18日／14:00\\
    $s_2$ & 读书会／地点／海棠楼201\\
    $s_3$ & 校外来宾／4月16日／报名\\
    $s_4$ & 读书会／欢迎／参加\\
    \bottomrule
  \end{tabular}
\end{center}
ROUGE-1的F值可写成$2\times\text{匹配词数}/(\text{摘要词数}+\text{参考词数})$，
匹配时每个参考词至多计一次。
单句的F值分别为$0.600,0.444,0.444,0.222$，先选$s_1$。
第二轮，$s_1{+}s_2$共7词、与参考匹配4个（“读书会”只计一次），F值为$8/13\approx0.615$；
$s_1{+}s_3$匹配5个，F值为$10/13\approx0.769$；$s_1{+}s_4$为$6/13\approx0.462$；选$s_3$。
第三轮，再加$s_2$虽把召回率提到1，精确率却降到0.6，F值为0.75，低于0.769，
算法停止。oracle为$\{s_1,s_3\}$，含地点的$s_2$被标为0。
可见oracle标签只是近似：它依赖评分函数和贪心次序，参考中出现的信息也可能因精确率惩罚而落选。

\paragraph{SummaRuNNer}
SummaRuNNer用两层双向GRU编码文档\scite{nallapati2017summarunner}
词级GRU读入每句的词向量，句内各位置隐状态的平均作为句子输入；
句级GRU在句子序列上运行，得到第$j$句的表示$\bm h_j\in\mathbb{R}^{2d}$（前后两向拼接）。
文档表示为$\bm d=\tanh\bigl(\bm W_d\frac{1}{N}\sum_{j=1}^{N}\bm h_j+\bm b_d\bigr)$。
第$j$句的入选概率由五项相加后经逻辑函数得到：
\begin{equation}
  P(z_j=1)=\sigma\Bigl(\underbrace{\bm w_c^{\top}\bm h_j}_{\text{内容}}
  +\underbrace{\bm h_j^{\top}\bm W_s\bm d}_{\text{显著性}}
  -\underbrace{\bm h_j^{\top}\bm W_r\tanh(\bm s_j)}_{\text{冗余}}
  +\underbrace{\bm w_a^{\top}\bm p^{\mathrm{abs}}_j+\bm w_p^{\top}\bm p^{\mathrm{rel}}_j}_{\text{位置}}
  +b\Bigr),
  \label{eq:nlp-generation-summarunner}
\end{equation}
其中$\bm s_j=\sum_{i<j}\bm h_iP(z_i=1)$是到第$j$句为止的“已选摘要”表示，按各句入选概率加权；
$\bm p^{\mathrm{abs}}_j$与$\bm p^{\mathrm{rel}}_j$是句子绝对位置与相对位置（文档分段后的段号）的嵌入。
内容项衡量句子本身的信息量，显著性项衡量它与全文主旨的关联，冗余项扣除它与已选内容的重合，
位置项吸收新闻开头更重要的先验。$\bm W_s,\bm W_r\in\mathbb{R}^{2d\times2d}$为双线性参数。
训练最小化oracle标签上的二元交叉熵，推断时按概率从高到低选句，直到预算用尽。
冗余项让抽取本身带有类似MMR的去冗余，但它是从oracle标签中间接学到的，不是显式约束。
式~\eqref{eq:nlp-generation-summarunner}的各项都可解释，这是它相对黑盒分类器的优点。

\paragraph{BertSum}
预训练编码器出现后，句子表示可以直接来自BERT。困难在于BERT的输入原本只有一个或两个句段，
其\code{[CLS]}向量表示整个输入。Liu与Lapata的BertSum改造了输入格式
\scite{liu2019text}：在每个句子前插入\code{[CLS]}、句后插入\code{[SEP]}；
\term{区间分段嵌入}（interval segment embedding）按句子的奇偶交替使用$E_A$和$E_B$，
使模型能区分相邻句子；第$i$个\code{[CLS]}位置的顶层向量$\bm t_i\in\mathbb{R}^{768}$作为第$i$句的表示。
图~\ref{fig:nlp-generation-bertsum}画出这一结构。

\begin{figure*}[htbp]
  \centering
  % BertSum抽取结构：每句前加[CLS]、后加[SEP]，分段嵌入按句子奇偶交替；
% 各[CLS]顶层向量经句间Transformer层与逻辑回归给出入选概率。
\begin{tikzpicture}[
  x=1mm,y=1mm,figure picture,
  tok/.style={figure node,minimum width=10mm,minimum height=6mm,
    inner sep=1pt,draw=FigureInput,fill=FigureInputFill,font=\figurefont\scriptsize},
  seg/.style={font=\figurefont\scriptsize,text=FigureMuted},
  enc/.style={figure node,minimum height=7mm,inner sep=1.5pt,
    draw=FigureModel,fill=FigureModelFill,font=\figurefont\scriptsize},
  vec/.style={font=\figurefont\scriptsize},
  outc/.style={figure node,minimum width=12mm,minimum height=6mm,
    inner sep=1pt,draw=FigureOutput,fill=FigureOutputFill,font=\figurefont\scriptsize},
  row label/.style={font=\figurefont\scriptsize,text=FigureMuted,anchor=east},
  flow/.style={figure flow}
]
  % 三个句子，每句三个位置：[CLS]、句子词元、[SEP]
  \foreach \k/\s/\e in {1/{$s_1$词元}/E_A,2/{$s_2$词元}/E_B,3/{$s_3$词元}/E_A}{
    \pgfmathsetmacro{\xo}{(\k-1)*36}
    \node[tok] (c\k) at (\xo,0) {[CLS]};
    \node[tok,minimum width=12mm] (w\k) at (\xo+12.5,0) {\s};
    \node[tok] (p\k) at (\xo+25,0) {[SEP]};
    \node[seg] at (\xo+12.5,-6) {分段$\e$};
    \draw[flow] (c\k) -- (c\k |- 0,8.5);
    \draw[flow] (w\k) -- (w\k |- 0,8.5);
    \draw[flow] (p\k) -- (p\k |- 0,8.5);
    \node[vec] (t\k) at (\xo,23) {$\bm t_{\k}$};
    \draw[flow] (\xo,15.5) -- (t\k);
    \node[outc] (z\k) at (\xo,48) {$\hat z_{\k}$};
  }
  \node[enc,minimum width=108mm] (bert) at (48.5,12)
    {BERT编码器（词元嵌入＋分段嵌入＋位置嵌入）};
  \node[enc,minimum width=84mm] (inter) at (36,35)
    {句间Transformer层（$L=2$），输入加句子位置嵌入};
  \foreach \k in {1,2,3}{
    \pgfmathsetmacro{\xo}{(\k-1)*36}
    \draw[flow] (t\k) -- (\xo,31.5);
    \draw[flow] (\xo,38.5) -- (z\k);
  }
  \node[row label] at (-7,0) {输入};
  \node[row label] at (-7,23) {句向量};
  \node[row label] at (-7,48) {$\sigma(\bm w_o^{\top}\bm h_i+b_o)$};
\end{tikzpicture}

  \caption{BertSum的抽取结构。每句前插入\code{[CLS]}、句后插入\code{[SEP]}，分段嵌入按句子奇偶交替为$E_A$与$E_B$；
  各\code{[CLS]}位置的顶层向量$\bm t_i$作为句子表示，经句间Transformer层建模句子之间的关系，
  再由逻辑回归给出每句的入选概率。图中只画三个句子，词元数受预训练位置嵌入的长度限制。}
  \label{fig:nlp-generation-bertsum}
\end{figure*}

句向量之上再叠加$L$层句间Transformer，让每句的表示看到其他句子，
相当于SummaRuNNer句级GRU的作用。记$\bm T=(\bm t_1,\ldots,\bm t_N)$，
$\bm H^{(0)}=\bm T+\operatorname{PosEmb}(\bm T)$，每层执行
$\tilde{\bm H}^{(\ell)}=\operatorname{LN}\bigl(\bm H^{(\ell-1)}+\operatorname{MHAtt}(\bm H^{(\ell-1)})\bigr)$、
$\bm H^{(\ell)}=\operatorname{LN}\bigl(\tilde{\bm H}^{(\ell)}+\operatorname{FFN}(\tilde{\bm H}^{(\ell)})\bigr)$，
最后$\hat z_i=\sigma(\bm w_o^{\top}\bm h^{(L)}_i+b_o)$。原文比较了若干层数，两层效果最好。
训练同样用oracle标签上的二元交叉熵。推断时选概率最高的若干句，并做\term{三元组阻断}
（trigram blocking）：若候选句与已选句共有一个三元组，就跳过它。
三元组阻断是一种简单的显式去冗余，作用与SummaRuNNer的冗余项相同，但不需要学习。

从SummaRuNNer到BertSum，变化在于句子表示的来源：前者从头训练循环编码器，
后者借预训练编码器的上下文表示，只训练少量句间层与输出层。两者都受oracle标签近似性的限制，
输出也都不能压缩句内信息。BertSum的输入长度受预训练位置嵌入限制，长文档需要截断，
这一问题在第~\ref{sec:nlp-generation-long-summary}小节处理。

\subsection{生成式摘要}
\label{sec:nlp-generation-abstractive}

\term{生成式摘要}（abstractive summarization）逐词写出摘要，可以删去句内冗余、合并多句信息，
也能用原文没有的词概括。它带来三个新问题：罕见的名称和数字不在词表中，
解码器会反复关注同一片段而重复输出，生成的句子可能与来源矛盾。
谱系依次是：注意力编码器—解码器首先在标题生成上展示生成式压缩的可行性；
pointer-generator用复制处理罕见词、用覆盖机制抑制重复；
去噪预训练让模型在大规模文本上先学会重建，再微调到摘要。
三者都以式~\eqref{eq:nlp-generation-nll}的负对数似然为主要训练目标，
差别在结构与初始化。

\subsubsection{注意力编码器—解码器摘要}
\label{sec:nlp-generation-rush-abs}

Rush等把句子级摘要定义为：给定新闻首句$x_{1:M}$，生成一个更短的标题$y_{1:N}$，
用约四百万对首句与标题训练\scite{rush2015neural}
解码器是一个前馈神经语言模型，只看已生成摘要的最后$C$个词$\bm y_c$。
设词嵌入矩阵为$\bm E\in\mathbb{R}^{D\times V}$，
\begin{equation}
  \begin{aligned}
    \tilde{\bm y}_c&=[\bm E y_{i-C+1};\ldots;\bm E y_i]\in\mathbb{R}^{CD},\qquad
    \bm h=\tanh(\bm U\tilde{\bm y}_c),\\
    p(y_{i+1}\mid\bm y_c,x)&\propto\exp\bigl(\bm V\bm h+\bm W\operatorname{enc}(x,\bm y_c)\bigr),
  \end{aligned}
  \label{eq:nlp-generation-rush-decoder}
\end{equation}
其中$\bm h\in\mathbb{R}^{H}$，$\bm V,\bm W\in\mathbb{R}^{V\times H}$。
编码器$\operatorname{enc}$是一个注意力模块：把源词嵌入为$\tilde{\bm x}_q=\bm F x_q\in\mathbb{R}^{H}$，
对相邻$\pm Q$个位置取平均得到平滑表示$\bar{\bm x}_q$，再以摘要上下文$\tilde{\bm y}'_c$
（用另一组嵌入）为查询计算注意力：
\begin{equation}
  \bm p=\operatorname{softmax}\bigl(\tilde{\bm X}^{\top}\bm P\tilde{\bm y}'_c\bigr)\in\mathbb{R}^{M},\qquad
  \operatorname{enc}(x,\bm y_c)=\sum_{q=1}^{M}p_q\,\bar{\bm x}_q,
  \label{eq:nlp-generation-rush-attention}
\end{equation}
$\tilde{\bm X}\in\mathbb{R}^{H\times M}$的各列为$\tilde{\bm x}_q$，$\bm P$为双线性参数。
解码器每生成一个词，注意力就重新分配到源句的不同位置，作用与
式~\eqref{eq:nlp-generation-bahdanau}的注意力相同；区别在于打分是双线性形式，
编码器也没有循环结构，源词表示只经过局部平均。推断用束搜索。

这一模型表明，不依赖句法规则、端到端学习的生成式压缩是可行的。
它也暴露了复制问题：名称、数字等罕见词在词表中缺失或训练不足，模型常把它们换成常见词。
原文的补救是ABS+：在对数线性框架中加入源句与摘要的一至三元词匹配特征，
用最小错误率训练调节权重，鼓励模型照抄源句中的词。
随后的工作把前馈解码器换成循环解码器，Nallapati等还加入在生成与指向源位置之间切换的开关，
并把任务从单句标题扩展到多句新闻摘要\scite{nallapati2016abstractive}
这一开关是下面pointer-generator的前身。

\subsubsection{pointer-generator复制与覆盖}
\label{sec:nlp-generation-pointer}

See等的pointer-generator在CNN/DailyMail多句新闻摘要上同时处理复制与重复两个问题
\scite{nlp-intro-see2017}
复制部分就是第~\ref{sec:nlp-generation-copy-pointer}小节的指针生成网络：
生成门按式~\eqref{eq:nlp-generation-pointer-generator}混合词表分布与注意力给出的复制分布，
扩展词表让源端未登录词也有概率；门的输入、扩展词表的构造及数值算例都在那里给出。
摘要中的具体配置是：编码器为单层双向LSTM，$\bm h_i\in\mathbb{R}^{512}$（每向256维），
解码器为单层LSTM，源端与目标端共用最常见的五万词。复制使固定词表可以保持较小，
而文章中的人名、地名和数字仍能原样写出。摘要特有的问题是重复，下面的覆盖机制处理它。

\paragraph{覆盖机制}
多句摘要中，解码器常在不同步骤反复关注同一片段，输出重复的句子。
\term{覆盖机制}（coverage mechanism）维护此前各步注意力的累加，即覆盖向量
\begin{equation}
  \bm c_t=\sum_{t'<t}\bm a_{t'},\qquad c_{t,i}=\sum_{t'<t}a_{t',i},
  \label{eq:nlp-generation-coverage}
\end{equation}
$c_{t,i}$表示源位置$i$到第$t$步为止已经被关注的总量，$\bm c_1=\bm 0$。
覆盖向量以两种方式起作用。其一，它作为注意力打分的额外输入，
把式~\eqref{eq:nlp-generation-bahdanau}改为
\begin{equation}
  e_{t,i}=\bm v^{\top}\tanh\bigl(\bm W_h\bm h_i+\bm W_s\bm s_t+\bm w_c\,c_{t,i}+\bm b\bigr),
  \label{eq:nlp-generation-coverage-attention}
\end{equation}
让模型在决定当前注意哪里时知道哪里已经读过。其二，\term{覆盖损失}（coverage loss）
惩罚当前注意力与已有覆盖的重叠：
\begin{equation}
  L_t=-\log p_t(y^{\ast}_t)+\lambda\sum_{i}\min(a_{t,i},c_{t,i}).
  \label{eq:nlp-generation-coverage-loss}
\end{equation}
第一项是第$t$步的负对数似然，第二项在当前注意力$a_{t,i}$与覆盖$c_{t,i}$都大的位置才大，
即只惩罚“再次关注已经关注过的位置”。用$\min$而不是乘积或$a_{t,i}$本身，
是为了让惩罚有界：每步覆盖损失不超过$\sum_i a_{t,i}=1$。
它不要求每个源位置都被关注到，只要求不重复关注，因为摘要本来就应跳过大量背景。
总损失为各步$L_t$的平均，原文取$\lambda=1$。

用数值核对这两项。设教学来源词元为“读书会／地点／海棠楼201／请／到／海棠楼201”，
此前各步注意力累加得覆盖$\bm c_t=(0,0,0.6,0.1,0,0.3)$，
本步注意力为$\bm a_t=(0.05,0.05,0.45,0.10,0.05,0.30)$，$p_{\mathrm{gen},t}=0.4$。
若“海棠楼201”是未登录词，它的输出概率为$0.6\times(0.45+0.30)=0.45$，
负对数似然为$-\log0.45\approx0.799$。
覆盖损失为$\min(0.05,0)+\min(0.05,0)+\min(0.45,0.6)+\min(0.10,0.1)+\min(0.05,0)+\min(0.30,0.3)=0.85$，
本步总损失为$0.799+0.85\approx1.649$。这一步正在重复关注两处“海棠楼201”，覆盖损失相应很大；
若注意力转到尚未读过的“读书会”，重叠会降到接近零。

覆盖损失要在模型学会对齐之后才加入。原文先不带覆盖训练至收敛，再加入覆盖机制和覆盖损失继续训练约三千步；
从一开始就加入覆盖目标会干扰主目标、降低整体效果，
而只加覆盖向量、不加覆盖损失时，看不出重复有所减少。
覆盖机制只防止“重复关注”，不判断信息是否该写：它可以避免重复，
却不能阻止模型复制一个正确名称后把它接到错误的活动上。

\subsubsection{去噪预训练摘要：BART、T5与PEGASUS}
\label{sec:nlp-generation-conditional-summary}

pointer-generator的全部参数都从摘要语料学起，几十万篇文档难以教会模型丰富的语言知识。
预训练的编码器—解码器先在大规模无标注文本上学习“从损坏的输入重建文本”，
再以“来源→摘要”的平行数据微调。摘要所需的能力，即读懂长输入、复制关键片段、
改写和压缩，都在重建任务中得到练习。第~\ref{sec:nlp-llm-objective-t5-bart}小节已给出
BART与T5两种去噪目标的构造，这里说明它们与摘要的关系，并补充专为摘要设计的PEGASUS。

\paragraph{BART与T5}
BART比较了五种噪声：词元掩码、词元删除、文本填充（按泊松分布抽取跨度长度，
每个跨度换成一个掩码符号）、句子置换和文档旋转，解码器总是重建完整原文
\scite{lewis-etal-2020-bart}
重建全文意味着解码器大量练习照抄未损坏的部分、补全缺失的部分，
这与摘要“多数内容来自来源、少数需要改写”的形态接近。最终模型组合了文本填充与句子置换，
后者要求模型理解句间顺序，对需要重组信息的摘要尤其有关。
T5用跨度损坏，只生成被删的跨度，并把所有任务写成带前缀的文本到文本格式，
摘要输入写作“summarize: 原文”\scite{raffel2020t5}
两者微调到摘要时都只需把来源作为编码器输入、摘要作为解码器目标，
复制不再需要显式的指针结构：预训练已经让交叉注意力与输出层习惯于照抄输入中的词元，
子词切分也使罕见名称可以按片段拼出。

\paragraph{PEGASUS}
BART和T5的预训练目标是通用的。Zhang等认为，预训练目标越接近下游任务，
微调所需的数据越少，于是设计了\term{间隔句生成}（gap sentences generation, GSG）：
从文档中选出若干完整句子，在输入中各替换为一个掩码符号，目标是把这些句子按原顺序拼接生成出来
\scite{zhang2020pegasus}
被挖去的句子由剩余文档生成，这正像从文档生成一个“伪摘要”。

选哪些句子是关键。随机选句和取开头几句都不如选“重要”的句子：
PEGASUS对每个句子$s_i$计算它与文档其余部分的ROUGE-1 F值
$\operatorname{R}_1(s_i,D\setminus\{s_i\})$，按这一分数独立排序，选前$m$句，
$m$按句数的固定比例确定，原文的消融实验以约30\%为宜。
与其余部分重合越多的句子越能概括全文，被挖去后也越能从剩余文本推断。
用四句教学文档手算：$s_1$“读书会／改至／周五／下午／举行”，
$s_2$“读书会／地点／海棠楼201”，$s_3$“周五／下午／读书会／在／海棠楼201／举行”，
$s_4$“欢迎／参加”。按词计数、每个词的匹配次数不超过它在两侧出现的较小次数，
$s_3$与其余三句共10个词匹配5个，
精确率$5/6$、召回率$5/10$，F值0.625；$s_1$对其余11个词匹配4个，F值0.5；
$s_2$为0.25，$s_4$为0。取$m=1$，挖去$s_3$，
输入为“$s_1$ \code{[MASK1]} $s_2$ $s_4$”，目标为$s_3$。
$s_3$同时包含时间与地点，正是这份文档最像摘要的一句。

\paragraph{微调目标的构造}
预训练决定模型会什么，微调目标决定它学会保留什么。会议纪要若只以“会议改至周五”
这类结论句为目标，模型就学不到“建议”与“决定”的区别；目标应把二者分开写，
并把行动项写成“负责人、动作、截止时间、决定依据”的固定格式。例如教学会议决定
“由林岚在4月16日18:00前汇总报名信息”，行动项就是（林岚，汇总报名信息，4月16日18:00，会议决定）；
原通知没有这位负责人，摘要不能把会议补充的事实写成通知已说过。
目标格式越明确，生成后越容易逐项回查来源。

\paragraph{比较}
三者共享编码器—解码器结构，区别在预训练目标与摘要的距离：
T5的跨度损坏目标短、计算省，最为通用；BART重建全文，更接近摘要中的复制与改写；
PEGASUS直接以伪摘要为目标，在微调数据很少时优势最明显：原文只用一千个微调样本，
就在十二个摘要数据集中的六个上超过了此前的最好结果。代价是PEGASUS的预训练专为摘要设计，
原文也只在摘要任务上评测；三者的输入长度都受位置编码与注意力计算的限制，
常见设置为数百到一千余词元。微调后的模型仍可能产生与来源矛盾的陈述，
预训练提高了流畅度，并不直接约束忠实性，忠实性核验见第~\ref{sec:nlp-generation-faithfulness}小节。

\subsection{长文档摘要}
\label{sec:nlp-generation-long-summary}

论文、报告和书籍章节动辄数万词元，超过上述模型的长度上限。
完整自注意力每层计算$n\times n$的得分矩阵，计算和显存都随长度平方增长；
简单截断又会丢掉结尾的结论。长文档摘要的路线可按“在哪里丢信息”区分。

\paragraph{截断与先抽取后生成}
最简单的是只读前$n_{\max}$个词元。新闻的开头偏重使这一做法在新闻上损失不大，
在论文上则会丢掉方法和结论。先抽取后生成用一个廉价的抽取器（如TextRank或有监督句子打分）
从全文选出若干句，拼成不超过长度上限的输入，再交给生成模型
\scite{liu2018generating}
信息在抽取阶段丢失，生成器无法恢复被抽取器误删的条件。

\paragraph{分段摘要与合并}
分治法把文档切成若干段，分别摘要再合并或再摘要一次。每段都在长度上限内，
但段与段之间的关系在第一阶段就被切断：会议材料若先写“林岚建议改至周五”，
后面另一段才写“与会者否决该建议”，第一阶段若只保存“会议改至周五”，
第二阶段已经无从恢复建议状态。为此，第一阶段的输出应保存“陈述、原文范围、
是否含条件或待决状态”，合并阶段可以回查原文。

\paragraph{分层编码}
分层编码让模型看到全文的结构。Cohan等为论文设计了两级编码器：
词级双向LSTM在每个章节内编码，章节级LSTM在章节表示的序列上编码；
解码时，注意力先分配到章节，再分配到章节内的词，词$i$（属于章节$j$）的权重为
$\alpha_{t,(j,i)}\propto\beta_{t,j}\exp\bigl(\operatorname{score}(\bm h_{(j,i)},\bm s_t)\bigr)$，
$\beta_{t,j}$是对章节的注意力\scite{cohan2018discourse}
这样的分层注意力利用了“摘要通常按章节组织”的篇章结构，计算量也只与章节数和章节长度有关。

\paragraph{稀疏注意力：LED}
另一条路线保留全文作为输入，但让每个位置只读取一部分位置。
Longformer把完整自注意力换成滑动窗口局部注意力加少量全局注意力：
每个位置只读取左右各$w/2$个邻居，指定的少数位置则读取所有位置、也被所有位置读取，
每层代价从$O(n^2)$降为$O(nw)$\scite{nlp-llm-longformer2020}
其机制见第~\ref{sec:nlp-llm-sliding-window}小节。
Longformer-Encoder-Decoder（LED）把这一编码器接到BART上：
编码器自注意力换成窗口加全局注意力，解码器保持完整的因果自注意力和对编码器的完整交叉注意力；
BART原有一千零二十四个位置嵌入，LED把它们复制十六次，把位置扩展到16384，
免去从头学习长位置。摘要时通常只给首个词元设全局注意力，使它汇总全文。
取$n=16384$、$w=1024$，每层得分计算从$n^2\approx2.7\times10^8$降到$nw\approx1.7\times10^7$，
约为十六分之一；解码器的交叉注意力仍是$m\times n$，因输出长度$m$远小于$n$而可以承受。

\paragraph{比较}
截断与先抽取后生成在输入端丢信息，丢失位置固定或取决于抽取器；
分段摘要在段边界丢失跨段关系；分层编码保留全文但把细节压到章节表示中；
稀疏注意力保留全文的每个词元，远距离信息只能经全局位置或多层窗口接力传递。
从前到后，保留下来的原文信息依次增多，计算量也随之上升。
即使输入能放进上下文，模型也不一定用好全部内容，长上下文中段信息利用不足的问题
见第~\ref{sec:nlp-inference-lost-middle}小节。

\subsection{基于人类反馈的摘要}
\label{sec:nlp-generation-feedback-summary}

前面的方法都以模仿参考摘要为目标，用ROUGE评价。两者都与人们对摘要质量的判断不完全一致：
参考摘要本身质量参差，一份摘要可以与参考措辞不同而同样好；
ROUGE奖励词面重合，对事实错误不敏感。Stiennon等把训练目标换成人类偏好，
用第~\ref{chap:nlp-preference-alignment}章的RLHF流程训练摘要模型
\scite{nlp-align-stiennon2020}

流程分三步。第一步，在经过筛选的Reddit帖子与作者自写摘要（TL;DR）上监督微调，
得到策略$\pi^{\mathrm{SFT}}$。第二步，对同一帖子从若干策略采样两份摘要，请标注者选出更好的一份，
得到比较数据$(x,y_w,y_l)$；用第~\ref{sec:nlp-align-bradley-terry}小节的
Bradley–Terry模型训练奖励模型$r_{\bm\phi}$，损失为
\begin{equation}
  L_{\mathrm{RM}}(\bm\phi)=-\mathbb{E}_{(x,y_w,y_l)}
  \Bigl[\log\sigma\bigl(r_{\bm\phi}(x,y_w)-r_{\bm\phi}(x,y_l)\bigr)\Bigr].
  \label{eq:nlp-generation-summary-rm}
\end{equation}
第三步，用PPO最大化奖励减去与监督模型的KL惩罚：
\begin{equation}
  R(x,y)=r_{\bm\phi}(x,y)-\beta\log\frac{\pi^{\mathrm{RL}}(y\mid x)}{\pi^{\mathrm{SFT}}(y\mid x)}.
  \label{eq:nlp-generation-summary-ppo}
\end{equation}
KL项把策略拴在监督模型附近，既防止它生成奖励模型没见过的文本，也保留一定的输出多样性。
PPO的裁剪目标与价值函数见第~\ref{sec:nlp-align-ppo}小节，
策略梯度与RLHF的一般推导可参考\BookChapterRef{chap:rl-reward}{卷七的奖励学习章}。
用数值看一次更新的信号：设奖励模型给两份摘要打分$1.2$与$0.4$，
则它判断前者更好的概率为$\sigma(0.8)\approx0.69$，这一对的损失为$-\log0.69\approx0.37$；
策略生成的某份摘要奖励为$1.0$，其对数概率比监督模型高$2$，取$\beta=0.05$，
扣除KL惩罚后$R=1.0-0.1=0.9$。

摘要任务中有三点现象。其一，以人类偏好训练的13亿参数模型，在人工比较中优于十倍规模的监督模型，
也常优于数据集中的参考摘要。用best-of-N拒绝采样比较优化对象时，
以ROUGE为选择标准的摘要质量比以奖励模型为标准时更早见顶，峰值也明显更低。
其二，人们倾向于偏好较长的摘要，奖励模型会学到这一偏好；原文控制长度后再比较，
67亿参数模型的摘要仍在约65\%的比较中胜过参考摘要，但比较摘要质量时必须交代长度。
其三，\term{过度优化}（overoptimization）：沿奖励模型持续优化、与监督模型的KL距离越拉越大时，
奖励模型分数继续上升，真实的人类偏好却在某一点之后下降，策略开始利用奖励模型的漏洞。
奖励模型过度优化的一般讨论见第~\ref{sec:nlp-align-overoptimization}小节。
这一方法的代价是需要持续收集人工比较，且奖励模型只覆盖标注者关注的维度；
一份摘要可能读起来令人满意，却漏掉了不显眼但关键的条件。

\subsection{多来源信息综合}
\label{sec:nlp-generation-multi-summary}

\term{多文档摘要}（multi-document summarization）的输入是讨论同一主题的多份材料。
来源增多后，重复不能当作多份独立证据，表面不一致也不一定是矛盾：
旧通知的4月17日和改期通知的4月18日对应不同版本；另一活动也在4月18日举行，则属于不同对象。
因此要先对齐“对象、属性、值、适用时点、来源”，再判断是转载、互补、条件差异还是冲突。
确认版本关系后可以表达最新状态；两个同样有效而相互冲突的来源应并列表述或请求澄清，
不能按转载数量投票。图~\ref{fig:nlp-generation-provenance}展示同一活动的三份教学通知中
的重复、互补和版本差异。

\begin{figure*}[htbp]
  \centering
  % Adapted from academic-drawing/templates/tikz_template.tex to the book palette.
% Source fragments are grouped by claim, not by chronological reading order.
\begin{tikzpicture}[
  every node/.style={font=\figurefont\small,align=center,line width=.6pt},
  source/.style={book node,text width=59mm,minimum height=12mm,
    fill=FigureInput!10,draw=FigureInput,font=\figurefont\small},
  claim/.style={book node,text width=65mm,minimum height=14mm,
    fill=FigureOutput!10,draw=FigureOutput,font=\figurefont\small},
  support/.style={book arrow,rounded corners=1mm},
  qualifier/.style={support,dashed,draw=BookAmber}
]
  \node[font=\figurefont\bfseries] at (-47mm,10mm) {来源片段};
  \node[font=\figurefont\bfseries] at (45mm,10mm) {摘要陈述及保留状态};
  \node[source] (date1) at (-47mm,0)
    {$D_1$：4月15日官方改期通知\\4月18日14:00};
  \node[source] (date2) at (-47mm,-15mm)
    {$D_2$：同日转载\\4月18日14:00};
  \node[claim,minimum height=23mm] (date) at (45mm,-7.5mm)
    {当前时间：4月18日14:00\\$D_1,D_2$支持同一陈述\\转载去重，不增加独立证据};
  \node[source] (place) at (-47mm,-32mm)
    {$D_2$：地点片段\\海棠楼201};
  \node[claim] (placeout) at (45mm,-32mm)
    {地点：海棠楼201\\与日期互补};
  \node[source,minimum height=19mm] (rule) at (-47mm,-51mm)
    {$D_1$：校外参加条件\\须在4月16日18:00前\\提交报名信息};
  \node[claim,minimum height=19mm] (ruleout) at (45mm,-51mm)
    {校外来宾须在4月16日18:00前\\提交报名信息\\保留角色、义务与截止条件};
  \node[source,draw=BookAmber,fill=BookAmber!8] (old) at (-47mm,-70mm)
    {$D_3$：4月14日旧通知，现已归档\\原定4月17日};
  \node[claim,draw=BookAmber,fill=BookAmber!8] (oldout) at (45mm,-70mm)
    {旧版日期：4月17日\\不作为当前安排};
  \node[source] (inside) at (-47mm,-86mm)
    {$D_3$：校内参加规则\\校内师生可直接参加};
  \node[claim,draw=BookMuted,fill=BookPaper] (omission) at (45mm,-86mm)
    {本次校外指南可省略\\仍保留原文，供其他用途调用};
  \node[draw=FigureLoss,fill=FigureLoss!7,rounded corners=1mm,
    text width=150mm,minimum height=11mm] at (-1mm,-103mm)
    {无支持新增：“主讲人为某教授”。三份资料均未给出主讲人，不能写入事实摘要。};
  \draw[support] (date1.east) -- (-3mm,0) |- ([yshift=4mm]date.west);
  \draw[support] (date2.east) -- (-3mm,-15mm) |- ([yshift=-4mm]date.west);
  \draw[support] (place.east) -- (placeout.west);
  \draw[support] (rule.east) -- (ruleout.west);
  \draw[qualifier] (old.east) -- node[above,font=\figurefont\footnotesize] {旧版} (oldout.west);
  \draw[support,draw=BookMuted] (inside.east) -- (omission.west);
\end{tikzpicture}

  \caption{摘要陈述与来源的对应。三份教学通知按活动编号\code{ML-0418}对齐；
  实线表示支持，虚线表示需要版本解释的日期差异。转载的相同日期合为一个事实，
  地点与报名条件分别保留来源。灰框标明用途允许省略的校内规则，
  底部警示框标明无来源支持的新增陈述，二者不能混为一类。}
  \label{fig:nlp-generation-provenance}
\end{figure*}
\FloatBarrier

图中$D_1$为4月15日官方改期通知，含当前时间和校外规则；
$D_2$为同日转载，重复当前时间并补列地点；
$D_3$为4月14日发布、现已归档的旧通知，给出原定4月17日及校内规则。
它们都是本节构造的来源片段。支持与完整需要分开检查：
摘要中每条陈述都有来源支持，即支持率高，也不说明该写的信息都已写出，即完整率高。
方法谱系有三层：选择式方法从各来源选片段，来源可追溯；
词图融合把近重复句压成一条更短的句子；分层编码与多文档预训练直接生成融合文本，
表达能力最强，却需要事后回查来源。

\subsubsection{MMR与最大覆盖的选择式综合}
\label{sec:nlp-generation-mmr}

多份材料的高相关片段往往彼此重复，只按相关性排序，同一条日期就会占满摘要。
\term{最大边际相关性}（maximal marginal relevance, MMR）每次选择时既考虑片段对查询的相关性，
也扣除它与已选内容的最大相似度\scite{nlp-generation-carbonell1998}
对已选集合$S$和尚未选入的片段$i$，
\begin{equation}
  i^*=\argmax_{i\notin S,\;\operatorname{cost}(S\cup\{i\})\leq B}
  \left[\lambda\,\operatorname{rel}(i,q)
  -(1-\lambda)\max_{j\in S}\operatorname{sim}(i,j)\right].
  \label{eq:nlp-generation-mmr}
\end{equation}
其中$q$为查询，$B$为预算，$\lambda\in[0,1]$权衡相关性与去冗余；$S$为空时最大相似度定义为零。
每轮取可容纳且分数为正的最高项，直到预算不足、没有正分项或候选耗尽。

在三份资料中，先把$D_1,D_2$逐字重复的当前时间聚成一个片段，保留来源集合$\{D_1,D_2\}$；
确认$D_3$为旧版本后，将旧日期留在历史记录中，不列为当前安排候选。
教学候选与固定分词成本如下，片段乙是$D_2$中的一句同义重述。

\begin{fullwidth}
\begin{center}
  \begin{tabular}{@{}clrr@{}}
    \toprule
    片段 & 按斜线分词的内容 & 成本 & 相关性\\
    \midrule
    甲 & 读书会／改至／4月18日／14:00 & 4 & 0.95\\
    乙 & 读书会／时间／已更新／至／4月18日 & 5 & 0.92\\
    丙 & 地点／为／海棠楼201 & 3 & 0.80\\
    丁 & 校外来宾／须／在／4月16日18:00／前／提交／报名信息 & 7 & 0.85\\
    \bottomrule
  \end{tabular}
\end{center}
\end{fullwidth}
片段间的相似度为：甲乙0.90，甲丙0.10，甲丁0.20，乙丙0.10，乙丁0.15，丙丁0.10，均为教学输入。
甲来自$D_1,D_2$，乙、丙来自$D_2$，丁来自$D_1$。

取$\lambda=0.7$、$B=14$。首轮分数为甲0.665、乙0.644、丙0.560、丁0.595，选甲。
第二轮乙减去$0.3\times0.9$后为0.374，丙为0.530，丁为0.535，选丁，已用成本11。
第三轮只有丙仍可容纳，得分0.530，选入后成本恰为14。
若只按相关性从高到低扫描可容纳项，会选甲、乙，剩余5容不下丁，再选丙，遗漏校外条件。

\term{最大覆盖}（maximum coverage）换一个角度计分：不惩罚片段相似，而是奖励新覆盖的事实。
设日期、时刻、地点、完整校外规则四个事实的权重分别为$2,1,2,4$；
甲覆盖前两项，乙只覆盖日期，丙、丁分别覆盖后两项。若$c_i$是片段成本，目标是
\begin{equation}
  \max_{S:\;\sum_{i\in S}c_i\leq B}\;\sum_{f}w_f\,\mathbb{I}\bigl[f\text{被}S\text{中某个片段覆盖}\bigr].
  \label{eq:nlp-generation-max-coverage}
\end{equation}
已覆盖的事实再被覆盖不加分，冗余因此自动不受奖励。
以“新增权重／成本”贪心选择，初值为甲$3/4$、乙$2/5$、丙$2/3$、丁$4/7$，先选甲；
此后乙的新增权重为0，再依次选丙、丁，覆盖权重9，成本14；相关性基线的甲、乙、丙只覆盖权重5。
带预算的最大覆盖是NP困难问题，按性价比贪心的结果需要再与“单选价值最高的一项”比较取优，
才有常数近似比保证；小例子上的贪心结果不能推广为一般最优。

选取次序甲、丁、丙不是阅读次序。把甲、丙组合为“何时何地”，再表述丁的参加前提，
按陈述保存来源：甲来自$\{D_1,D_2\}$，丙来自$D_2$，丁来自$D_1$。
若$D_3$不是归档旧版而是同日有效的通知，冲突尚未解决，就不能在聚合时悄悄删除它，
应保留冲突标记，并在最终表达中说明不确定性。

\subsubsection{词图最短路径压缩与多句融合}
\label{sec:nlp-generation-word-graph}

选择式综合保留原句的冗长部分。对已经确认表达同一事件的近重复句簇，
Filippova把各句的词及相邻关系构成\term{词图}（word graph），
在图上寻找一条较短且受来源支持的路径，实现\term{多句压缩}（multi-sentence compression）
\scite{filippova-2010-multi}
词形和词性相同的词优先共用节点，有多个可能位置时依据相邻上下文区分；
起止节点连接每个句子的首尾。图可能有环。一条边在句簇中出现$f(u,v)$次，
基础版本的成本为$1/f(u,v)$，出现越多的相邻关系越便宜；
枚举若干条短路径，按是否含动词、是否达到最小长度过滤，再选成本最低者。

输入三句，分词固定为
\begin{quote}
读书会／于／周五／在／海棠楼201／举行。\\
本次／读书会／周五／在／海棠楼201／举行。\\
读书会／将／于／周五／在／海棠楼201／举行。
\end{quote}
“起点到读书会”出现2次，成本$1/2$；“于到周五”出现2次，成本$1/2$；
共同尾部“周五到在、在到海棠楼201、海棠楼201到举行、举行到终点”各出现3次，各成本$1/3$；
其他边各出现1次，成本1。
路径“起点、读书会、周五、在、海棠楼201、举行、终点”的成本为$1/2+1+4/3=17/6$，
恢复为“读书会周五在海棠楼201举行”。它使用第一、三句的起点边和第二句的“读书会到周五”边，
并非照抄某一原句。经“于”的路径成本为$1/2+1+1/2+4/3=10/3$，经“本次”的路径为$13/3$。
最短候选有五个词元并含动词“举行”，通过最小长度为五的教学过滤条件。

词图方法不需训练，输出的每条边都来自某个来源句。但连通性、低成本和含动词只检查了图与表面形式，
不保证关系正确：若把不同角色的句子错误聚为一簇，合法路径也可能生成“校内师生须提前提交”。
句簇的划分因此必须先于压缩完成，融合互补事实也超出了它的能力。

\subsubsection{分层编码与多文档预训练}
\label{sec:nlp-generation-multi-pretraining}

选择与词图都只在原句层面操作，无法把不同来源的互补信息写进一句话。
神经生成模型可以融合，但多文档输入更长，文档之间的关系也需要建模。

\paragraph{分层Transformer}
Liu与Lapata针对由多个网页生成维基百科首段的任务，先用一个排序器给全部段落打分，
保留排名靠前的若干段，再用分层Transformer编码\scite{liu2019hierarchical}
每层包含两部分：局部层在每个段落内部做普通的自注意力；
全局层先用多头池化把每个段落压成每头一个向量，即对段内各词元
$a^{z}_{k,j}=\operatorname{softmax}_j(\bm w^{z\top}_a\bm x_{k,j})$、
$\bm u^{z}_k=\sum_j a^{z}_{k,j}\bm W^{z}_b\bm x_{k,j}$，
其中$\bm x_{k,j}$是第$k$段第$j$个词元的表示，$z$为头的编号；
再在段落向量之间做多头注意力，得到每段的跨段上下文，并加回该段的每个词元。
这样，词元级计算只在段内进行，跨段交互发生在段落向量这一级，
计算量随段数和段长增长，而不是随总长度平方增长。
它的思路与第~\ref{sec:nlp-generation-long-summary}小节的分层编码相同，区别在于这里的“段”来自不同文档。

\paragraph{PRIMERA}
分层Transformer需要在每个多文档数据集上从头学习跨文档关系。
PRIMERA借鉴PEGASUS，设计专为多文档摘要的预训练\scite{xiao2022primera}
结构上，它以LED为基础，把同一簇的各文档拼接，文档之间插入分隔符\code{<doc-sep>}，
并给这些分隔符设全局注意力，使它们成为文档间交换信息的枢纽。
预训练目标仍是挖去句子再生成，选句策略称为\term{实体金字塔}（entity pyramid）：
抽取簇内所有实体，按出现在多少份文档中排序，频繁跨文档出现的实体代表簇的共同主题；
对排名靠前的每个实体，在含有它的句子中选与簇内其他文档ROUGE最高的一句，挖去后作为生成目标。
与PEGASUS按“与本文档其余部分的重合”选句相比，实体金字塔按“与其他文档的重合”选句，
迫使模型从其他文档推断被挖去的内容，从而练习跨文档的信息汇聚。

\paragraph{比较与来源回查}
选择式方法每条陈述都对应原文片段，来源可追溯，但不能融合；
词图压缩能删去冗余，却只限于同一事件的近重复句；
分层编码与多文档预训练能把互补信息写进一句，代价是输出不再与来源一一对应，
跨文档的对象绑定也可能出错。例如另一条教学记录说明，编号\code{IR-0418}的
“检索系统研讨会”也在4月18日14:00举行，地点为松柏楼305；
若模型只按日期把两份材料对齐，就可能生成“读书会在松柏楼305举行”这一语法正确的错误陈述。
引用一份确实含有“松柏楼305”的资料并不能支持这一陈述，支持对象是完整的关系，而不是孤立的词语。
因此，融合式生成之后需要把每条陈述拆回“对象、属性、值、来源”，逐项回查，
具体核验见第~\ref{sec:nlp-generation-faithfulness}小节和第~\ref{chap:nlp-question-answering-dialogue}章。

\paragraph{摘要方法的对照}
表~\ref{tab:nlp-generation-summarization-methods}汇总本节方法。
从上到下，输出越来越依赖生成而不是选择，压缩与融合能力上升，
对来源的忠实则越来越需要额外的核验来保证；可处理的输入长度与训练数据需求，
则分别由结构与预训练决定。

\begin{table*}[htbp]
  \centering
  \small
  \begin{tabularx}{\linewidth}{@{}>{\raggedright\arraybackslash}p{26mm}
    >{\raggedright\arraybackslash}p{28mm}>{\raggedright\arraybackslash}p{38mm}
    >{\raggedright\arraybackslash}p{26mm}>{\raggedright\arraybackslash}X@{}}
    \toprule
    方法 & 解决的问题 & 核心机制 & 输入长度与数据 & 忠实性与新代价\\
    \midrule
    Lead、TextRank & 无训练数据 & 位置先验；句子图上的中心性传播 & 任意长度；无需标注 & 原句忠实；不能压缩，易漏独特条件\\
    SummaRuNNer、BertSum & 规则打分不可学习 & oracle标签上的句子二分类；冗余项或三元组阻断 & BertSum受位置嵌入限制；需摘要语料 & 原句忠实；oracle标签近似\\
    注意力编码器—解码器 & 抽取不能压缩 & 每步注意力读取源句，逐词生成 & 单句级；大规模平行对 & 罕见词复制困难\\
    pointer-generator & 罕见词与重复 & 生成门混合复制分布；覆盖向量与覆盖损失 & 数百词元；摘要语料 & 复制正确仍可能绑定错误\\
    BART、T5、PEGASUS & 摘要语料不足以学语言知识 & 去噪或间隔句预训练后微调 & 约一千词元；少量微调数据即可 & 流畅不等于忠实\\
    长文档方法 & 超出长度上限 & 截断、抽取后生成、分段、分层、稀疏注意力 & 至数万词元 & 丢信息的位置各不相同\\
    人类反馈 & 模仿参考与ROUGE偏离人类判断 & 偏好比较训练奖励模型，PPO加KL & 需持续人工比较 & 长度偏好，奖励过度优化\\
    MMR、最大覆盖 & 多来源重复 & 相关性减冗余；覆盖事实的权重 & 片段级；无需训练 & 来源可追溯；不能融合\\
    词图压缩 & 近重复句冗长 & 词图上的最短路径 & 同一事件的句簇 & 可能接错主体\\
    分层编码、PRIMERA & 跨文档融合 & 段内与段间注意力；实体金字塔预训练 & 多文档拼接，数千词元 & 需回查来源与对象绑定\\
    \bottomrule
  \end{tabularx}
  \caption{摘要与信息综合方法对照。“输入长度”只给典型量级，具体取决于模型配置；
  所有生成式方法的输出都需要按第~\ref{sec:nlp-generation-faithfulness}小节的方法核验。}
  \label{tab:nlp-generation-summarization-methods}
\end{table*}
\FloatBarrier

\section{机器翻译}
\label{sec:nlp-generation-translation}

\term{机器翻译}（machine translation, MT）把源语言句子$x=(x_1,\ldots,x_S)$变成目标语言句子$y=(y_1,\ldots,y_T)$，要求传达同样的意思，允许语序调整、词数变化和句子拆合。共同通知的第二句“校外来宾须在4月16日18:00前提交报名信息”译成英文时，“须”要变成must，“18:00前”要变成before 18:00并移到句末，“校外来宾”是两个汉语词对应一个名词短语external guests。这一句已经包含了翻译的大部分困难。把整个研究问题展开，可以列出七类挑战：
\begin{enumerate}
  \item \emph{词序与对齐}：源词与目标词之间是一对多、多对一或无对应的关系，平行句子又不标明谁对应谁；中文状语前置、英文后置，日语动词在句末，局部和长距离的重排都需要建模。
  \item \emph{罕见词与开放词表}：人名、机构名、术语和形态变化使词表无法封闭，按词建模时总有训练中没见过的词。
  \item \emph{定长瓶颈}：神经模型若把整句源文压成一个固定维数的向量，长句中的信息难以完整保留。
  \item \emph{平行数据稀缺}：只有少数语向拥有上千万句对，多数语向只有几万句或没有平行数据，单语文本却多得多。
  \item \emph{多语言}：几百个语向各训练一个模型不现实，共享一个模型又带来容量竞争和输出语言控制问题。
  \item \emph{实时性}：同声传译在源句读完之前就要输出，后到的否定词或动词可能推翻已经说出的译文。
  \item \emph{篇章一致性}：逐句翻译看不到前文，跨句的指代、省略和术语统一会出错。
\end{enumerate}

各类方法大致按这些挑战依次出现。统计机器翻译用隐变量模型估计词对齐，在对齐上抽取短语、组合多种特征并搜索译文，处理了词序与对齐，但组件多、各自优化；神经机器翻译用一个条件模型端到端替代这些组件，并逐步去掉定长瓶颈与循环依赖、处理开放词表；回译、多语模型和无监督翻译依次放宽对平行数据的要求；大语言模型把翻译变成按指令生成，再借助质量估计和最小贝叶斯风险解码从候选中选择译文；同声传译与篇章翻译改变了模型能看到的条件，前者只见源前缀，后者要读当前句之外的上下文。表~\ref{tab:nlp-generation-translation-methods}按这条谱系列出各方法解决的问题与引入的新代价，后文逐项展开。

\begin{table*}[htbp]
  \small
  \begin{tabularx}{\linewidth}{@{}>{\raggedright\arraybackslash}p{32mm}>{\raggedright\arraybackslash}p{31mm}>{\raggedright\arraybackslash}p{44mm}>{\raggedright\arraybackslash}X@{}}
    \toprule
    方法 & 解决的问题 & 核心机制 & 新代价\\
    \midrule
    IBM词对齐 & 平行句子不标词对应 & 对齐作隐变量，EM估计词翻译表 & 不建模语序，只能用作对齐\\
    短语统计翻译 & 逐词替换缺少局部语序，特征难以组合 & 对称化对齐抽取短语，对数线性打分，栈解码，MERT调权 & 组件多且分别优化；长距离重排弱\\
    循环编码器—解码器 & 特征工程与组件拼接 & 端到端条件语言模型 & 源句压成定长向量\\
    加性注意力 & 定长瓶颈 & 每步对全部源状态加权求和 & 循环串行，训练难以并行\\
    Transformer & 循环依赖 & 自注意力、交叉注意力，训练时目标位置并行 & 注意力代价随长度平方增长\\
    子词与标签平滑 & 开放词表；过度自信 & 联合BPE与共享词表；软化的目标分布 & 切分可能违背词义；困惑度变差\\
    长度与覆盖惩罚 & 束搜索偏短、漏译 & 长度项除对数概率，覆盖项惩罚未被注意的源词 & 两个超参数；依赖注意力权重\\
    回译 & 平行数据少 & 反向模型把目标语单语译成合成源句 & 合成源句的噪声\\
    多语模型 & 语向多、低资源语向 & 共享参数，目标语言标记指定输出 & 容量竞争；译成错误语言\\
    无监督翻译 & 没有平行数据 & 跨语言初始化、去噪自编码、在线回译 & 远距离语对效果差\\
    大语言模型翻译 & 领域知识与指令要求 & 提示、少样本示例、两阶段微调 & 低资源语向弱；推断成本高\\
    质量估计重排与MBR & 最大概率译文未必最好 & 用学习型度量在候选中选择 & $N$或$N^2$次度量计算；放大度量偏好\\
    wait-$k$ & 实时输出 & 固定读写节奏，前缀到前缀训练 & 不随语序差异调整等待\\
    MILk与MMA & 固定节奏等待过多或过早 & 单调注意力学习何时读、何时写 & 训练复杂，延迟有波动\\
    上下文翻译 & 跨句指代、省略与术语 & 拼接前文或单独编码上下文 & 输入变长；未必真用到上下文\\
    DocRepair & 篇章级平行数据少 & 用目标语单语往返翻译训练修订模型 & 看不到源文\\
    \bottomrule
  \end{tabularx}
  \caption{机器翻译方法的对照。各行依次分为统计翻译、神经翻译、数据条件和大语言模型四组，最后四行处理实时与篇章两种新增条件。“新代价”列出方法本身引入的成本或局限，它们往往就是下一行方法的出发点。}
  \label{tab:nlp-generation-translation-methods}
\end{table*}

\subsection{统计机器翻译}
\label{sec:nlp-generation-smt}

\term{统计机器翻译}（statistical machine translation, SMT）从\term{噪声信道模型}（noisy channel model）出发：把源句看作目标句经过一条“信道”变形后的观测，翻译就是反推最可能的目标句，
\begin{equation}
  \hat y=\argmax_{y}\,p(y\mid x)=\argmax_{y}\,p(x\mid y)\,p_{\mathrm{LM}}(y).
  \label{eq:nlp-generation-noisy-channel}
\end{equation}
第二个等号用了贝叶斯公式并丢掉与$y$无关的$p(x)$。这样分解的好处是两部分可以分开估计：翻译模型$p(x\mid y)$只需负责“意思对不对得上”，可以从平行句对估计；语言模型$p_{\mathrm{LM}}(y)$负责“目标语是否通顺”，可以从大量单语文本估计。要把这一分解变成可用的系统，还需回答三个问题：怎样从不标对应关系的平行句子估计词对齐，怎样在对齐上构造短语翻译模型并组合更多特征，怎样在指数大的候选空间里搜索译文并调节特征权重。

\subsubsection{IBM词对齐}
\label{sec:nlp-generation-ibm}

平行语料只告诉我们哪两句互为翻译，不告诉哪个源词对应哪个目标词。\term{词对齐}（word alignment）就是这种对应关系。IBM Model 1把它作为隐变量，用词翻译概率解释观察到的目标句，再依据当前模型推断对应关系并更新概率\scite{brown-etal-1993-mathematics}。按噪声信道的方向，翻译模型应写作$p(x\mid y)$；IBM模型可以在任一方向上训练，下面为了与后文记号一致写成从$x$生成$y$，实际系统两个方向都训练，后面的对称化要用到这两份结果。为了不把句长模型与对齐计算混在一起，这里固定两句的长度。

在源句前增加$x_0=\mathrm{NULL}$，表示目标词没有显式的源词对应。隐变量$a_j\in\{0,\ldots,S\}$表示目标词$y_j$对齐的源位置，$t(v\mid u)$为源词$u$生成目标词$v$的概率，满足$\sum_v t(v\mid u)=1$。Model 1假定每个目标位置独立、均匀地选择源位置，于是
\begin{align}
  p(y,\bm{a}\mid x)
    &=\frac{1}{(S+1)^T}\prod_{j=1}^{T}t(y_j\mid x_{a_j}),\\
  p(y\mid x)
    &=\frac{1}{(S+1)^T}\prod_{j=1}^{T}
       \sum_{i=0}^{S}t(y_j\mid x_i).
  \label{eq:nlp-generation-ibm-likelihood}
\end{align}
第二式对所有$(S+1)^T$种对齐求和；由于每个目标位置独立选择源位置，求和可以逐位置分解成$T$个长为$S+1$的和。这也暴露了Model 1的限制：打乱源词次序，句子概率不变，它完全不知道语序。

参数估计用期望最大化（EM）算法。E步用当前$t$计算对齐后验
\begin{equation}
  \gamma_{j,i}
  =p(a_j=i\mid x,y)
  =\frac{t(y_j\mid x_i)}{\sum_{\ell=0}^{S}t(y_j\mid x_\ell)}.
  \label{eq:nlp-generation-ibm-posterior}
\end{equation}
同一个源词重复出现时，每个位置都进入分母和计数。把全语料中源词$u$与目标词$v$的后验相加，得到期望计数$c(u,v)$。M步在每个源词下归一化，$t_{\mathrm{new}}(v\mid u)=c(u,v)/\sum_{v'}c(u,v')$；未在语料中出现的源词分母为零，需要另定平滑或保留规则。EM每轮不降低训练语料的似然，其一般推导（以对数似然的下界交替优化）可参考\BookChapterRef{chap:pgm-parameter-learning}{卷六参数学习章}。Model 1还有一个特别的性质：对数似然关于$t$是凹函数，EM找到的是全局最大值点之一（不一定唯一），因此它常被用来给更复杂的对齐模型提供初值。

\Needspace{12\baselineskip}
\begin{example}[含空词的一轮EM]
  \label{ex:nlp-generation-em}
  教学语料只有两对句子：第一对为“我／走”与“I／go”；第二对为“我”与“I”。目标词表仅有I和go，初始概率及累计期望计数如下：
  \begin{center}
    \begin{tabular}{@{}lrrrr@{}}
      \toprule
      源词 & $t(\mathrm{I}\mid u)$ & $t(\mathrm{go}\mid u)$
        & $c(u,\mathrm{I})$ & $c(u,\mathrm{go})$\\
      \midrule
      NULL & $1/2$ & $1/2$ & $28/39$ & $1/3$\\
      我 & $4/5$ & $1/5$ & $224/195$ & $2/15$\\
      走 & $1/5$ & $4/5$ & $2/15$ & $8/15$\\
      \bottomrule
    \end{tabular}
  \end{center}
  第一对中，I的分母为$1/2+4/5+1/5=3/2$，在NULL、我、走上的后验为$(1/3,8/15,2/15)$；go的分母也为$3/2$，后验为$(1/3,2/15,8/15)$。第二对只有NULL和我，I的分母为$1/2+4/5=13/10$，后验为$(5/13,8/13)$。因此NULL到I累计$1/3+5/13=28/39$，我到I累计$8/15+8/13=224/195$，其余计数直接来自第一对。

  三个源词的计数总量分别为$41/39,50/39,2/3$，相加为3，正好等于全语料三个目标位置。M步得到
  \[
    \begin{array}{c|cc}
      &t_{\mathrm{new}}(\mathrm{I}\mid u)&t_{\mathrm{new}}(\mathrm{go}\mid u)\\
      \hline
      \mathrm{NULL}&28/41&13/41\\
      \text{我}&112/125&13/125\\
      \text{走}&1/5&4/5
    \end{array}
  \]
  “我”生成I的概率从0.8增加到0.896，因为第二对额外提供了对应证据。固定句长的语料似然从$0.25\times0.65=0.1625$变为约0.19054。第一对各目标词取最高后验时，I对齐“我”，go对齐“走”；NULL仍保留计数，下一轮EM继续参与竞争。
\end{example}

Model 1不知道语序，后续模型逐步补上位置信息。IBM Model 2增加对齐位置分布$a(i\mid j,S,T)$，使目标位置$j$更可能对齐到相对位置相近的源位置。HMM对齐模型进一步让$a_j$依赖前一个对齐$a_{j-1}$，转移概率只取决于跳跃宽度$a_j-a_{j-1}$，从而偏好“相邻目标词对齐到相邻源词”的单调片段\cite{nlp-gen-vogel1996hmm}；这是一个以源位置为状态的隐马尔可夫模型，E步用前向后向算法计算期望计数，算法细节可参考\BookChapterRef{chap:pgm-dynamic}{卷六动态有向模型章}。IBM Model 3至5再加入繁衍数（一个源词生成几个目标词）等成分。这些模型的产物是对齐，不是译文：对齐准确只说明词对应找得对，整句翻译还需要下面的短语模型和搜索。

\subsubsection{短语翻译的对数线性模型}
\label{sec:nlp-generation-phrase-loglinear}

逐词替换有两个缺陷：它不能表达“校外来宾”到external guests这种多词对一词的对应，也不能利用局部固定的语序变化。短语翻译以连续词块为翻译单位，这里的“短语”只指连续词块，不必是句法短语\scite{nlp-intro-koehn2003}。要抽取短语，先要有一份尽量好的对齐；要组合短语翻译、语言模型和重排等多种证据，又需要比式\eqref{eq:nlp-generation-noisy-channel}的两项乘积更灵活的打分形式。下面按这两个需要展开。

\paragraph{对齐对称化。}IBM模型规定每个目标词只对齐一个源词，因此单向对齐不能表示一个源词对多个目标词。两个方向各训练一次，得到$A_{x\to y}$与$A_{y\to x}$：交集精确率高但稀疏，并集召回高但含错误点。\term{对齐对称化}（alignment symmetrization）从交集出发，有选择地加入并集中的点。常用的grow-diag-final启发式如算法~\ref{alg:nlp-generation-grow-diag-final}所示\cite{nlp-gen-och2003systematic}：生长阶段只加入与已有点相邻（含对角相邻）、且至少一端尚未对齐的点；最后阶段再把并集中至少一端尚未对齐的剩余点补进来。

\Needspace{14\baselineskip}
\begin{algorithm}[grow-diag-final对齐对称化]
  \label{alg:nlp-generation-grow-diag-final}
  \begin{algorithmic}[1]
    \Require 两个方向的对齐点集$A_{x\to y}$、$A_{y\to x}$
    \Ensure 对称化对齐$A$
    \State $A\gets A_{x\to y}\cap A_{y\to x}$；$U\gets A_{x\to y}\cup A_{y\to x}$
    \Repeat\Comment{grow-diag：沿邻域生长}
      \For{每个$(i,j)\in A$及其八个邻点$(i',j')$}
        \If{$(i',j')\in U\setminus A$，且源位置$i'$或目标位置$j'$在$A$中尚未对齐}
          \State $A\gets A\cup\{(i',j')\}$
        \EndIf
      \EndFor
    \Until{$A$不再变化}
    \For{每个$(i,j)\in U\setminus A$}\Comment{final：补充孤立点}
      \If{源位置$i$或目标位置$j$在$A$中尚未对齐}
        \State $A\gets A\cup\{(i,j)\}$
      \EndIf
    \EndFor
    \State \Return $A$
  \end{algorithmic}
\end{algorithm}

以教学句对“我／今天／出发”与“I／set／off／today”为例。目标到源方向（每个目标词选一个源词）给出I--我、set--出发、off--出发、today--今天；源到目标方向（每个源词选一个目标词）给出我--I、今天--today、出发--set。交集是三个一对一的点，并集多出出发--off。生长阶段检查出发--set的邻点，出发--off在并集中且off尚未对齐，于是被加入。对称化结果同时含有“出发”对set与off，两个单向对齐都无法单独表示这个一对多的对应。

\paragraph{短语抽取与计数估计。}给定对齐$A$，源片段$\bar x=x_{i_1:i_2}$与目标片段$\bar y=y_{j_1:j_2}$与对齐一致（consistent with the alignment），当且仅当
\begin{equation}
  \forall (i,j)\in A:\ \bigl(i_1\le i\le i_2\iff j_1\le j\le j_2\bigr),
  \qquad
  \exists (i,j)\in A:\ i_1\le i\le i_2.
  \label{eq:nlp-generation-phrase-consistency}
\end{equation}
第一个条件要求片段内的词不能对齐到对方片段之外，第二个条件要求片段至少含一个对齐点。仍用上例，与对齐一致的短语对有（我，I）、（今天，today）、（出发，set off）、（今天／出发，set off today）和整句；（我／今天，I set off today）不一致，因为目标片段中的set对齐到源片段之外的“出发”；（出发，set）也不一致，因为off对齐到“出发”却落在目标片段之外。短语对（今天／出发，set off today）在一个单位里完成了“今天”从前到后的移动，逐词替换本身没有这种次序信息。片段两端未对齐的词可以并入也可以不并入，抽取时会枚举这些扩展，并限制短语最大长度（常取7）。

从全语料抽取的短语对构成短语表。双向短语翻译概率用相对频率估计，
\begin{equation}
  \phi(\bar y\mid\bar x)=\frac{c(\bar x,\bar y)}{\sum_{\bar y'}c(\bar x,\bar y')},\qquad
  \phi(\bar x\mid\bar y)=\frac{c(\bar x,\bar y)}{\sum_{\bar x'}c(\bar x',\bar y)},
  \label{eq:nlp-generation-phrase-prob}
\end{equation}
其中$c(\bar x,\bar y)$是短语对被抽取的次数。只出现过一次的长短语对会得到$\phi=1$，这样的估计并不可靠。\term{词汇加权}（lexical weighting）用词级翻译概率给短语对另打一个分：先由对称化对齐按相对频率估计词翻译表$w(y\mid x)$，再对目标短语中每个词取其所对齐源词的平均概率并连乘，
\begin{equation}
  \operatorname{lex}(\bar y\mid\bar x,A)
  =\prod_{j=1}^{|\bar y|}\frac{1}{|\{i:(i,j)\in A\}|}\sum_{(i,j)\in A}w(\bar y_j\mid\bar x_i),
  \label{eq:nlp-generation-lexical-weighting}
\end{equation}
未对齐的目标词改用$w(\bar y_j\mid\mathrm{NULL})$。例如（出发，set off）在语料中只出现一次，$\phi(\text{set off}\mid\text{出发})=1$；若$w(\text{set}\mid\text{出发})=0.3$、$w(\text{off}\mid\text{出发})=0.2$（教学数值），词汇权重为$0.06$，提醒系统这对短语的词级证据并不强。词汇加权同样双向计算。

\paragraph{对数线性打分。}有了多种证据，式\eqref{eq:nlp-generation-noisy-channel}只允许两个概率相乘，难以再加入其他特征。\term{对数线性模型}（log-linear model）把句子得分写成特征函数的加权和\cite{nlp-gen-och2002loglinear}：
\begin{equation}
  \hat y=\argmax_{y}\max_{d\in D(x,y)}\sum_{k=1}^{K}\lambda_k h_k(x,y,d).
  \label{eq:nlp-generation-loglinear}
\end{equation}
其中$d$是一个\emph{推导}：把源句切成若干短语、为每个短语选一个译文并决定它们在目标句中的顺序；$D(x,y)$是能产生$y$的所有推导，$\max_d$用最好推导近似对推导求和。当$K=2$、$h_1=\log p(x\mid y)$、$h_2=\log p_{\mathrm{LM}}(y)$、$\lambda_1=\lambda_2=1$时，式\eqref{eq:nlp-generation-loglinear}退化为噪声信道模型，因此对数线性模型是它的推广。典型的特征包括：两个方向的短语翻译对数概率$\sum\log\phi$与两个方向的词汇加权$\sum\log\operatorname{lex}$，共四项；目标语言模型$\log p_{\mathrm{LM}}(y)$；\term{扭曲代价}（distortion cost）
\begin{equation}
  h_{\mathrm{dist}}(d)=-\sum_{r=1}^{R}\bigl|\mathrm{start}_r-\mathrm{end}_{r-1}-1\bigr|,
  \label{eq:nlp-generation-distortion}
\end{equation}
其中$\mathrm{start}_r$与$\mathrm{end}_{r-1}$分别是第$r$个被翻译的源短语的起点和前一个源短语的终点，$\mathrm{end}_0=0$，源短语按顺序接续时代价为零；以及词数惩罚$h=T$和短语数惩罚$h=R$，二者的权重用来调节译文长短和切分粗细。

用“我／今天／出发”比较两个推导。推导甲逐词翻译并在目标端重排：依次翻译“我”（源位置1）、“出发”（位置3）、“今天”（位置2），三次跳跃代价为$|1-0-1|+|3-1-1|+|2-3-1|=0+1+2=3$。推导乙用短语（今天／出发，set off today），依次覆盖位置1和位置2至3，跳跃代价为0。两个推导产生同一个译文，乙的重排藏在短语内部，扭曲代价更低。这正是短语模型的优势：局部的语序变化由短语表免费提供，扭曲代价只负责短语之间的移动。反过来，扭曲代价只看跳跃距离，不看跳的是什么词，长距离重排（如日语句末动词提到英语主语之后）只能靠语言模型勉强补救，这是短语模型的主要弱点之一。

\subsubsection{栈解码与MERT调参}
\label{sec:nlp-generation-stack-mert}

即使只考虑单词级的重排，在式\eqref{eq:nlp-generation-loglinear}下精确求最优译文也是NP难的\cite{nlp-gen-knight1999decoding}，实际系统采用近似搜索。\term{栈解码}（stack decoding）的思路是从左到右逐步生成目标句，每一步选一个尚未翻译的源短语并给出译文；部分译文称为\emph{假设}，记录已覆盖的源位置集合、已生成的目标词和累计得分。假设按已覆盖的源词数放入$S+1$个栈，第$s$个栈只比较覆盖了$s$个源词的假设，以免把刚开始的假设与快要完成的假设放在一起比较。

只按累计得分在同一栈内比较仍不公平：覆盖相同词数的两个假设，可能一个先翻译了容易的词，一个先翻译了难的词，前者的累计得分高只是因为难的部分还没算。\term{未来代价}（future cost）为每个假设估计“剩余源词至少还要付出多少得分”，比较时用累计得分加未来代价。估计时忽略扭曲代价与语言模型的上下文，对每个源区间$[i,j]$取能覆盖它的最好短语选项的得分，再用动态规划允许把区间切成两段，
\begin{equation}
  \mathrm{fc}(i,j)=\max\Bigl\{\max_{\text{覆盖}[i,j]\text{的短语选项}}\mathrm{score},\ \max_{i\le k<j}\bigl[\mathrm{fc}(i,k)+\mathrm{fc}(k+1,j)\bigr]\Bigr\}.
  \label{eq:nlp-generation-future-cost}
\end{equation}
一个假设的未来代价是其所有未覆盖的极大连续区间的$\mathrm{fc}$之和。这张表对每个源句只需算一次。教学例子：两个假设都在第1个栈里，假设甲先翻译了容易的“出发”，累计$-0.4$，剩下“我／今天”，未来代价$-2.0$；假设乙翻译了“我”，累计$-1.0$，剩下“今天／出发”，未来代价$-0.8$。只看累计得分会保留甲；加上未来代价后乙为$-1.8$，优于甲的$-2.4$，乙其实是更有希望的部分译文。未覆盖部分不连续时，例如只覆盖了中间的“今天”，未来代价取“我”与“出发”两个区间的$\mathrm{fc}$之和。

除剪枝外，栈解码还做\term{假设重组}（hypothesis recombination）：两个假设若覆盖集合相同、最后$n-1$个目标词相同（$n$为语言模型阶数）、最后一个源短语的终点相同，它们此后的所有扩展得分完全相同，只保留得分高的一个不会丢掉最优推导。算法~\ref{alg:nlp-generation-stack-decoding}汇总全过程\scite{nlp-gen-koehn2010smt}。

\Needspace{16\baselineskip}
\begin{algorithm}[短语翻译的栈解码]
  \label{alg:nlp-generation-stack-decoding}
  \begin{algorithmic}[1]
    \Require 源句$x$，短语表，特征权重$\bm\lambda$，栈容量$b$，扭曲上限$\delta$
    \Ensure 得分最高的完整译文
    \State 按式\eqref{eq:nlp-generation-future-cost}计算所有区间的未来代价表$\mathrm{fc}$
    \State 栈$\mathcal S_0\gets\{$空假设$\}$；$\mathcal S_1,\ldots,\mathcal S_S\gets\varnothing$
    \For{$s=0,\ldots,S-1$}
      \State 把$\mathcal S_s$中的假设按“累计得分$+$未来代价”排序，只保留前$b$个
      \For{每个假设$h\in\mathcal S_s$}
        \For{每个与$h$的覆盖不相交、跳跃距离不超过$\delta$的源区间及其短语选项$o$}
          \State $h'\gets$扩展$h$：并入覆盖，追加译文，按式\eqref{eq:nlp-generation-loglinear}累加特征得分
          \State 记$|h'|$为$h'$已覆盖的源词数；若$\mathcal S_{|h'|}$中已有可与$h'$重组的假设，只保留得分较高者，否则把$h'$加入$\mathcal S_{|h'|}$
        \EndFor
      \EndFor
    \EndFor
    \State \Return $\mathcal S_S$中累计得分最高的假设所对应的译文
  \end{algorithmic}
\end{algorithm}

每个栈至多$b$个假设，每个假设的扩展数受扭曲上限和短语选项数限制，因此解码代价约为$O(S\cdot b\cdot S\delta\cdot o_{\max})$，其中$o_{\max}$是每个源区间保留的短语选项数。栈容量和扭曲上限越小越快，搜索错误也越多。

特征权重$\bm\lambda$还需要确定。最终评价看的是BLEU（定义见第~\ref{sec:nlp-generation-overlap}小节）等译文质量指标，而不是似然，所以希望直接选使开发集BLEU最高的权重。困难在于BLEU作为$\bm\lambda$的函数是分段常数：权重小幅变化时，每句选出的最优译文不变，BLEU也不变，梯度处处为零或不存在。\term{最小错误率训练}（minimum error rate training, MERT）利用这一结构做精确的线搜索\scite{nlp-gen-och2003mert}。固定当前权重$\bm\lambda$和搜索方向$\bm d$，对开发集每句的$n$-best列表，候选$e$在直线$\bm\lambda+\gamma\bm d$上的得分是$\gamma$的一次函数
\begin{equation}
  s_e(\gamma)=\underbrace{\bm\lambda^{\top}\bm h(e)}_{a_e}+\gamma\,\underbrace{\bm d^{\top}\bm h(e)}_{b_e}.
  \label{eq:nlp-generation-mert-line}
\end{equation}
每句被选中的候选是这些直线的上包络，随$\gamma$变化只在包络的交点处切换。把所有句子的切换点合并排序，相邻切换点之间每句的选择都不变，BLEU的充分统计量（各阶n元的匹配数与长度）可以逐区间增量更新，于是能精确找到这条直线上BLEU最高的区间，取区间中点作为新的$\gamma$。

教学例子：某句有三个候选，$(a_e,b_e)$分别为$(-2.0,0)$、$(-2.5,1)$、$(-4.0,2)$，句级错误（教学中以$1-$句级质量代替）分别为$0.6,0.2,0.4$。第一、二条直线在$\gamma=0.5$相交，第二、三条在$\gamma=1.5$相交，第一、三条在$\gamma=1$相交但此时第二条更高。上包络因此为：$\gamma<0.5$选候选一，$0.5<\gamma<1.5$选候选二，$\gamma>1.5$选候选三。错误最低的区间是$(0.5,1.5)$，取$\gamma=1$。真实MERT对整个开发集的语料级BLEU做同样的合并。

\Needspace{12\baselineskip}
\begin{algorithm}[最小错误率训练]
  \label{alg:nlp-generation-mert}
  \begin{algorithmic}[1]
    \Require 开发集，初始权重$\bm\lambda$，解码器
    \Ensure 调好的权重$\bm\lambda$
    \State 用$\bm\lambda$解码开发集，得到各句$n$-best列表$C$
    \Repeat
      \For{若干方向$\bm d$（各坐标方向及随机方向）}
        \State 对每句候选按式\eqref{eq:nlp-generation-mert-line}计算$(a_e,b_e)$，求上包络的切换点
        \State 合并全部切换点，逐区间更新BLEU统计量，取BLEU最高区间的中点$\gamma^\ast$
        \State $\bm\lambda\gets\bm\lambda+\gamma^\ast\bm d$
      \EndFor
      \State 用新$\bm\lambda$重新解码，把新的$n$-best候选并入$C$
    \Until{$C$不再增加新候选}
    \State \Return $\bm\lambda$
  \end{algorithmic}
\end{algorithm}

外层循环必不可少：$n$-best列表只是全部候选的一个样本，权重变了，解码器会找到列表外的新候选。MERT的目标函数非凸、分段常数，只适合十几到几十个稠密特征；特征多到上千时，常改用成对排序（PRO）或在线大间隔方法（MIRA）。

回看整套系统：对齐、短语抽取、各概率表、语言模型、扭曲模型分别训练，只有十几个权重在最终指标上联合调节；某一环节的错误（例如对齐错导致短语表错）后面无法纠正，长距离重排也只能由距离代价和$n$元语言模型间接处理。神经翻译要解决的正是这种分别优化的问题。

\subsection{神经机器翻译}
\label{sec:nlp-generation-neural-translation}

\term{神经机器翻译}（neural machine translation, NMT）直接建模条件分布
\begin{equation}
  p_{\bm\theta}(y\mid x)=\prod_{t=1}^{T}p_{\bm\theta}(y_t\mid y_{<t},x),
  \label{eq:nlp-generation-nmt}
\end{equation}
用一个网络同时承担短语表、重排模型和语言模型的作用，所有参数在平行语料上按最大似然联合训练，训练时以参考前缀为条件逐位置计算交叉熵（教师强制，掩码损失的写法见第~\ref{sec:nlp-sft-masked-loss}小节）。这正好回应了统计翻译分别优化的问题。编码器—解码器与交叉注意力的一般计算、批张量形状和键值缓存已在第~\ref{sec:nlp-generation-encoder-decoder}小节给出，下面按结构演进只写翻译模型特有的选择，以及每一步解决了前一步的哪个问题。

\subsubsection{循环编码器—解码器与注意力}
\label{sec:nlp-generation-rnn-attention-mt}

最早的端到端翻译模型用两个循环网络：编码器逐词读入源句，把最后的隐状态当作整句的表示；解码器以这个表示为初始条件，逐词生成目标句\cite{sutskever2014seq2seq,cho2014}。设源词嵌入$\bm e_s=\bm E_{\mathrm{src}}[x_s]\in\mathbb R^{d_e}$，编码器与解码器的隐状态维数为$d$，
\begin{equation}
  \begin{gathered}
  \bm h_s=f_{\mathrm{enc}}(\bm h_{s-1},\bm e_s),\qquad \bm c=\bm h_S,\\
  \bm s_t=f_{\mathrm{dec}}\bigl(\bm s_{t-1},\bm E_{\mathrm{tgt}}[y_{t-1}],\bm c\bigr),\qquad
  p(y_t\mid y_{<t},x)=\operatorname{softmax}(\bm W_o\bm s_t+\bm b_o),
  \end{gathered}
  \label{eq:nlp-generation-rnn-seq2seq}
\end{equation}
其中$f_{\mathrm{enc}},f_{\mathrm{dec}}$是LSTM或GRU单元，$\bm W_o\in\mathbb R^{V\times d}$把状态映射到目标词表的$V$个得分，$y_0$为起始标记BOS。

这个模型的问题在$\bm c$：无论源句有5个词还是50个词，解码器能用的源信息都只有一个$d$维向量，即第~\ref{sec:nlp-generation-encoder-decoder}小节所说的\term{定长瓶颈}（fixed-length bottleneck）。Sutskever等发现把源句倒序输入能明显提升效果，倒序后句首词离解码起点更近；这个技巧有效，本身就说明信息在长距离传递中衰减。

加性注意力（additive attention）让解码器每一步重新读取全部源位置，而不是只读$\bm c$。第~\ref{sec:nlp-generation-encoder-decoder}小节式\eqref{eq:nlp-generation-bahdanau}给出了它的一般形式，这里写出Bahdanau等原始翻译模型的完整计算\scite{bahdanau2015}。编码器改为双向循环网络，第$i$个位置的表示$\bm h_i=[\overrightarrow{\bm h}_i;\overleftarrow{\bm h}_i]\in\mathbb R^{2d}$同时含左右上下文。生成第$t$个目标词时，先用上一步的解码状态$\bm s_{t-1}$给每个源位置打分，归一化为权重，再加权求和得到本步的\term{上下文向量}（context vector）：
\begin{equation}
  \begin{gathered}
  e_{t,i}=\bm v_a^{\top}\tanh\bigl(\bm W_a\bm s_{t-1}+\bm U_a\bm h_i\bigr),\qquad
  \alpha_{t,i}=\frac{\exp(e_{t,i})}{\sum_{k=1}^{S}\exp(e_{t,k})},\\
  \bm c_t=\sum_{i=1}^{S}\alpha_{t,i}\bm h_i,\qquad
  \bm s_t=f_{\mathrm{dec}}\bigl(\bm s_{t-1},\bm E_{\mathrm{tgt}}[y_{t-1}],\bm c_t\bigr),\\
  p(y_t\mid y_{<t},x)=g\bigl(\bm s_t,\bm E_{\mathrm{tgt}}[y_{t-1}],\bm c_t\bigr).
  \end{gathered}
  \label{eq:nlp-generation-bahdanau-mt}
\end{equation}
其中$\bm W_a\in\mathbb R^{d_a\times d}$、$\bm U_a\in\mathbb R^{d_a\times 2d}$把解码状态与源表示投到同一个$d_a$维打分空间，$\bm v_a\in\mathbb R^{d_a}$把它压成标量；$g$是含线性层与softmax的输出网络。与式\eqref{eq:nlp-generation-rnn-seq2seq}相比，唯一的改变是固定的$\bm c$换成了每步不同的$\bm c_t$。与一般形式相比，原始模型有两处具体选择：查询用上一步状态$\bm s_{t-1}$，因此可以先算出$\bm c_t$再更新状态，计算次序不会互相等待；$\bm c_t$既进入状态更新，也直接进入输出层。

按批大小$B$写出一步的形状：源表示$\bm H\in\mathbb R^{B\times S\times 2d}$，投影$\bm H\bm U_a^{\top}\in\mathbb R^{B\times S\times d_a}$与解码步无关，对每句只算一次；每步只需计算$\bm W_a\bm s_{t-1}\in\mathbb R^{B\times d_a}$，广播到$S$个位置相加后过$\tanh$并与$\bm v_a$内积，得到$B\times S$的得分和权重，再与$\bm H$相乘得到$B\times 2d$的$\bm c_t$。每步代价$O(Sd_a)$，整句$O(STd_a)$，换来的是解码器在每一步都能直接访问每个源位置。

权重$\alpha_{t,\cdot}$构成一个$T\times S$的矩阵，第$t$行说明生成$y_t$时主要读了哪些源词，可视化后常与词对齐大体吻合，因此也称\emph{软对齐}。与IBM模型对比：IBM的对齐是离散隐变量，单独训练后交给短语抽取；注意力权重是连续的，与翻译目标一起训练，不需要对齐标注，但也不受“必须是对齐”的约束，个别词上可能偏移一两个位置。Luong等把打分函数换成点积或双线性形式$\bm s_t^{\top}\bm W\bm h_i$，并用当前状态$\bm s_t$而非$\bm s_{t-1}$做查询，计算更简单\scite{luong2015}。

注意力解决了瓶颈，循环结构的另一个代价仍在：编码第$s$个位置必须等第$s-1$个位置算完，解码训练时也要按$t=1,\ldots,T$顺序展开。一个句对的前向需要约$S+T$个串行步骤，句内无法并行，长语料上的训练时间因此很长。

\subsubsection{Transformer翻译模型}
\label{sec:nlp-generation-transformer-mt}

Transformer去掉循环，把“读取其他位置”全部交给注意力\scite{vaswani2017}。它最初正是一个翻译模型：编码器与解码器各$L$层，编码器层由自注意力和前馈网络组成，解码器层由因果自注意力、交叉注意力和前馈网络组成，每个子层带残差连接与层归一化（LayerNorm, LN）。图~\ref{fig:nlp-generation-transformer-mt}画出整体结构，下面逐层写出计算。

\begin{figure}[htbp]
  \centering
  % 第九卷 12.6.2.2 节：Transformer 翻译模型（编码器—解码器）。
% 左列编码器、右列解码器，主链路自下而上；灰字为省略批维的张量形状。
% 蓝色：注意力子层；黄色：前馈子层；红色：输出；白底灰框：嵌入与位置编码。
\begin{tikzpicture}[
  x=1mm,y=1mm,
  every node/.style={font=\figurefont\footnotesize,text=FigureInk},
  op/.style={rectangle,rounded corners=1.8pt,draw=FigureStroke,line width=0.6pt,fill=white,
    minimum height=7.5mm,minimum width=36mm,inner sep=1.2mm,align=center},
  attnop/.style={op,draw=FigureBlue,fill=FigureBlueFill},
  ffnop/.style={op,draw=FigureYellow,fill=FigureYellowFill},
  outop/.style={op,draw=FigureRed,fill=FigureRedFill},
  tok/.style={font=\figurefont\footnotesize,text=FigureInk,inner sep=0.8pt},
  flow/.style={-{Stealth[length=1.8mm,width=1.3mm]},draw=FigureStroke,line width=0.7pt},
  kv/.style={-{Stealth[length=1.8mm,width=1.3mm]},draw=FigureBlue,line width=0.9pt,rounded corners=2pt},
  dim/.style={font=\figurefont\scriptsize,text=FigureMuted,inner sep=0.8pt},
  panel/.style={rounded corners=2.5pt,draw=FigureGrid,line width=0.5pt,fill=FigurePanel},
  ptitle/.style={font=\figurefont\scriptsize\bfseries,text=FigureMuted}
]
  % 层堆叠背景。
  \draw[panel] (0,20) rectangle (44,51);
  \draw[panel] (58,20) rectangle (102,70.5);
  \node[ptitle,anchor=north west] at (0.8,50.2) {编码器层$\times L$};
  \node[ptitle,anchor=north west] at (58.8,69.7) {解码器层$\times L$};

  % 编码器。
  \node[tok] (src) at (22,0) {源句：他\ 不\ 参加\ 会议};
  \node[op] (eemb) at (22,11) {$\sqrt d\,\bm E[x]+\bm P$};
  \node[attnop] (esa) at (22,28) {多头自注意力（双向）\\残差与层归一化};
  \node[ffnop] (effn) at (22,42) {前馈网络\\残差与层归一化};
  \node[op,draw=none,fill=none] (henc) at (22,60) {编码器输出$\bm H_{\mathrm{enc}}$};
  \draw[flow] (src) -- (eemb);
  \draw[flow] (eemb) -- (esa);
  \draw[flow] (esa) -- (effn);
  \draw[flow] (effn) -- (henc);
  \node[dim,anchor=west] at (23,17.5) {$S\times d$};
  \node[dim,anchor=west] at (23,54.5) {$S\times d$};
  \node[dim,anchor=east] at (21,21.8) {权重$S\times S$};

  % 解码器。
  \node[tok] (tgt) at (80,0) {右移目标：BOS\ He\ will\ not};
  \node[op] (demb) at (80,11) {$\sqrt d\,\bm E[y_{<t}]+\bm P$};
  \node[attnop] (dsa) at (80,28) {因果自注意力\\残差与层归一化};
  \node[attnop] (dca) at (80,44) {交叉注意力\\残差与层归一化};
  \node[ffnop] (dffn) at (80,59) {前馈网络\\残差与层归一化};
  \node[outop] (out) at (80,78) {乘$\bm E^{\top}$，softmax};
  \node[tok] (pred) at (80,90) {预测：He\ will\ not\ attend};
  \draw[flow] (tgt) -- (demb);
  \draw[flow] (demb) -- (dsa);
  \draw[flow] (dsa) -- node[dim,right=0.4mm,pos=0.5] {查询} (dca);
  \draw[flow] (dca) -- (dffn);
  \draw[flow] (dffn) -- (out);
  \draw[flow] (out) -- (pred);
  \node[dim,anchor=west] at (81,17.5) {$T\times d$};
  \node[dim,anchor=east] at (79,21.8) {权重$T\times T$};
  \node[dim,anchor=east] at (79,37.2) {权重$T\times S$};
  \node[dim,anchor=west] at (81,72.5) {$T\times d$};
  \node[dim,anchor=west] at (81,84) {$T\times V$};

  % 编码器输出送入每个解码层的交叉注意力。
  \draw[kv] (henc.east) -- (51,60) -- (51,44) -- (dca.west);
  \node[dim,anchor=south,text=FigureBlue,align=center] at (44.5,60.6) {键、值\\各层共用};
\end{tikzpicture}

  \caption{Transformer翻译模型。左列为编码器，读入完整源句，各位置可以双向互相注意；右列为解码器，输入右移一位的目标前缀，因果自注意力只看已生成的位置，交叉注意力以解码器状态为查询、以编码器最终输出为键和值。两列中的大框各重复$L$层。灰字为省略批维后的张量形状，$S$与$T$分别为源、目标长度；输出层与两端嵌入共享同一矩阵$\bm E$。}
  \label{fig:nlp-generation-transformer-mt}
\end{figure}

\paragraph{编码器。}源词元经嵌入并乘以$\sqrt d$，加上正弦位置编码$\bm P$，得到$\bm H^{(0)}=\sqrt d\,\bm E[x]+\bm P_{1:S}\in\mathbb R^{S\times d}$。第$\ell$层（原始论文的后置归一化写法）为
\begin{equation}
  \begin{aligned}
  \bm U^{(\ell)}&=\operatorname{LN}\bigl(\bm H^{(\ell)}+\operatorname{MHA}(\bm H^{(\ell)},\bm H^{(\ell)};\bm M_{\mathrm{src}})\bigr),\\
  \bm H^{(\ell+1)}&=\operatorname{LN}\bigl(\bm U^{(\ell)}+\operatorname{FFN}(\bm U^{(\ell)})\bigr),
  \end{aligned}
  \label{eq:nlp-generation-transformer-encoder}
\end{equation}
其中$\operatorname{MHA}(\bm Z,\bm C;\bm M)$表示以$\bm Z$产生查询、以$\bm C$产生键和值、在得分上加掩码$\bm M$的多头注意力：每个头按式\eqref{eq:nlp-llm-cross-attention}用$\bm W_Q^a,\bm W_K^a,\bm W_V^a\in\mathbb R^{d\times d_h}$计算缩放点积注意力，$h$个头的输出按式\eqref{eq:nlp-llm-mha}拼接后乘$\bm W_O\in\mathbb R^{hd_h\times d}$；$\bm Z=\bm C$时就是自注意力。前馈网络$\operatorname{FFN}(\bm u)=\bm W_2\operatorname{ReLU}(\bm W_1\bm u+\bm b_1)+\bm b_2$逐位置作用，$\bm W_1\in\mathbb R^{d_{\mathrm{ff}}\times d}$。编码器的掩码$\bm M_{\mathrm{src}}$只把批内填充位置所在的列置为$-\infty$，其余为0，因此每个源位置可以看到整句的左右两侧。$L$层之后的输出记为$\bm H_{\mathrm{enc}}=\bm H^{(L)}$。

\paragraph{解码器。}目标端输入是右移一位的参考$(\mathrm{BOS},y_1,\ldots,y_{T-1})$，同样嵌入并加位置编码得$\bm G^{(0)}\in\mathbb R^{T\times d}$。第$\ell$层依次计算
\begin{equation}
  \begin{aligned}
    \bm A^{(\ell)}&=\operatorname{LN}\bigl(\bm G^{(\ell)}+\operatorname{MHA}(\bm G^{(\ell)},\bm G^{(\ell)};\bm M_{\mathrm{causal}})\bigr),\\
    \bm B^{(\ell)}&=\operatorname{LN}\bigl(\bm A^{(\ell)}+\operatorname{MHA}(\bm A^{(\ell)},\bm H_{\mathrm{enc}};\bm M_{\mathrm{src}})\bigr),\\
    \bm G^{(\ell+1)}&=\operatorname{LN}\bigl(\bm B^{(\ell)}+\operatorname{FFN}(\bm B^{(\ell)})\bigr).
  \end{aligned}
  \label{eq:nlp-generation-transformer-decoder}
\end{equation}
第一行的因果掩码$(\bm M_{\mathrm{causal}})_{t,j}=-\infty\ (j>t)$保证位置$t$只看自己和之前的目标词；第二行是交叉注意力，查询来自解码器的$\bm A^{(\ell)}$，键和值来自编码器输出$\bm H_{\mathrm{enc}}$，每个解码层都读一次同一个$\bm H_{\mathrm{enc}}$。最后一层输出经线性层和softmax给出下一词分布，
\begin{equation}
  p(y_t\mid y_{<t},x)=\operatorname{softmax}\bigl(\bm G^{(L)}_{t,:}\bm E^{\top}\bigr),
  \label{eq:nlp-generation-transformer-output}
\end{equation}
这里输出投影直接用嵌入矩阵$\bm E\in\mathbb R^{V\times d}$的转置；原论文在源、目标共享子词词表时让源嵌入、目标嵌入与输出投影三处共用同一矩阵。

\paragraph{参数量。}批张量形状见图~\ref{fig:nlp-generation-transformer-mt}与第~\ref{sec:nlp-generation-encoder-decoder}小节，编码器自注意力权重另为$B\times h\times S\times S$。这里用原论文的基础配置核算参数量：$L=6$，$d=512$，$h=8$，$d_h=64$，$d_{\mathrm{ff}}=2048$，英德任务共享约37000个子词。共享的嵌入矩阵约$37000\times512\approx1.89\times10^7$个参数；一个注意力子层有$\bm W_Q,\bm W_K,\bm W_V,\bm W_O$四个$512\times512$矩阵，约$1.05\times10^6$；一个前馈子层有两个$512\times2048$矩阵，约$2.10\times10^6$。编码器层含一个注意力和一个前馈，约$3.15\times10^6$，6层共$1.89\times10^7$；解码器层多一个交叉注意力，约$4.19\times10^6$，6层共$2.52\times10^7$。三项相加约$6.3\times10^7$，偏置与归一化参数只再增加十万个左右，与原论文报告的约6500万相近。其中嵌入占近三成，这也说明共享嵌入矩阵能省下可观的参数。

\paragraph{训练与推断。}训练时目标句整句已知，式\eqref{eq:nlp-generation-transformer-decoder}对全部$T$个位置一次算完，因果掩码保证每个位置的预测只依赖它之前的参考词；所有位置的交叉熵相加即为损失。与循环模型的$S+T$个串行步骤相比，句内全部并行。推断时解码器仍逐词生成，编码器只运行一次，交叉注意力的键值可缓存复用。现在的实现多把层归一化移到子层之前（前置归一化），深层训练更稳定；式\eqref{eq:nlp-generation-transformer-decoder}与式\eqref{eq:nlp-generation-decoder-layer}的区别只在归一化的位置。

与加性注意力相比，Transformer的交叉注意力有三处变化：打分从加性网络换成缩放点积，便于批量矩阵乘；单头变成$h$头，各头可以关注不同的源位置；每个解码层都有一次交叉注意力，而不是整个解码器只有一次。代价是每层注意力的计算与显存分别随$S^2$、$T^2$和$ST$增长，长文档翻译时这一项会成为主要开销；同时模型不再有循环带来的语序偏置，位置信息完全依赖位置编码。

\subsubsection{子词词表与标签平滑}
\label{sec:nlp-generation-subword-label-smoothing}

前面的模型都假定词表封闭。按词建模时，词表通常截断到三五万个高频词，其余词替换为未知词标记UNK；人名、机构名、数字串和形态变化词恰好多在截断之外，而这些正是翻译最不能错的地方。Sennrich等把字节对编码（BPE）引入翻译，反复合并语料中最常见的相邻符号对得到子词词表，高频词成为单个子词，罕见词被拆成几个较常见的片段，任何词都能被表示\scite{nlp-llm-sennrich2016bpe}。合并算法与词表大小的选择见第~\ref{sec:nlp-llm-bpe}小节。

翻译场景多出一个选择：分别为源语和目标语学两套BPE，还是在两种语言拼接后的语料上学一套。\term{联合BPE}（joint BPE）采用后者，使同一个人名、数字或代码在两侧被切成相同的片段。例如活动编号\code{ML-0418}在中文源句和英文译文中原样出现，联合切分后两侧的子词序列完全一致，模型只需学会逐片段复制，而不必为两套不同的切分学一个映射。联合BPE得到的\term{共享词表}（shared vocabulary）还允许源嵌入、目标嵌入和输出投影共用一个矩阵（式\eqref{eq:nlp-generation-transformer-output}），参数更少；输出投影与输入嵌入共享权重的做法最早在语言模型中提出\scite{nlp-llm-press2017tying}。对中英这种文字不重叠的语对，共享词表的收益主要来自数字、拉丁字母专名和代码，代价是两种语言分享一个容量固定的词表，每种语言分到的子词变少。

第二个问题出在训练目标上。逐词交叉熵的目标分布是独热的：只有参考词概率为1。最小化它会把参考词的概率推向1，使模型输出过度自信，而翻译中同一位置常有多个合理的词。\term{标签平滑}（label smoothing）把目标分布换成独热分布与均匀分布的混合\cite{nlp-gen-szegedy2016inception}：
\begin{equation}
  \begin{gathered}
  q_\varepsilon(v)=(1-\varepsilon)\,\mathbb 1[v=y_t]+\frac{\varepsilon}{V},\\
  \mathcal L_{\mathrm{LS}}=-\sum_{v=1}^{V}q_\varepsilon(v)\log p(v)
  =(1-\varepsilon)\bigl(-\log p(y_t)\bigr)+\varepsilon\cdot\frac1V\sum_{v=1}^{V}\bigl(-\log p(v)\bigr).
  \end{gathered}
  \label{eq:nlp-generation-label-smoothing}
\end{equation}
第一项是原来的交叉熵，第二项是均匀分布对模型分布的交叉熵，它等于均匀分布的熵加上$\mathrm{KL}(u\,\|\,p)$，惩罚离均匀分布太远的输出。交叉熵在$p=q_\varepsilon$时取最小，因此参考词的最优概率是$1-\varepsilon+\varepsilon/V$而不是1。教学数值：$V=4$、$\varepsilon=0.1$时，最优分布为$(0.925,0.025,0.025,0.025)$，此时$\mathcal L_{\mathrm{LS}}\approx0.349$。若模型输出$(0.97,0.01,0.01,0.01)$，普通交叉熵为$0.030$，比最优分布的$-\log0.925\approx0.078$更低；平滑后的损失却升到$0.374$，高于$0.349$。标签平滑因此使模型停在更不自信的位置。

Transformer翻译取$\varepsilon=0.1$，原论文报告这使困惑度变差（模型不再追求参考词概率最大），BLEU却提高\scite{vaswani2017}。Müller等的分析表明，标签平滑改善了概率校准，并使同一类别的倒数第二层表示聚得更紧；后一点也有代价，用平滑训练的模型当知识蒸馏的教师时，传给学生的类间相似信息变少\scite{nlp-gen-muller2019labelsmoothing}。

\subsubsection{束搜索的长度与覆盖惩罚}
\label{sec:nlp-generation-mt-beam-penalty}

训练好的模型仍需解码。翻译通常用束搜索（第~\ref{sec:nlp-generation-beam}小节）近似求式\eqref{eq:nlp-generation-nmt}的最大概率译文，第~\ref{sec:nlp-inference-length-penalty}小节已说明对数概率逐词相加使它偏好短句。翻译中这一偏差随束宽放大：Koehn与Knowles观察到，束宽超过一定大小后BLEU反而下降，译文明显变短\scite{nlp-inference-koehn2017six}。翻译还有第二种偏差，即漏译：模型可能流畅地跳过一个源词，尤其是否定词、数量词这类短词，而漏掉它的候选概率并不一定低。

GNMT的排序分数式\eqref{eq:nlp-inference-gnmt}同时处理这两点\cite{nlp-inference-wu2016gnmt}：累计对数概率$\ell(y)$除以长度项$\mathrm{lp}(y)=((5+|y|)/6)^{\alpha}$，再加覆盖项$\mathrm{cp}(x;y)=\beta\sum_i\log\min(\sum_ja_{i,j},1)$。覆盖项中$\sum_ja_{i,j}$是整条译文对第$i$个源词的累计注意力，正好利用了编码器—解码器的交叉注意力：累计达到1的源词不扣分，从未被注意的源词使对数项趋向$-\infty$。在Transformer中注意力有多层多头，需要约定取哪一层、哪些头的平均作为$a_{i,j}$。

教学例子：源句“他／不／参加／会议”，两个完成候选为A“He will attend the meeting EOS”（6个词元，漏译“不”）和B“He will not attend the meeting EOS”（7个词元），长度按第~\ref{sec:nlp-inference-length-penalty}小节的约定计入EOS。设累计对数概率$\ell_A=-2.0$、$\ell_B=-2.4$，四个源词得到的累计注意力在A中为$(1.1,0.3,1.2,1.4)$，在B中为$(1.0,0.9,1.1,1.0)$，取$\alpha=0.6$、$\beta=0.2$（均为教学数值）。
\begin{center}
  \begin{tabular}{@{}lrrrrr@{}}
    \toprule
    候选 & $\ell$ & $\mathrm{lp}$ & $\ell/\mathrm{lp}$ & $\mathrm{cp}$ & 总分\\
    \midrule
    A（漏译） & $-2.000$ & $1.439$ & $-1.390$ & $-0.241$ & $-1.631$\\
    B（完整） & $-2.400$ & $1.516$ & $-1.583$ & $-0.021$ & $-1.604$\\
    \bottomrule
  \end{tabular}
\end{center}
其中$\mathrm{cp}_A=0.2\log0.3$，$\mathrm{cp}_B=0.2\log0.9$。只按$\ell$或只加长度惩罚，A都胜出；加入覆盖惩罚后B以微弱优势胜出。这个例子也说明了覆盖惩罚的局限：它依赖注意力权重确实反映“翻译了哪个词”，而注意力只是软对齐；$\beta$过大时覆盖项会压过模型概率，偏好把注意力摊到每个源词上的候选，即使它们译得并不好。$\alpha,\beta$都要在开发集上调节。

\subsection{低资源与多语翻译}
\label{sec:nlp-generation-low-resource-mt}

上面的神经模型需要几百万句对才能训练好，而世界上多数语向的平行数据只有几万句甚至没有。NLLB项目为200种语言构建数据、模型和评测集，WMT每年的通用翻译评测则覆盖高、中、低资源的若干语向，二者是这一问题的数据与评测入口\cite{nlp-intro-nllb2022,nlp-intro-kocmi2024}。回译、多语模型与无监督翻译依次放宽数据要求，代价分别是合成数据的噪声、语言间的容量竞争，以及对两种语言单语数据可比性的依赖。

\subsubsection{回译及其变体}
\label{sec:nlp-generation-back-translation}

平行语料不足而目标语单语较多时，\term{回译}（back-translation）用反向模型构造额外的训练对\scite{sennrich-etal-2016-improving}。先用已有平行数据$\mathcal D_{\mathrm{par}}$训练反向模型$q_{\bm\phi}(x\mid y)$，把目标语单语集合$\mathcal M_y$中的每句$y$译成合成源句$\tilde x$，再与真实数据一起训练正向模型：
\begin{equation}
  \mathcal L(\bm\theta)=\sum_{(x,y)\in\mathcal D_{\mathrm{par}}}-\log p_{\bm\theta}(y\mid x)
  +\lambda\sum_{y\in\mathcal M_y,\ \tilde x\sim q_{\bm\phi}(\cdot\mid y)}-\log p_{\bm\theta}(y\mid\tilde x).
  \label{eq:nlp-generation-back-translation}
\end{equation}
关键在于目标端$y$是真实的人写文本，噪声只在源端。正向模型学习的是“从可能有误的源句生成通顺的目标句”，解码器得到了大量真实目标语的训练，这与统计翻译中单独训练目标语言模型的作用相近。$\lambda$或真实与合成数据的采样比例在开发集上确定。例如英文单语“Do not modify the 12 records”经反向模型得到“不要修改12份记录”，正向训练仍以原英文为目标。

合成源句怎样生成影响很大。用束搜索或贪心生成$\tilde x$，得到的是反向模型的高概率输出，句式规整、彼此相似，对正向模型而言太容易预测。Edunov等比较了几种生成方式\cite{nlp-gen-edunov2018bt}：从$q_{\bm\phi}$中\emph{采样}，或对束搜索输出\emph{加噪}（以0.1的概率删词、以0.1的概率把词替换为占位符、在3个位置的范围内打乱词序），在大规模英德实验中都明显优于束搜索与贪心；他们的解释是，采样和加噪让合成源句更多样、更难预测，提供了更强的训练信号。这一结论有条件：在只用约8万句平行数据的模拟低资源设置中，反向模型本身较差，采样产生的噪声过大，束搜索反而更好。

Caswell等提出另一种解释\cite{nlp-gen-caswell2019tagged}：加噪之所以有效，主要是让模型能区分合成数据与真实数据。\term{带标记回译}（tagged back-translation）不加噪，只在每个合成源句前加一个保留标记\code{<BT>}，模型据此知道这对数据的源端不可靠；推断时输入不带标记。他们的实验中，带标记的束搜索回译与加噪回译效果相当甚至更好。

反向模型越好，合成数据越好；而正向模型变好后又能为反向模型合成数据。\term{迭代回译}（iterative back-translation）利用这一点交替改进两个方向\cite{nlp-gen-hoang2018iterative}，如算法~\ref{alg:nlp-generation-iterative-bt}所示。

\Needspace{12\baselineskip}
\begin{algorithm}[迭代回译]
  \label{alg:nlp-generation-iterative-bt}
  \begin{algorithmic}[1]
    \Require 平行数据$\mathcal D_{\mathrm{par}}$，源语单语$\mathcal M_x$，目标语单语$\mathcal M_y$，轮数$R$
    \Ensure 正向模型$p_{\bm\theta}(y\mid x)$与反向模型$q_{\bm\phi}(x\mid y)$
    \State 在$\mathcal D_{\mathrm{par}}$上分别训练$p_{\bm\theta}$与$q_{\bm\phi}$
    \For{$r=1,\ldots,R$}
      \State 用$q_{\bm\phi}$把$\mathcal M_y$译成$\{(\tilde x,y)\}$，在$\mathcal D_{\mathrm{par}}\cup\{(\tilde x,y)\}$上重训$p_{\bm\theta}$
      \State 用$p_{\bm\theta}$把$\mathcal M_x$译成$\{(x,\tilde y)\}$，在$\mathcal D_{\mathrm{par}}\cup\{(\tilde y,x)\}$上重训$q_{\bm\phi}$
    \EndFor
    \State \Return $p_{\bm\theta},q_{\bm\phi}$
  \end{algorithmic}
\end{algorithm}

回译的风险在合成源句的错误。若反向模型漏掉“不”，得到的训练对就把肯定的源句与否定的目标句配在一起；这类错误需要用质量估计分数筛选或降权来控制，筛选器不能使用测试集参考。回译也要与运行时的\term{枢轴翻译}（pivot translation）区分：后者在推断时经第三种语言两步翻译，例如把“不要修改12份记录”先译成英语再译成德语。回译的噪声通过训练数据影响参数，枢轴的错误则沿两步推断直接传播，第一步漏掉的否定，第二步没有依据恢复。

\subsubsection{多语模型与目标语言标记}
\label{sec:nlp-generation-multilingual-mt}

为每个语向训练一个模型，$n$种语言就需要$n(n-1)$个模型，低资源语向的模型也无法借用其他语向的数据。Johnson等让所有语向共享同一个编码器—解码器和同一份子词词表，只在源句前加一个\term{目标语言标记}（target language token），例如\code{<2es>}表示“译成西班牙语”\scite{nlp-gen-johnson2017multilingual}。模型结构、损失和解码都不变，训练数据是所有语向句对的混合。高低资源语向的数据量相差几个数量级，直接混合会让低资源语向几乎采样不到，通常用温度采样调节各语向的比例，公式与权衡见第~\ref{sec:nlp-llm-temperature-sampling}小节。

共享参数带来两种效果。一是迁移：编码器被迫把不同语言的句子映射到可共用的表示上，低资源语向得益于高资源语向学到的结构。Johnson等报告，多对一（多种源语言译成同一目标语言）的设置中多语模型普遍持平或优于单语向模型；一对多的设置收益较小，个别语向（如WMT英法）明显下降，他们推测原因是解码器要在一个共享词表中生成多种目标语言。二是\term{零样本翻译}（zero-shot translation）：只用葡萄牙语—英语与英语—西班牙语训练，推断时给葡萄牙语源句加\code{<2es>}，模型能直接输出西班牙语，尽管训练中从未见过葡萄牙语—西班牙语句对。在他们的实验中，零样本译文的质量低于先译成英语再译成西班牙语的枢轴翻译，但证明了共享表示确有跨语言的成分。

这一做法的问题也来自共享。所有语向分享固定容量，语言越多，每种语言分到的容量越少，高资源语向的质量可能下降；目标语言只由一个标记指定，零样本语向和低资源语向常出现\term{偏离目标语言}（off-target）的输出，即译成了英语或源语言本身。后续工作通过加大模型、加入语言特定的参数，以及增加非英语中心的平行数据来缓解，NLLB的模型就沿这一方向扩展到200种语言。

\subsubsection{无监督机器翻译}
\label{sec:nlp-generation-unsupervised-mt}

若两种语言之间完全没有平行数据，回译的反向模型无从训练。\term{无监督机器翻译}（unsupervised machine translation）只用两种语言各自的单语语料$\mathcal M_1,\mathcal M_2$，依靠三个要素\cite{nlp-gen-lample2018unsupervised,nlp-gen-artetxe2018unsupervised,nlp-gen-lample2018phrase}：一个让两种语言起步就处在共同表示空间中的初始化；一个让模型学会生成每种语言的去噪目标；一个把单语数据变成伪平行数据的在线回译。

\paragraph{跨语言初始化。}一种做法是先学\term{跨语言词嵌入}（cross-lingual word embedding）。两种语言的单语词向量空间形状相似，可以用一个正交矩阵近似对齐。给定一份种子词典$\{(x_k,y_k)\}$，把对应词向量按列排成$\bm X,\bm Y\in\mathbb R^{d\times n}$，正交普鲁克问题
\begin{equation}
  \bm W^\ast=\argmin_{\bm W^{\top}\bm W=\bm I}\bigl\|\bm W\bm X-\bm Y\bigr\|_F
  =\bm U\bm V^{\top},\qquad \bm U\bm\Sigma\bm V^{\top}=\operatorname{SVD}(\bm Y\bm X^{\top})
  \label{eq:nlp-generation-procrustes}
\end{equation}
有闭式解。没有种子词典时，Conneau等先用对抗训练得到初始映射，再取映射后互为最近邻的词对作为词典，反复用式\eqref{eq:nlp-generation-procrustes}精化\scite{nlp-gen-conneau2018muse}。对齐后的词嵌入用来初始化共享编码器—解码器的嵌入层。另一种更强的做法是在两种语言的拼接语料上用共享子词词表做掩码语言模型预训练，用得到的跨语言模型初始化整个编码器与解码器\scite{nlp-gen-conneau2019xlm}。

\paragraph{去噪自编码与在线回译。}模型以语言标记区分输出语言。去噪目标让模型从加噪的句子恢复原句，
\begin{equation}
  \mathcal L_{\mathrm{dae}}=\sum_{\ell\in\{1,2\}}\mathbb E_{z\sim\mathcal M_\ell}\bigl[-\log p_{\bm\theta}(z\mid C(z),\ell)\bigr],
  \label{eq:nlp-generation-unsup-dae}
\end{equation}
其中加噪函数$C$以一定概率删去每个词，并在小窗口内（例如不超过3个位置）打乱词序。单纯复制输入无法完成这个任务，模型必须学会该语言的词序与搭配，这使解码器成为两种语言各自的语言模型。在线回译用当前模型把单语句子译到另一种语言，再要求模型从这个合成译文恢复原句，
\begin{equation}
  \mathcal L_{\mathrm{bt}}=\mathbb E_{x\sim\mathcal M_1}\bigl[-\log p_{\bm\theta}(x\mid\hat y(x),1)\bigr]
  +\mathbb E_{y\sim\mathcal M_2}\bigl[-\log p_{\bm\theta}(y\mid\hat x(y),2)\bigr],
  \label{eq:nlp-generation-unsup-bt}
\end{equation}
$\hat y(x)$由当前模型贪心或束搜索生成，不对生成过程求梯度。这正是式\eqref{eq:nlp-generation-back-translation}的回译，只是反向模型就是同一个模型、且每一步都在改进。算法~\ref{alg:nlp-generation-unsupervised-mt}给出训练循环。

\Needspace{12\baselineskip}
\begin{algorithm}[无监督机器翻译的训练循环]
  \label{alg:nlp-generation-unsupervised-mt}
  \begin{algorithmic}[1]
    \Require 单语语料$\mathcal M_1,\mathcal M_2$，共享子词词表，加噪函数$C$
    \Ensure 双向翻译模型$p_{\bm\theta}$
    \State 用跨语言词嵌入或跨语言预训练模型初始化$\bm\theta$
    \Repeat
      \State 从$\mathcal M_1,\mathcal M_2$各取一批句子，按式\eqref{eq:nlp-generation-unsup-dae}计算去噪损失
      \State 用当前$\bm\theta$把两批句子分别译到另一种语言，得到伪平行对（不回传梯度）
      \State 按式\eqref{eq:nlp-generation-unsup-bt}计算回译损失，对两项损失之和更新$\bm\theta$
    \Until{开发集上的往返重构或无监督选择准则不再改善}
    \State \Return $p_{\bm\theta}$
  \end{algorithmic}
\end{algorithm}

三个要素缺一不可：没有好的初始化，第一轮回译的输出与原句几乎无关，回译损失学不到对应；没有去噪目标，模型只靠回译容易学会把输入原样复制或输出退化句子；没有回译，去噪只训练了单语重构，从未要求跨语言生成。Lample等的第一版方法还在编码器输出上加了一个判别器，用对抗损失让两种语言的编码难以区分，同一作者的后续版本去掉了这一项，效果并未变差。

无监督翻译的效果高度依赖两种语言的相近程度和单语语料的可比性。报告较好结果的多是英法、英德这类语言相近、新闻领域语料充足的语对；当语言差异大、单语数据少或领域不匹配时，效果明显下降，甚至不如用少量平行数据训练的有监督模型\scite{nlp-gen-kim2020unsupervised}。因此在实际低资源场景中，它更多作为初始化或与回译、多语模型结合使用。

\subsection{大语言模型翻译}
\label{sec:nlp-generation-llm-translation}

大语言模型在多语文本上预训练后，按指令就能翻译。与专用的编码器—解码器翻译模型相比，它的优势在于能读入长上下文和自然语言要求（术语表、语体、目标读者），对高资源语向和需要领域知识的文本表现突出；不足在于预训练数据以英语为主，低资源语向较弱，也容易出现漏译、添加解释性内容、译成错误语言等问题。WMT24的通用翻译评测中，表现最好的系统多基于大语言模型，但评测报告的标题本身就强调了翻译问题并未解决\scite{nlp-intro-kocmi2024}。下面先讲如何形成翻译能力，再讲如何从多个候选中选出译文。

\subsubsection{提示翻译与翻译微调}
\label{sec:nlp-generation-llm-mt-prompt-ft}

最直接的用法是提示。零样本提示只给指令和源句，例如“将下面的中文译成英文。中文：$x$ 英文：”；少样本提示在前面附上$k$个示例句对，模型按
\begin{equation}
  p_{\bm\theta}\bigl(y\mid I,(x^{(1)},y^{(1)}),\ldots,(x^{(k)},y^{(k)}),x\bigr)
  \label{eq:nlp-generation-llm-mt-prompt}
\end{equation}
生成，$I$为指令。示例的作用是指定语向、格式与语体，也可以示范术语的译法。示例选择的一般方法（按与输入的相似度检索、兼顾多样性）见第~\ref{sec:nlp-inference-icl-shots}小节。在翻译上，Vilar等对PaLM的系统比较发现，示例本身的翻译质量是影响最大的因素，与源句语义相似的示例带来的额外收益并不稳定\scite{nlp-gen-vilar2023palm}。因此示例池应先保证质量，再考虑相似度；术语要求直接写进指令（例如“海棠楼译为Haitang Building”）往往比指望示例中恰好出现更可靠。

提示不改变参数，模型在低资源语言上的能力上限仍由预训练决定。用大量平行数据直接微调大模型，又可能冲淡它在预训练中获得的通用能力。Xu等的ALMA把翻译微调分成两个阶段\scite{nlp-gen-xu2024alma}。第一阶段在所涉及语言的单语文本上继续做下一词元预测，
\begin{equation}
  \mathcal L_1=-\sum_{z\in\mathcal M}\sum_{t}\log p_{\bm\theta}(z_t\mid z_{<t}),
  \label{eq:nlp-generation-alma-stage1}
\end{equation}
补强模型对非英语语言的掌握；第二阶段只用少量（数万句对量级）高质量的人工翻译，把句对写成翻译提示，只对译文部分计算损失，
\begin{equation}
  \mathcal L_2=-\sum_{(x,y)}\sum_{t=1}^{|y|}\log p_{\bm\theta}\bigl(y_t\mid \operatorname{prompt}(x),y_{<t}\bigr).
  \label{eq:nlp-generation-alma-stage2}
\end{equation}
第二阶段的损失掩码与第~\ref{sec:nlp-sft-masked-loss}小节的有监督微调相同。ALMA的核心判断是：大模型的翻译能力主要受限于对目标语言本身的掌握，而不是缺少海量平行数据；用单语数据补语言能力、用少而精的平行数据教会翻译格式，比用大量质量参差的平行数据微调更有效。原文报告，基于7B与13B参数LLaMA-2的模型在其评测语向上超过了此前更大的专用翻译模型。

与专用神经翻译相比，提示翻译不需要训练，但每次推断都要处理更长的输入，且难以覆盖低资源语向；翻译微调的大模型质量高，但推断成本远高于几亿参数的编码器—解码器模型，对语向覆盖也仍受预训练语料限制。实际系统常把二者结合：用专用模型处理大批量的常规文本，用大模型处理需要上下文与指令的文本。

\subsubsection{质量估计重排与MBR解码}
\label{sec:nlp-generation-qe-mbr}

无论用哪种模型，按模型概率选出的译文未必是质量最好的译文。第~\ref{sec:nlp-inference-reranking}小节的重排与第~\ref{sec:nlp-inference-mbr}小节的最小贝叶斯风险解码提供了两种从候选中选择的方式，翻译是它们最成熟的应用，因为这里有与人工判断较一致的学习型度量（COMET等，见第~\ref{sec:nlp-generation-comet}小节）。

先生成$N$个候选$\mathcal C=\{y_1,\ldots,y_N\}$，可以是束搜索的$N$-best，也可以是采样结果（采样更多样，通常更适合后续选择）。\term{质量估计}（quality estimation, QE）模型只读源句与候选、不需要参考译文，直接预测人工质量分，重排取
\begin{equation}
  \hat y_{\mathrm{QE}}=\argmax_{y\in\mathcal C}\operatorname{QE}(x,y).
  \label{eq:nlp-generation-qe-rerank}
\end{equation}
MBR解码则把其余候选当作伪参考，以COMET为效用，取与其他候选平均效用最高的那个：
\begin{equation}
  \hat y_{\mathrm{MBR}}=\argmax_{y_i\in\mathcal C}\frac{1}{N}\sum_{j=1}^{N}\operatorname{COMET}(x,y_i,y_j),
  \label{eq:nlp-generation-comet-mbr}
\end{equation}
这是式\eqref{eq:nlp-inference-mbr-sample}在效用$u(y_i,y_j)=\operatorname{COMET}(x,y_i,y_j)$下的特例。直观上，QE重排相信一个外部评审对单个候选的打分；MBR选的是候选集合中的“共识”译文，与多数候选在意思上最接近，因此不容易选到偶然出错的离群候选。Freitag等的人工评价显示，以神经度量为效用的MBR优于束搜索\cite{nlp-inference-freitag2022mbr}；Fernandes等系统比较了重排、MBR及二者的组合，其中一种两阶段做法先用QE把候选筛到前$k$个，再在其中做MBR，把$N^2$次效用计算降为$k^2$次\scite{nlp-gen-fernandes2022quality}。

两种方法的代价不同。QE重排需要$N$次度量计算；MBR需要$N^2$次，$N=64$时约4000次。对按句分别编码源句、候选和参考的COMET模型，编码器的前向只需对$N$个候选各做一次并缓存句向量，$N^2$次计算只落在其后的小型前馈网络上，实际开销远低于$N^2$次完整前向。更值得注意的是度量偏好的放大：选择本身是在优化某个度量，Fernandes等发现，用于选择的度量上的提升明显大于其他度量上的提升。评价这类系统时，应换用未参与解码的度量，并以人工评价为准；评价方法的讨论见第~\ref{sec:nlp-generation-evaluation}节与第~\ref{chap:nlp-evaluation}章。

\subsection{同声传译}
\label{sec:nlp-generation-simultaneous}

\term{同声传译}（simultaneous translation）要求在源句尚未读完时就开始输出。模型每一步在两个动作之间选择：READ读入一个源词，WRITE写出一个目标词。记$g(t)$为写出第$t$个目标词时已读入的源词数，读写策略就由$g$描述。挑战在于：只能看到源前缀；等待越久质量越好、延迟越高；语序差异大时（如日语、德语动词在句末），模型可能被迫在关键信息到达之前作出承诺。方法谱系从固定策略到自适应策略：wait-$k$按固定节奏读写，简单且延迟可预期；单调注意力类方法让模型按内容决定何时读、何时写。评价必须同时报告质量与延迟。声音输入输出时的语音识别与合成环节见第~\ref{chap:nlp-speech-tasks}章。

\subsubsection{STACL与wait-k前缀翻译}
\label{sec:nlp-generation-wait-k}

STACL的wait-$k$策略先读$k$个源词，此后每写一个词读一个词，源句读完后连续写到结束\cite{nlp-intro-ma2019stacl}：
\begin{equation}
  g(t)=\min(k+t-1,S),\qquad
  p(y\mid x)=\prod_{t=1}^{T}p(y_t\mid x_{1:g(t)},y_{<t}).
  \label{eq:nlp-generation-wait-k}
\end{equation}
关键是训练也服从同样的可见性限制，即\term{前缀到前缀}（prefix-to-prefix）训练：模型学的是从源前缀预测目标前缀，因此会学到一定的预测能力，例如在看到动词之前，根据主语和宾语猜测可能的动作。若用整句训练的普通模型只在测试时按wait-$k$读写，模型从未见过这种残缺输入，效果明显更差。

在Transformer中，前缀到前缀训练用掩码实现，仍然可以一次算完所有目标位置。解码器交叉注意力的掩码为
\begin{equation}
  (\bm M_k)_{t,i}=\begin{cases}0,&i\le g(t),\\-\infty,&i>g(t),\end{cases}
  \label{eq:nlp-generation-wait-k-mask}
\end{equation}
它替代式\eqref{eq:nlp-generation-transformer-decoder}第二行中的$\bm M_{\mathrm{src}}$。编码器有两种选择。若用双向自注意力，源位置$i$的表示会读到$i$之后的词，第$t$个目标位置应看到的是“只编码了前$g(t)$个词”的表示，不同$t$需要不同的编码，训练代价成倍增加。若把编码器自注意力也改为因果掩码（单向编码），每个源位置的表示只依赖它之前的词，新词到来时旧表示不变，训练时一次编码即可满足所有$t$的可见性，推断时也能增量计算；代价是每个位置失去右侧上下文。STACL原文的训练实现就采用这种单向编码器。无论哪种实现，都不能让编码器读到后缀而只在交叉注意力上遮挡，否则训练时的信息会多于推断。

延迟需要一个能比较不同策略的指标。\term{平均滞后}（average lagging, AL）把策略与一个“理想的同步译员”比较：后者按长度比$r=T/S$匀速前进，写第$t$个词时已读$(t-1)/r$个源词。AL对源句读完之前的各步平均滞后量，
\begin{equation}
  \mathrm{AL}=\frac{1}{\tau}\sum_{t=1}^{\tau}\Bigl[g(t)-\frac{t-1}{r}\Bigr],\qquad
  \tau=\min\{t:g(t)=S\}.
  \label{eq:nlp-generation-average-lagging}
\end{equation}
只算到$\tau$为止，是因为源句读完后的输出已不存在等待，计入它们会让长译文显得延迟更低。当$r=1$时，wait-$k$的每一项都是$k$，$\mathrm{AL}=k$，指标的单位就是“落后多少个源词”。

用日语教学源句“彼は／会議に／参加／し／ない”说明。五个源词块依次表示“他、会议、参加、做、否定”，完整句意是“He will not attend the meeting”，$T=6$，$r=6/5$。设$k=2$，每步只保留一个候选：
\begin{center}
  \begin{tabular}{@{}cllcl@{}}
    \toprule
    目标步$t$ & 新READ内容 & 已读$g(t)$ & $g(t)-\frac{t-1}{r}$ & WRITE\\
    \midrule
    1 & 彼は、会議に & 2 & 2.000 & He\\
    2 & 参加 & 3 & 2.167 & will\\
    3 & し & 4 & 2.333 & attend\\
    4 & ない & 5 & 2.500 & the\\
    5 & 无 & 5 & — & meeting\\
    6 & 无 & 5 & — & EOS\\
    \bottomrule
  \end{tabular}
\end{center}
$\tau=4$，$\mathrm{AL}=(2+2.167+2.333+2.5)/4=2.25$。第3步写出attend时还没有看到否定；第4步虽然读到ない，在只能追加词元的接口中已无法把not插回will与attend之间，这条轨迹漏译了否定。改为$k=5$时，第一次写出之前已读完整句，$g(t)\equiv5$，$\tau=1$，$\mathrm{AL}=5$，可以完整译出“He will not attend the meeting”。两种设置分别代表低延迟与高质量的两端：句末否定这类语序差异，正是固定策略难以兼顾的情形。真实系统的延迟还取决于模型计算时间和输入粒度，报告时不能只给$k$。

\subsubsection{自适应读写策略}
\label{sec:nlp-generation-adaptive-policy}

wait-$k$对所有句子用同样的节奏：语序一致的句子它等得太久，动词在句末的句子它又写得太早。自适应策略让模型根据已读内容决定READ还是WRITE。直接把读写决策当作离散动作，可以用强化学习训练，但奖励稀疏、方差大。单调注意力类方法把决策嵌入注意力，使它可以用可微的期望形式训练。

\term{单调注意力}（monotonic attention）的推断过程是：生成第$t$个目标词时，从上一步停下的源位置开始向右扫描，在每个位置$i$计算停止概率$p_{t,i}=\sigma\bigl(\mathrm{Energy}(\bm s_{t-1},\bm h_i)\bigr)$；若停下，就只注意$\bm h_i$并写出$y_t$，否则READ下一个源词\scite{nlp-gen-raffel2017monotonic}。停止位置随$t$单调不减，这正是同传所需的“只向前读”。离散的停止决策不能求导，训练时改用停止位置的期望分布，它满足递推
\begin{equation}
  \alpha_{t,i}=p_{t,i}\sum_{k=1}^{i}\Bigl(\alpha_{t-1,k}\prod_{l=k}^{i-1}(1-p_{t,l})\Bigr),
  \label{eq:nlp-generation-monotonic-attention}
\end{equation}
含义是：上一步停在$k$的概率，乘以本步从$k$起一路不停直到$i$才停的概率，对所有$k\le i$求和。这个递推可以用累积乘积并行计算，复杂度为每句$O(ST)$。

只注意停止位置一个源词，信息太少，质量明显低于普通注意力。Arivazhagan等的\term{单调无限回看注意力}（monotonic infinite lookback attention, MILk）保留单调的停止机制决定“读到哪里”，停下后对已读的全部源词$\bm h_1,\ldots,\bm h_{g(t)}$做普通软注意力\scite{nlp-gen-arivazhagan2019milk}。训练时把式\eqref{eq:nlp-generation-monotonic-attention}的停止分布与软注意力组合，得到期望注意力权重；期望延迟$\mathbb E[g(t)]=\sum_i i\,\alpha_{t,i}$也可微，于是在似然之外加入延迟损失，
\begin{equation}
  \mathcal L=-\log p(y\mid x)+\lambda\,\mathcal C\bigl(\mathbb E[g(1)],\ldots,\mathbb E[g(T)]\bigr),
  \label{eq:nlp-generation-milk-loss}
\end{equation}
$\mathcal C$取平均滞后的可微版本，$\lambda$控制质量与延迟的折中。$\lambda=0$时模型倾向于读完整句再写，$\lambda$越大，模型越早开始写。

Ma等的\term{单调多头注意力}（monotonic multihead attention, MMA）把这一机制搬到Transformer\cite{nlp-gen-ma2020mma}：解码器每层交叉注意力的每个头各自做单调注意力（只看停止位置，或像MILk那样无限回看），各有自己的停止位置。只有当所有头都决定停止时才写出下一个词，因此实际的$g(t)$由读得最远的头决定。为了不让个别头一直读到句末、拖高整体延迟，训练损失中除平均延迟项外，还加入各头期望延迟的方差项，约束各头的读取进度相近。

与wait-$k$相比，自适应策略在相同平均延迟下通常质量更高，因为它在语序一致处少等、在需要等待的地方多等；代价是训练要计算期望注意力的递推、要调延迟权重，且单句的延迟随内容波动，不如wait-$k$可预期。对需要保证最长等待的应用，有时仍选用固定策略。

\subsection{篇章翻译}
\label{sec:nlp-generation-document-translation}

前面的方法都逐句翻译。逐句看似正确的译文放在一起，仍可能指代不明、主语缺失或术语前后不一致：前句的“海棠研究所”在后句写成“该所”，英文要决定用the institute还是it；中文省略的主语，英文必须补出；同一个术语在一篇文档中应保持同一个译名。这些都需要当前句之外的信息。解决思路有两条：翻译时把上下文作为额外条件；或者先逐句翻译，再只在目标语端修订篇章一致性。评价需要专门的对比测试集，因为这类错误在整篇中只占少数词元，语料级BLEU几乎不受影响。

\subsubsection{上下文拼接与多编码器翻译}
\label{sec:nlp-generation-context-translation}

最简单的做法不改模型，只改数据。Tiedemann与Scherrer的\emph{2-to-2}拼接把前一句与当前句用分隔标记连起来作为输入，目标端同样拼接前一句与当前句的译文\cite{nlp-gen-tiedemann2017context}：
\[
  x_{k-1}\ \langle\mathrm{BRK}\rangle\ x_k\ \longrightarrow\ y_{k-1}\ \langle\mathrm{BRK}\rangle\ y_k .
\]
推断时只保留最后一个分隔标记之后的部分作为当前句译文。另一种\emph{2-to-1}设置只在源端加上下文、目标端只输出$y_k$。拼接法的好处是任何句级模型都能直接用；代价是输入长度加倍，上下文与当前句在注意力中地位相同，模型需要自己学会区分。

多编码器方法为上下文单独编码，再与当前句的表示融合。Voita等的上下文感知Transformer让上下文编码器与源句编码器共享前$L-1$层，在源句编码器的最后一层中，每个位置除自注意力外再对上下文编码做一次注意力，二者用门控融合\cite{nlp-gen-voita2018anaphora}：
\begin{equation}
  \bm g_i=\sigma\bigl(\bm W_g[\bm c^{(s)}_i;\bm c^{(c)}_i]+\bm b_g\bigr),\qquad
  \bm c_i=\bm g_i\odot\bm c^{(s)}_i+(1-\bm g_i)\odot\bm c^{(c)}_i,
  \label{eq:nlp-generation-context-gate}
\end{equation}
其中$\bm c^{(s)}_i$是源句自注意力的输出，$\bm c^{(c)}_i$是对上下文注意力的输出，门$\bm g_i\in(0,1)^d$逐维决定当前位置取多少上下文信息。他们分析注意力权重发现，模型学会了让代词更多地注意上下文中的先行词，即在没有显式监督的情况下学到了部分指代消解。与拼接法相比，多编码器把“当前句”和“上下文”分开，门控值可以检查；代价是要修改结构，且通常只用前一两句上下文。

怎样知道模型真的用了上下文？\term{对比测试集}（contrastive test set）为每个样本准备上下文、源句、正确译文$y^{+}$和一个只在篇章现象上出错的对比译文$y^{-}$（例如代词性别错、省略的动词形式错），检查模型是否给正确译文更高的概率：
\begin{equation}
  \mathrm{Acc}=\frac1N\sum_{n=1}^{N}\mathbb 1\Bigl[\log p\bigl(y^{+}_n\mid x_n,\mathrm{ctx}_n\bigr)>\log p\bigl(y^{-}_n\mid x_n,\mathrm{ctx}_n\bigr)\Bigr].
  \label{eq:nlp-generation-contrastive-accuracy}
\end{equation}
它评的是打分而不是生成，因此能精确地隔离某一种现象。Bawden等为英法翻译构造了指代与衔接的对比集\cite{nlp-gen-bawden2018discourse}，Voita等为英俄翻译构造了指示、省略与词汇衔接的对比集\scite{nlp-gen-voita2019goodwrong}。这些集的正误译文只在篇章现象上不同，BLEU相近的模型在其上的准确率可以相差很大。

用两个教学篇章说明这些现象。第一篇为“海棠研究所发布了新通知。该所还说明了报名规则。”，机构译名已确认为Haitang Institute。句级系统只读第二句，可能写出“It also explained the registration rules”，it指代不明；读取前句后可写成“The institute also explained the registration rules”。第二篇为“林岚提交了申请。随后收到了确认。”，第二句省略的接收者应由前句恢复为林岚。只看第二句的译文“An acknowledgement was received”保留了事件，却丢了人物；上下文模型可以写成“Lin Lan then received an acknowledgement”。对比测试集正是把这类成对的译文放在一起比较概率。

\subsubsection{DocRepair目标语篇章修订}
\label{sec:nlp-generation-docrepair}

上下文模型需要篇章级平行数据训练，而这类数据远少于句级平行数据。Voita等的DocRepair换了一个角度：句级翻译的篇章错误主要表现为目标语篇章读起来不连贯，那么可以只在目标语端学一个修订模型，把不一致的逐句译文改成一致的篇章\scite{voita-etal-2019-context}。这个模型只需要目标语的单语文档。

训练数据这样构造：取目标语单语文档中连续的$K$句$Y=(y_1,\ldots,y_K)$（原文取$K=4$），把每句独立地往返翻译一次，
\begin{equation}
  \tilde y_k=\operatorname{MT}_{x\to y}\bigl(\operatorname{MT}_{y\to x}(y_k)\bigr),\qquad k=1,\ldots,K,
  \label{eq:nlp-generation-docrepair-roundtrip}
\end{equation}
逐句往返会丢掉跨句线索，人为制造出句级翻译特有的指代与衔接错误。修订模型$q_{\bm\psi}$是一个普通的编码器—解码器，输入为$\tilde Y=(\tilde y_1,\ldots,\tilde y_K)$拼接成的序列，输出为原文档$Y$，按条件最大似然$\mathcal L=-\log q_{\bm\psi}(Y\mid\tilde Y)$训练。推断时，先用句级模型逐句翻译源文档得到$\hat y_k=\operatorname{MT}_{x\to y}(x_k)$，再取
\begin{equation}
  \hat Y=\argmax_{Y}\,q_{\bm\psi}\bigl(Y\mid\hat y_1,\ldots,\hat y_K\bigr).
  \label{eq:nlp-generation-docrepair}
\end{equation}

教学输入为“Haitang Institute issued a new notice. Begonia Institute said he would update it.”第二句的机构名与代词都与前句不一致，修订候选为“Haitang Institute issued a new notice. The institute said it would update the notice.”

DocRepair的优点是只用单语数据、可以接在任何句级系统之后；局限也来自它看不到源文。若第二句原文确实指另一个机构，修订模型按“连贯”的标准统一名称就会改错；往返翻译制造的错误与真实系统的错误分布也不完全相同。因此它适合修正衔接与一致性，源意是否保留仍需对照源文另行核验。

\section{约束生成}
\label{sec:nlp-generation-constraints}

把改期通知译成英文时，学校的术语表规定“海棠楼”须译为Haitang Hall，“校外来宾”须译为
external guests；同一份通知还要写进日历系统，系统只接受带固定字段的JSON，日期字段必须是
通知中的新日期。前一种要求规定输出必须出现哪些短语，称为词汇约束；后一种要求规定输出的
格式与字段取值，称为结构约束。两者都是硬要求：违反一次，整条输出就不能交付。

满足硬要求有四个困难。第一，普通束搜索只按模型分数剪枝，术语的首词元若概率较低，含它的
前缀可能在头几步就被剪掉，生成结束后再筛选已经找不回来。第二，生成后再把术语强行替换或
插入，会破坏词形、语序与上下文的衔接。第三，结构合法不等于取值正确：一个能被解析的JSON
仍可能把报名截止日期填进活动日期字段。第四，约束之间、约束与长度预算之间可能互相冲突，
此时应当报告失败，而不是交付一个看似可用的近似结果。

本节按约束的类型组织方法。词汇约束沿着“生成后筛选或占位替换→在搜索状态中跟踪约束进度
的网格束搜索→固定总束宽、在进度组间分配名额的动态束分配”展开，并以训练时注入术语的软
方法作对照。结构约束的机制已在第~\ref{sec:nlp-inference-fsm}与第~\ref{sec:nlp-inference-grammar}小节
用有限状态机和下推自动机讲过，这里只补生成任务特有的取值绑定检查。正式、礼貌这类软属性
无法写成必含短语，改用属性分类器重加权，见第~\ref{sec:nlp-generation-fudge}小节。

\subsection{词汇约束解码}
\label{sec:nlp-generation-lexical-constraints}

词汇约束最常见的用途是术语翻译。机构名、产品名、法律与医学术语在同一份文件里必须使用
客户或术语库指定的译法，模型自己选出的近义译法即使正确，也会造成同一文件内的前后不一致。
单语改写中保留关键词、数据到文本中必须出现的实体名，也是同类要求。

把要求写成搜索问题。给定输入$x$与约束集合$\mathcal Z=\{z_1,\ldots,z_K\}$，每条约束
$z_i=(z_{i,1},\ldots,z_{i,\ell_i})$是长度为$\ell_i$的词元序列，总约束词元数为
$C=\sum_{i}\ell_i$。所求输出为
\begin{equation}
  \widehat y=\argmax_{y\in\mathcal Y_{\mathcal Z}}\log p_{\bm\theta}(y\mid x),
  \label{eq:nlp-generation-lexical-objective}
\end{equation}
其中可行集合$\mathcal Y_{\mathcal Z}$由所有以EOS结束、且每个$z_i$都作为连续子串出现的
输出组成。
它与第~\ref{sec:nlp-generation-beam}小节束搜索的目标相同，只是可行集合缩小了。困难在于
$\mathcal Y_{\mathcal Z}$不能逐词元判断：前缀“The meeting”还没有出现任何术语，却完全可能
延续成合格输出，因此不能像格式约束那样在每一步屏蔽非法词元，只能在搜索状态中记住每条
约束完成了多少。

一个搜索假设记为$h=(y_{<t},s,S,\pi)$：$y_{<t}$是已生成前缀，$s$是累计对数概率，
$S\subseteq\{1,\ldots,K\}$是已经完整出现的约束集合，$\pi=(i,j)$表示正处在第$i$条约束内部、
已匹配前$j$个词元，$\pi=\varnothing$表示没有未完成的短语，称为开放状态。假设的约束进度为
\[
  g(h)=\sum_{i\in S}\ell_i+j,
\]
开放状态下$j=0$。只有$g(h)=C$、处于开放状态的假设才允许生成EOS并交付。

在引入专门的搜索之前，有三种更简单的做法。生成后筛选先取$n$个最佳候选，再保留满足
约束的那些；当候选全部漏掉术语时它无从补救。占位替换先把源文术语换成特殊占位符，译完
再把占位符换成指定译法；它总能放进术语，却让模型看不到术语本身，在需要变格、变数或
调整冠词的语言里容易产生不合语法的句子。训练时注入术语则把指定译法直接拼在源文术语
旁边，并用额外的输入标记区分原文与注入词，让模型在训练中学会复制它们
\scite{dinu-etal-2019-training}
这种做法不改变解码器，也能让模型自行处理词形与衔接，但满足约束只是高概率行为，并没有
保证。下面两种解码方法给出硬保证：只要搜索在预算内找到可行路径，交付的输出一定包含全部
约束。

为在两种方法之间比较，固定一个教学模型。用两个双词元术语演示：$\mathtt{ab}$代表
Haitang Hall，$\mathtt{cd}$代表external guests，$\mathtt{o}$代表其他文字片段，故$K=2$、
$C=4$。下一词元的分布只依赖前缀结尾，按下表给出，空格表示概率为零：
\begin{center}
  \small
  \begin{tabular}{@{}lcccccc@{}}
    \toprule
    前缀结尾 & $\mathtt o$ & $\mathtt a$ & $\mathtt b$ & $\mathtt c$ & $\mathtt d$ & EOS\\
    \midrule
    起始 & 0.50 & 0.30 & & 0.15 & & 0.05\\
    $\mathtt o$ & 0.40 & 0.25 & & 0.25 & & 0.10\\
    $\mathtt a$ & 0.30 & & 0.60 & & & 0.10\\
    $\mathtt c$ & 0.10 & & & & 0.80 & 0.10\\
    $\mathtt{ab}$ & 0.30 & & & 0.50 & & 0.20\\
    $\mathtt{cd}$ & 0.10 & & & & & 0.90\\
    \bottomrule
  \end{tabular}
\end{center}
最多五步，EOS占一步，平局按$\mathtt a<\mathtt c<\mathtt o$打破。这些数值只用于演示搜索
状态，不是任何翻译模型的输出。

宽度为1的普通束搜索第一步选$\mathtt o$，此后每步都是$\mathtt o$比$\mathtt a$、$\mathtt c$
概率高，于是一路走成$\mathtt{oooo}\,\mathrm{EOS}$，两个术语都没有出现；生成后筛选只能丢弃
它，无法补回第一步就被剪掉的$\mathtt a$。

\subsubsection{Grid Beam Search}
\label{sec:nlp-generation-gbs}

普通束搜索让所有前缀争夺同一批名额。已经开始写术语的前缀往往含有低概率词元，累计分数
较低，会被只顾高分的普通前缀挤出束外。\term{网格束搜索}（Grid Beam Search, GBS）的思想
是按约束进度分组竞争：在生成步数$t$之外再加一维进度$g$，网格的每一格$G_{t,g}$独立保留
至多$k$个假设，只有进度相同的前缀互相比较\scite{hokamp-liu-2017-lexically}

格$G_{t,g}$的候选来自三种操作。\emph{生成}作用于开放状态的假设，由模型自由选择一个词元，
进度不变（即使选中的词元恰好是某个约束词元，也不计入进度），所以从$G_{t-1,g}$转移到$G_{t,g}$；\emph{开始}作用于开放状态的假设，追加某条
未完成约束$z_i$的首词元，进度加一，从$G_{t-1,g-1}$转移而来，若$\ell_i=1$该约束随即完成；
\emph{继续}作用于处在短语内部的假设，它只能追加下一个约束词元$z_{i,j+1}$，进度加一，
当$j+1=\ell_i$时把$i$并入$S$并回到开放状态。于是
\begin{equation}
  G_{t,g}=\operatorname{top}_k\Bigl(\mathrm{Gen}(G_{t-1,g})\cup
  \mathrm{Start}(G_{t-1,g-1})\cup\mathrm{Cont}(G_{t-1,g-1})\Bigr),
  \label{eq:nlp-generation-gbs}
\end{equation}
$\operatorname{top}_k$按累计对数概率取前$k$项。每个追加词元$v$都按模型原分布计分，
$s(hv)=s(h)+\log p_{\bm\theta}(v\mid x,y_{<t})$；约束只决定哪些扩展可以进入哪一格，
不改写模型概率。处于短语内部的假设不能执行生成操作，这保证了约束短语连续出现。
算法~\ref{alg:nlp-generation-gbs}写出完整过程。

\Needspace{16\baselineskip}
\begin{algorithm}[网格束搜索]
  \label{alg:nlp-generation-gbs}
  \begin{algorithmic}[1]
    \Require 模型$p_{\bm\theta}$，输入$x$，约束$z_1,\ldots,z_K$（总词元数$C$），每格宽度$k$，最大步数$T$
    \Ensure 满足全部约束的最高分输出，或失败状态
    \State $G_{0,0}\gets\{(\mathrm{BOS},0,\varnothing,\varnothing)\}$；完成集合$\mathcal F\gets\varnothing$
    \For{$t=1,\ldots,T$}
      \For{$g=0,\ldots,\min(t,C)$}
        \State $\mathcal A\gets$ 对$G_{t-1,g}$中开放假设取模型前$k$个扩展，不计入约束进度；EOS仅在$g=C$时允许
        \State $\mathcal A\gets\mathcal A\,\cup$ 对$G_{t-1,g-1}$中开放假设追加每条未完成约束的首词元
        \State $\mathcal A\gets\mathcal A\,\cup$ 对$G_{t-1,g-1}$中短语内部假设追加该短语的下一个词元
        \State $G_{t,g}\gets\mathcal A$中累计分数最高的$k$项，更新各自的$S$与$\pi$
      \EndFor
      \State 把$G_{t,C}$中以EOS结尾的假设移入$\mathcal F$
    \EndFor
    \If{$\mathcal F=\varnothing$}
      \State \Return 失败：预算内未找到满足全部约束的输出
    \EndIf
    \State \Return $\mathcal F$中累计分数最高的输出
  \end{algorithmic}
\end{algorithm}

在教学模型上取每格宽度$k=1$。图~\ref{fig:nlp-generation-gbs-grid}给出每格保留的假设。
第一步，$g=0$格保留$\mathtt o:0.50$；$g=1$格在开始操作得到的$\mathtt a:0.30$与
$\mathtt c:0.15$中保留$\mathtt a$，它的内部状态是“第一条约束已匹配一个词元，等待
$\mathtt b$”。第二步三格分别保留$\mathtt{oo}:0.20$、$\mathtt{oa}:0.125$（与$\mathtt{oc}$
同分，按平局规则取$\mathtt a$）和$\mathtt{ab}:0.18$；$\mathtt{ab}$完成了第一条短语，
但$g=2<C$，还不能接受EOS。沿右上对角线，
$\mathtt{ab}\to\mathtt{abc}\to\mathtt{abcd}\to\mathtt{abcd}\,\mathrm{EOS}$的累计概率依次为
$0.18,0.09,0.072,0.0648$，在第五步交付。$\mathtt{abc}$的进度为3，表示第二条短语仍缺
$\mathtt d$，而不是“两条约束都已完成”。其余格中的路径来不及在第五步前补齐术语和EOS。

\begin{figure*}[htbp]
  \centering
  % 第九卷 12.7.1.1：教学模型上的网格束搜索（每格宽度k=1，两条双词元约束ab、cd，C=4）。
% 格中数值为累计概率，由正文给出的教学分布逐步相乘得到；第5步只画出交付的假设。
\begin{tikzpicture}[
  x=0.9mm,y=1mm,figure picture,
  cell/.style={rectangle,rounded corners=1mm,draw=FigureStroke,fill=white,
    line width=.6pt,minimum width=19mm,minimum height=7mm,inner sep=1pt,
    font=\figurefont\scriptsize,align=center},
  best cell/.style={cell,draw=BookTeal,fill=FigureModelFill,line width=1pt},
  done cell/.style={cell,draw=BookTeal,fill=FigureInputFill,line width=1pt},
  gen/.style={draw=FigureMuted,line width=.5pt,-{Stealth[length=1.4mm,width=1mm]}},
  con/.style={draw=FigureMuted,line width=.5pt,dashed,-{Stealth[length=1.4mm,width=1mm]}},
  best/.style={draw=BookTeal,line width=1.2pt,-{Stealth[length=1.8mm,width=1.3mm]}},
  axis label/.style={font=\figurefont\scriptsize,text=FigureMuted}
]
  % 坐标：x=26*t，y=12*g
  \node[cell] (b) at (0,0) {BOS\\$1$};
  \node[cell] (o1) at (26,0) {$\mathtt{o}$\\$0.50$};
  \node[best cell] (a1) at (26,12) {$\mathtt{a}$\\$0.30$};
  \node[cell] (oo2) at (52,0) {$\mathtt{oo}$\\$0.20$};
  \node[cell] (oa2) at (52,12) {$\mathtt{oa}$\\$0.125$};
  \node[best cell] (ab2) at (52,24) {$\mathtt{ab}$\\$0.18$};
  \node[cell] (ooo3) at (78,0) {$\mathtt{ooo}$\\$0.08$};
  \node[cell] (ooa3) at (78,12) {$\mathtt{ooa}$\\$0.05$};
  \node[cell] (oab3) at (78,24) {$\mathtt{oab}$\\$0.075$};
  \node[best cell] (abc3) at (78,36) {$\mathtt{abc}$\\$0.09$};
  \node[cell] (oooo4) at (104,0) {$\mathtt{oooo}$\\$0.032$};
  \node[cell] (oooa4) at (104,12) {$\mathtt{oooa}$\\$0.02$};
  \node[cell] (ooab4) at (104,24) {$\mathtt{ooab}$\\$0.03$};
  \node[cell] (oabc4) at (104,36) {$\mathtt{oabc}$\\$0.0375$};
  \node[best cell] (abcd4) at (104,48) {$\mathtt{abcd}$\\$0.072$};
  \node[done cell] (eos5) at (130,48) {$\mathtt{abcd}\,\mathrm{EOS}$\\$0.0648$};

  % 生成操作（进度不变）
  \foreach \a/\b in {b/o1,o1/oo2,oo2/ooo3,ooo3/oooo4}{\draw[gen] (\a) -- (\b);}
  % 开始或继续约束（进度加一）
  \foreach \a/\b in {o1/oa2,oo2/ooa3,oa2/oab3,ooo3/oooa4,ooa3/ooab4,oab3/oabc4}{\draw[con] (\a) -- (\b);}
  % 交付路径
  \draw[best] (b) -- (a1);
  \draw[best] (a1) -- (ab2);
  \draw[best] (ab2) -- (abc3);
  \draw[best] (abc3) -- (abcd4);
  \draw[best] (abcd4) -- (eos5);

  % 坐标轴说明
  \foreach \t in {0,...,5}{\node[axis label] at ({26*\t},-8) {$t=\t$};}
  \foreach \g in {0,...,4}{\node[axis label,anchor=east] at (-12,{12*\g}) {$g=\g$};}
  \draw[gen] (60,58) -- (70,58) node[right,axis label,text=FigureInk] {生成};
  \draw[con] (84,58) -- (94,58) node[right,axis label,text=FigureInk] {开始或继续约束};
  \draw[best] (122,58) -- (132,58) node[right,axis label,text=FigureInk] {交付路径};
\end{tikzpicture}

  \caption{教学模型上的网格束搜索，每格宽度$k=1$。横轴是生成步$t$，纵轴是约束进度$g$
  （已匹配的约束词元数）；每格写出保留的前缀与累计概率。水平箭头是生成操作（进度不变），
  斜箭头是开始或继续约束（进度加一）。加粗路径$\mathtt{ab}\to\mathtt{abc}\to\mathtt{abcd}$在
  $g=C=4$的格中生成EOS后交付；宽度为1的普通束只会沿底行$g=0$前进，两条术语都不出现。第5步只画出交付的假设。}
  \label{fig:nlp-generation-gbs-grid}
\end{figure*}

例子同时显示了可行性的边界。两条不重叠的双词元短语至少需要四个正文词元和一个EOS，
把预算改成四步，任何路径都不可能完成；若再要求“必须含$\mathtt{ab}$”且“禁止出现
$\mathtt b$”，增加预算也无解。这两种情形下算法都返回失败，调用方应据此放宽预算或报告
约束冲突，而不是退回一条缺术语的输出。

网格共有$C+1$行，每步至多保存$(C+1)k$个假设，模型前向的批量也随之扩大，计算与存储
都按约束词元数线性增长。网格只是把束搜索的剪枝按进度分组，仍是近似搜索，找到的可行
输出不一定是式\eqref{eq:nlp-generation-lexical-objective}的全局最优。约束以固定词形出现，
需要变格的术语要么列出多种形式，要么交给训练时注入的方法；约束也只保证术语出现，
不保证出现的位置正确，external guests可能被放进“无需报名”的从句里，含义仍要另行核验。

\subsubsection{Dynamic Beam Allocation}
\label{sec:nlp-generation-dba}

网格束搜索的假设数随$C$增长，带来两个实际问题：约束多时解码变慢，而且同一批句子各自
的约束数不同，网格大小不一，难以组成整齐的批量。\term{动态束分配}（Dynamic Beam Allocation,
DBA）把总束宽固定为$k$，不论约束多少，每步都只对$k$个前缀调用模型；约束进度仍然参与
选择，但改为在一条束内部把名额分给不同进度的假设\scite{post-vilar-2018-fast}

每一步的候选来自三个来源：整个束的前$k$个高分扩展；每个前缀对每条未完成约束的约束扩展
（开始一条新约束，或继续当前短语）；每个前缀自身的最佳单一扩展，保证每个前缀至少有一个
后继参与竞争。候选按进度$g$分入$C+1$个组，原方法称之为银行（bank）。分配时先给每组
$\lfloor k/(C+1)\rfloor$个名额，整除剩下的名额给进度最高的组$C$；某组候选少于名额时，
把多出的名额从相邻组开始向外让给候选多于名额的组，组内按累计分数取优。原论文正文把每组
名额写作$\lfloor k/C\rfloor$，其参考实现按$C+1$组均分，这里按后者叙述。
与网格束搜索的另一处不同是允许“回退”：处在短语内部的假设若追加了非约束词元，它在该短语
上的部分进度清零，例如$\mathtt{ao}$不再算完成了$\mathtt{ab}$的一半，以后再出现$\mathtt a$
时重新开始匹配；已经完整出现的短语保留完成标记。算法~\ref{alg:nlp-generation-dba}写出一步
的更新。

\Needspace{14\baselineskip}
\begin{algorithm}[动态束分配的一步更新]
  \label{alg:nlp-generation-dba}
  \begin{algorithmic}[1]
    \Require 当前束$\mathcal B$（$|\mathcal B|\leq k$），约束$z_1,\ldots,z_K$（总词元数$C$）
    \Ensure 下一步的束$\mathcal B'$
    \State $\mathcal A\gets\mathcal B$中全部扩展里累计分数最高的$k$项
    \For{每个$h\in\mathcal B$}
      \State $\mathcal A\gets\mathcal A\,\cup$ $h$对每条未完成约束的开始或继续扩展
      \State $\mathcal A\gets\mathcal A\,\cup\{h\text{的最佳单一扩展}\}$
    \EndFor
    \State 更新$\mathcal A$中每个候选的$S$与$\pi$；打断短语的候选清零该短语的部分进度
    \State 按进度$g$把$\mathcal A$分为$\mathcal A_0,\ldots,\mathcal A_C$，每组分配$\lfloor k/(C+1)\rfloor$个名额，余数给$\mathcal A_C$
    \State 候选不足的组把空余名额由近及远转给仍有候选的组；各组按累计分数取优，合并为$\mathcal B'$
    \State 把$\mathcal B'$中进度为$C$且以EOS结尾的假设移入完成集合
  \end{algorithmic}
\end{algorithm}

沿用同一模型，取总宽度$k=2$。此时$k<C+1$，均分给每组零个名额，两个余数名额都归进度最高的空组，再由近及远
转出，结果偏向较高进度的组。为把取舍看得更清楚，本例改用一个教学分配规则：只有两组时各一项；有三组以上而只有两个名额时，给最高与最低进度组各一项。
这个规则只用来显示固定容量的取舍，不复现原方法的全部配额细节。第一步的候选分为
$\{\mathtt o:0.50\}$与$\{\mathtt a:0.30,\mathtt c:0.15\}$两组，保留$\mathtt o$与$\mathtt a$。
第二步候选包含$\mathtt{oo}:0.20$、$\mathtt{oa}:0.125$、$\mathtt{oc}:0.125$与
$\mathtt{ab}:0.18$，分属进度0、1、2三组；两个名额分给$\mathtt{oo}$与$\mathtt{ab}$，
网格束搜索在$g=1$格保留的$\mathtt{oa}$被丢掉。此后$\mathtt{ab}$仍能沿
$\mathtt{abc}\to\mathtt{abcd}\to\mathrm{EOS}$完成，在本例中结果与网格束搜索相同，但束内
已经少了一条备用路径；若后续分布不利于$\mathtt{ab}$，这条被丢掉的路径可能正是唯一可行解。

两种方法的关系可以这样理解：若取DBA的总宽度为$(C+1)k$，并且每组总有足够候选、每组恰好
分到$k$个名额，它的选择就接近每格宽度为$k$的网格束搜索，差别只在候选来源与回退规则。
DBA把“每组多少名额”从固定值变成随候选情况调整的量，换来每步$k$次模型前向、与$C$无关的
批量形状；构造约束扩展、更新匹配状态和分组仍有与约束数相关的开销，约为每步$O(kK)$个
额外候选。后续工作用前缀树表示多条约束的匹配状态，把这些操作向量化，使约束解码能在
GPU上与普通束搜索同样批量执行\scite{hu-etal-2019-improved}
总束宽必须至少为$C+1$，每个进度组才能分到名额；约束很多而束宽很小时，DBA会像本例一样
在中间进度组上失去候选。

表~\ref{tab:nlp-generation-lexical-methods}汇总词汇约束的几种做法。

\begin{table*}[htbp]
  \centering
  \small
  \renewcommand{\tabularxcolumn}[1]{>{\raggedright\arraybackslash}p{#1}}
  \begin{tabularx}{\linewidth}{@{}>{\raggedright\arraybackslash}p{22mm}XXX@{}}
    \toprule
    方法 & 解决的问题 & 核心机制 & 新代价\\
    \midrule
    生成后筛选 & 不改解码即可检查约束 & 取$n$个最佳候选，保留含全部约束者 & 候选全漏术语时无解\\
    占位替换 & 保证术语出现 & 源文术语换成占位符，译后替换 & 模型看不到术语，词形与衔接易错\\
    训练时注入 & 让模型自行处理词形 & 源文旁拼接指定译法并加标记，训练复制 & 只是高概率满足，无硬保证\\
    网格束搜索 & 术语前缀在剪枝中被挤掉 & 按约束进度分格，每格各留$k$项 & 假设数$(C+1)k$，随约束线性增长\\
    动态束分配 & 网格束的开销与批量问题 & 固定总束宽，按进度组分配名额，允许回退 & 束宽小于$C+1$时丢失中间进度路径\\
    \bottomrule
  \end{tabularx}
  \caption{词汇约束的处理方法。后两种在搜索找到可行路径时保证约束全部出现；所有方法都只
  保证术语在输出中出现，术语的位置与含义仍需另行核验。}
  \label{tab:nlp-generation-lexical-methods}
\end{table*}

\subsection{结构约束与值绑定检查}
\label{sec:nlp-generation-grammar}

日历系统要求的输出形如代码~\ref{code:nlp-generation-calendar-json}。

\begin{codeblock}[label={code:nlp-generation-calendar-json}]{日历系统接收的教学输出}
{"event": "ML-0418", "date": "2025-04-18",
 "weekday": "Fri", "time": "14:00", "room": "201"}
\end{codeblock}

第~\ref{chap:nlp-model-inference}章已经给出保证这类输出可解析的机制：把接口模式编译为
有限状态机或下推自动机，每一步只允许能延续成合法输出的词元，代价是局部重新归一化带来的
分布偏移和计算掩码的时间。这里不重复该机制，只讨论生成任务在结构合法之后还要检查的部分。

对一个结构化输出，要依次确认四件事。\emph{模式合法}：能被解析，必需字段齐全，类型正确；
这一层可由约束解码在构造时保证。\emph{值域合法}：日期是真实存在的日历日期，时间在
00:00至23:59之间，房间号属于允许的格式或枚举。\emph{事实绑定}：每个字段的值来自输入中
角色正确的那一项，或经过允许的规范化得到。\emph{跨字段一致}：星期与日期相符，报名截止
早于活动开始。后两层是生成任务特有的，第6章的约束解码不负责它们。

把事实绑定写成可以检查的条件。设从输入中抽取出带角色的事实集合
$R=\{(\rho,v)\}$，本例为
\[
  \begin{aligned}
  R=\{&(\text{新日期},\text{4月18日}),(\text{原日期},\text{4月17日}),
  (\text{截止},\text{4月16日18:00}),\\
  &(\text{时间},\text{14:00}),(\text{地点},\text{海棠楼201})\}.
  \end{aligned}
\]
接口为每个输出字段$f$规定它应绑定的角色$\rho(f)$和规范化函数$N_f$，例如日期字段绑定
“新日期”，$N_{\mathrm{date}}$按发布年份2025把“4月18日”写成\code{2025-04-18}；房间字段
绑定“地点”，$N_{\mathrm{room}}$只保留房间号。输出$o$通过事实绑定，当且仅当
\begin{equation}
  \forall f:\quad \exists\,(\rho(f),v)\in R,\quad N_f(o_f)=N_f(v).
  \label{eq:nlp-generation-binding}
\end{equation}
一个更弱的条件是“取值出现在输入中”，即$o_f\in\{N_f(v):(\rho,v)\in R\}$，不要求角色匹配。
弱条件能挡住凭空编造的日期，却挡不住角色互换：把报名截止日期\code{2025-04-16}填进
日期字段，值确实来自通知，绑定的角色却错了。

哪些检查可以前移到解码中？弱条件可以：为每个输入单独构造一份文法，让日期字段只能取
$\{\code{2025-04-16},\code{2025-04-17},\code{2025-04-18}\}$之一，第~\ref{sec:nlp-inference-grammar}小节
的机制就能保证输出不出现输入以外的日期。角色绑定通常不能，因为它需要知道哪个日期是
新日期，这一信息来自对输入的理解，可由信息抽取得到带角色的记录
（第~\ref{chap:nlp-information-extraction}章），再在生成后按式\eqref{eq:nlp-generation-binding}比对。
跨字段一致需要日历计算，也放在生成后检查：2025年4月18日是星期五，与\code{"Fri"}相符。
算法~\ref{alg:nlp-generation-binding-check}把这四层组织成一次核验。

\Needspace{14\baselineskip}
\begin{algorithm}[结构化输出的分层核验]
  \label{alg:nlp-generation-binding-check}
  \begin{algorithmic}[1]
    \Require 输出文本$y$，接口模式，带角色的输入事实$R$，字段绑定$\rho(f)$与规范化$N_f$，跨字段条件$\phi_1,\ldots,\phi_M$
    \Ensure 通过，或失败原因列表
    \State 解析$y$得到对象$o$；失败则返回“结构不合法”及位置
    \For{每个字段$f$}
      \State 检查$o_f$的类型与值域；不合法则记录“值域错误：$f$”
      \If{$R$中没有角色$\rho(f)$}
        \State 模式允许空值且$o_f$为空则通过，否则记录“输入缺少$\rho(f)$”
      \ElsIf{$N_f(o_f)$不等于任何$(\rho(f),v)\in R$的$N_f(v)$}
        \State 记录“绑定错误：$f$”，并附上$o_f$实际匹配到的角色（若有）
      \EndIf
    \EndFor
    \State 对每个$\phi_m$检查$\phi_m(o)$；不成立则记录“跨字段冲突：$\phi_m$”
    \State \Return 原因列表为空则通过，否则返回列表
  \end{algorithmic}
\end{algorithm}

逐层检查代码~\ref{code:nlp-generation-calendar-json}：模式合法；日期、时间与星期都在值域内；日期绑定“新日期”，房间绑定
“地点”且规范化后为201，事实绑定通过；星期与日期一致。若模型输出\code{"date":"2025-04-17"}，
前两层仍然通过，第三层报告日期实际匹配到“原日期”角色，这正是改期通知最容易出的错。
若接口还要求主讲人字段，而通知没有给出，模式允许空值时应输出空值；不允许时应返回
“输入缺少主讲人”，不能补一个看似合理的名字。约束解码在预算内无法闭合结构、或文法与
长度限制冲突时，同样返回失败及原因。

值绑定检查在多个生成任务中复用。数据到文本要求句中数值能回溯到记录；文本到SQL的
条件值须来自问题或数据库，见第~\ref{sec:nlp-analysis-text-to-sql}小节；工具调用的参数要
绑定到用户请求中的实体，见第~\ref{chap:nlp-solving-and-execution}章。这类检查的可靠性取决于
带角色的输入事实是否抽取正确，抽取错误会让一个正确输出被判为失败，或让错误输出通过，
因此抽取结果本身也要可追溯到原文位置。

\section{生成结果的评价}
\label{sec:nlp-generation-evaluation}

上一节的约束可以逐条判定对错，大部分生成质量却没有这样的判定程序。评价生成结果有五个
挑战。第一，一个输入有许多合法输出，参考文本只覆盖其中少数，合法改写可能因为与参考不同
而得低分。第二，流畅与忠实可以分离，一句通顺的译文可能改错了日期。第三，编辑任务的
输出与来源高度重合，只比较输出与参考，会奖励什么都不改。第四，自动指标与人工判断的
一致程度有限，而且随语言、领域和系统质量变化。第五，人工评价成本高，结果还依赖错误
类别的定义与评价者的一致性。

本节的方法谱系按“测量什么、需要什么输入”展开：先用来源检查忠实性；再看只需参考的
重合指标；编辑任务另用同时读取来源、输出与参考的编辑指标；为容忍合法改写，学习型指标
用预训练模型比较语义，无参考质量估计进一步去掉参考；人工评价是这些指标的最终参照，
元评价衡量指标与人工判断的一致程度；最后讨论用大语言模型按准则评判。评测协议、数据
独立性与统计比较已在第~\ref{chap:nlp-evaluation}章给出，这里只写各指标的定义、手算例子
与失效情形，节末表~\ref{tab:nlp-generation-metrics}列出它们的测量对象与输入。

\subsection{忠实性与无依据内容检测}
\label{sec:nlp-generation-faithfulness}

检查\term{忠实性}（faithfulness）之前，先要确定忠实于什么。翻译以源文原意为依据，
摘要以选定来源及用途为依据，编辑还要检查每处修改是否必要；开放故事中的虚构人物可以说
虚构的话，但引用的真实活动安排仍不能改成另一天。生成文本中没有来源支持的内容称为
幻觉（hallucination）。按与来源的关系可以分为两类\scite{nlp-generation-maynez2020faithfulness}
\term{内在幻觉}（intrinsic hallucination）错误地组合或改写了来源中的信息，与来源矛盾，
例如“活动仍于4月17日举行”，或把报名截止日期说成活动日期；\term{外在幻觉}（extrinsic
hallucination）加入了来源中无法核实的信息，例如“预计50人到场”。外在幻觉的内容在现实中
可能恰好为真，例如“海棠楼位于校园东区”，但它不能从来源得到支持，评价忠实性时仍记为
无依据；它是否可以接受，要看任务是否允许引入外部知识。两类错误的修订方式也不同：前者
要按来源改正，后者要删去或补上出处。

最直接的检查是规则：把实体、日期和数量规范化后逐项与来源比对。规则可解释，但只能覆盖
预先列出的类型，而且必须保留角色，否则把报名时间与开会时间互换，数值集合仍然一致；这与
上一节的值绑定是同一个问题。为覆盖任意陈述，学习式核验把问题化归为已有的语言理解任务，
沿两条路线发展：一条用自然语言推断判断来源能否推出生成文本，先有在合成错误样本上训练的
文档级分类器FactCC\scite{kryscinski-etal-2020-evaluating}，后有下面的句级聚合方法SummaC；
另一条从生成文本中提问，再用来源回答并比对，先有QAGS\scite{wang-etal-2020-asking}，后有
组件经过系统比较的QAFactEval。把生成文本进一步拆成原子事实逐条核验的方法见
第~\ref{chap:nlp-question-answering-dialogue}章。

\subsubsection{SummaC的句对推断核验}
\label{sec:nlp-generation-summac}

推断模型通常在单句对上训练，把整篇来源和整篇摘要一次性交给它，粒度与长度都超出训练
分布。SummaC先分句，为每个摘要句寻找来源中的支持，再汇总成文档分数
\scite{laban-etal-2022-summac}
设来源句为$x_1,\ldots,x_M$，摘要句为$y_1,\ldots,y_N$，推断模型（第~\ref{sec:nlp-classification-nli}小节）
给出蕴含概率$E_{i,j}=p(\text{蕴含}\mid x_i,y_j)$，构成$M\times N$的矩阵。零样本版本
SummaC-ZS对每个摘要句取最佳支持，再对摘要句平均：
\begin{equation}
  s_{\mathrm{ZS}}=\frac{1}{N}\sum_{j=1}^{N}\max_{1\leq i\leq M}E_{i,j}.
  \label{eq:nlp-generation-summac}
\end{equation}
取最大值意味着一个摘要句只要被某一来源句支持即可，平均则让每个摘要句同等计入。这一步
没有可学习参数。最大值只看一个数，丢掉了整列的分布：一列中有一个0.6，和一列中有多个
0.5，含义并不相同。卷积版本SummaC-Conv把第$j$列的$M$个蕴含分数离散成$H$个区间的直方图
$\bm u_j\in\mathbb R^H$，用一个可学习的一维卷积把直方图映射为该句的一致性分数，再平均：
\begin{equation}
  s_{\mathrm{Conv}}=\frac{1}{N}\sum_{j=1}^{N}g_{\bm\psi}(\bm u_j),
  \label{eq:nlp-generation-summac-conv}
\end{equation}
参数$\bm\psi$在带一致性标签的数据上训练。两个版本共用同一个推断模型，差别只在列聚合。

把通知第一句作为$x_1$，参加规则作为$x_2$，候选摘要两句为$y_1$：“读书会改至4月18日举行”，
$y_2$：“预计50人到场”。设教学蕴含矩阵为
\[
  \bm E=\begin{pmatrix}0.95&0.05\\0.10&0.02\end{pmatrix}.
\]
按列取最大值得$(0.95,0.05)$，平均为$0.50$，低分直接定位到人数句。若用四个等宽区间构造
直方图，第一列在最低、最高区间各有一个元素，第二列两个元素都落在最低区间，这就是
$g_{\bm\psi}$读取的分布信息。人数句蕴含分低，并不说明它与来源矛盾：它对$x_1$可能是中立
$0.92$、矛盾$0.03$，属于外在幻觉；若改成“活动仍在4月17日举行”，才会得到高矛盾分数，
属于内在幻觉。同时保存三类分数，就能区分两种幻觉。

分句也会带来问题。若来源把“读书会已经改期”和“新日期为4月18日”分成两句，摘要中的完整
陈述可能单靠任何一句都得不到充分支持；可以改用相邻句窗口重新核查，而不是把分句造成的
缺口判为幻觉。平均分还会冲淡唯一但严重的日期错误，报告时应同时给出最低支持的句子及其
对应来源。计算上，SummaC需要$MN$次推断模型调用，长文档的代价随句数乘积增长。

\subsubsection{QAFactEval的问答核验}
\label{sec:nlp-generation-qafacteval}

句对推断给出的是“能否推出”的整体判断，不说明错在哪个信息点。问答核验换一个角度：
从生成文本中的具体信息点出发提问，看来源能否给出相同答案。QAFactEval的流水线有四步：
从摘要中选出实体、名词短语等答案片段；以摘要和片段为条件生成问题；先用摘要回答这些
问题，过滤掉连摘要自己都答不回原片段的坏问题；再用来源回答保留的问题，用学习式答案
匹配器比较来源答案与原片段\scite{fabbri-etal-2022-qafacteval}
记过滤后保留的问题集为$\mathcal Q$，$m_q\in[0,1]$为问题$q$上两份答案的匹配分，来源
不可回答时$m_q=0$，则
\begin{equation}
  s_{\mathrm{QA}}=\frac{1}{|\mathcal Q|}\sum_{q\in\mathcal Q}m_q .
  \label{eq:nlp-generation-qafacteval}
\end{equation}
没有问题通过过滤时应报告“无法核验”，不能把空集合的平均当作满分。

沿用含错误人数的摘要，另加正确的第三句“校外来宾须在4月16日18:00前提交报名信息”。
选出的三个答案为“4月18日”“50人”和第三句的报名条件，对应问题为“读书会现在哪天举行”
“预计多少人到场”“校外来宾参加前需要在何时完成什么”。三题都能在摘要上答回原片段，全部
保留；到原通知中，第一题答出4月18日，第二题不可回答，第三题答出“在4月16日18:00前提交
报名信息”。若为教学采用二值匹配，得分为$(1+0+1)/3=2/3$；这不是QAFactEval真实匹配器
的输出尺度。

流水线的每一步都可能出错，核验日志要分别保留答案片段、问题、摘要答案、来源答案及出处。
人数确实是无依据新增；若问题生成成“50人何时到场”，问题本身就混入了未经证实的前提；
若来源阅读器把4月16日当成活动日期，失败发生在阅读一步。另一个局限是问题只覆盖被选中的
信息点：若选择漏掉了14:00或海棠楼201，得分高也不说明完整性已经通过。与句对推断相比，
问答核验能指出错在哪个信息点，代价是三个模型的串联误差和更多调用。

\subsection{参考重合指标}
\label{sec:nlp-generation-overlap}

忠实性核验需要来源和推断或问答模型。若手头有人工写的参考输出，更便宜的做法是统计输出与
参考共享多少单位。这类指标计算快、可复现，长期用于系统比较；它们的共同弱点是对合法改写
不宽容，对一个关键数字的错误又不够敏感。下面的例子统一小写、按空格分词、不含标点，
单参考为$r$=“the meeting is on april 18”，错日期候选$y_1$=“the meeting is on april 17”，
合法改写$y_2$=“the event takes place on april 18”。分词器、大小写、词干、参考数与聚合
层级都会改变分数，必须固定后才能比较。

\subsubsection{BLEU}
\label{sec:nlp-generation-bleu}

BLEU为机器翻译设计，衡量候选与参考共享多少词$n$元组\scite{papineni-etal-2002-bleu}
直接数候选中有多少$n$元组出现在参考里，会让重复一个常见词刷高分，所以每种$n$元组的匹配
次数要截断到它在参考中的出现次数，多参考时取各参考计数的最大值。记$\mathrm{cnt}_y(u)$为
$n$元组$u$在候选$y$中的出现次数，$R$为该句的参考集合。语料级$n$阶精确率先在句内截断，
再跨句累加：
\[
  p_n=\frac{\sum_{(y,R)}\sum_{u}\min\bigl(\mathrm{cnt}_y(u),\max_{r\in R}\mathrm{cnt}_r(u)\bigr)}
  {\sum_{(y,R)}\sum_{u}\mathrm{cnt}_y(u)} .
\]
精确率不惩罚过短的候选，只输出最有把握的几个词就能得高分，BLEU因此乘一个简短惩罚。
设语料候选总长为$\ell_y$，每句取长度最接近的参考后累加得有效参考总长$\ell_r$，则
\begin{equation}
  \mathrm{BLEU}_N=\mathrm{BP}\cdot\exp\!\Bigl(\sum_{n=1}^{N}w_n\log p_n\Bigr),\qquad
  \mathrm{BP}=\min\bigl(1,\,e^{1-\ell_r/\ell_y}\bigr),
  \label{eq:nlp-generation-bleu}
\end{equation}
其中$w_n$非负且和为一，标准设置为$N=4$、$w_n=1/4$；空候选另记零。几何平均使任何一阶
精确率为零都会让总分为零。

手算$y_1$：六个一元组匹配五个，五个二元组匹配四个，四个三元组匹配三个，三个四元组匹配
两个，长度相等使$\mathrm{BP}=1$，于是
$\mathrm{BLEU}_4=(\tfrac56\cdot\tfrac45\cdot\tfrac34\cdot\tfrac23)^{1/4}=(1/3)^{1/4}\approx0.760$。
日期错误只损失少量重合，得分仍高。$y_2$的七个一元组只匹配the、on、april、18四个，六个
二元组只匹配on april与april 18，取$N=2$时
$\mathrm{BLEU}_2=\sqrt{\tfrac47\cdot\tfrac13}\approx0.436$，远低于$y_1$的
$\mathrm{BLEU}_2=\sqrt{\tfrac56\cdot\tfrac45}\approx0.816$：保留了正确日期的改写反而得分低。
若候选只有“the meeting”，两阶精确率都是1，简短惩罚为$e^{1-6/2}=e^{-2}\approx0.135$；
它没有三元组和四元组，不平滑的$\mathrm{BLEU}_4$直接为零。若候选为“the the meeting”，
两个the只有一个能匹配，一元精确率为$2/3$而不是1。

最后一个例子说明了句级BLEU的不稳定：短句的高阶$n$元组很少，一个词的变化就能让某阶
精确率归零，必须选择某种平滑方法，不同平滑给出不同分数。BLEU原本按语料累加计数，
语料BLEU也不等于逐句BLEU的平均。分词与大小写的选择对分数影响很大，同一系统在不同
预处理下可相差数个点；为此SacreBLEU固定参考侧的预处理，并要求随分数报告一个记录全部
设置的签名\scite{post-2018-call}

\subsubsection{ROUGE-N与ROUGE-L}
\label{sec:nlp-generation-rouge}

摘要更关心参考中的要点有没有被覆盖，而不是候选里有多少词是对的。ROUGE因此以召回为
取向\scite{lin-2004-rouge}
ROUGE-N统计参考中的$n$元组有多少被候选覆盖：
\begin{equation}
  \mathrm{ROUGE\text{-}N}=\frac{\sum_{r\in R}\sum_{u\in r}\min\bigl(\mathrm{cnt}_y(u),\mathrm{cnt}_r(u)\bigr)}
  {\sum_{r\in R}\sum_{u\in r}\mathrm{cnt}_r(u)},
  \label{eq:nlp-generation-rouge-n}
\end{equation}
分母是参考的$n$元组总数，而BLEU的分母是候选的$n$元组总数，这是两者取向不同的来源。
多参考时，原定义在分子分母上对所有参考累加，参考中重复出现的要点权重更大；常用实现也
有对每个参考分别计算、再取最大值的做法，报告时需写明。

ROUGE-N要求词连续匹配，换一个词序就丢分。ROUGE-L改用\term{最长公共子序列}（longest
common subsequence, LCS），即两个序列中保持先后顺序、但不要求相邻的最长共同词串。设参考
长度为$m$，候选长度为$n$，LCS长度为$L$，则
\begin{equation}
  R_{\mathrm{L}}=\frac{L}{m},\qquad P_{\mathrm{L}}=\frac{L}{n},\qquad
  F_{\mathrm{L}}=\frac{(1+\beta^2)R_{\mathrm{L}}P_{\mathrm{L}}}{R_{\mathrm{L}}+\beta^2P_{\mathrm{L}}},
  \label{eq:nlp-generation-rouge-l}
\end{equation}
$\beta$越大越偏向召回，原文在摘要评测中取很大的$\beta$，常用工具包则报告$\beta=1$的F值。
多句摘要还要说明是逐句计算后平均，还是在整篇上取各参考句与候选句的LCS并集。

对$y_1$，参考的五个二元组被覆盖四个，ROUGE-1与ROUGE-2分别为$5/6$与$4/5$；LCS为
“the meeting is on april”，$L=5$，$F_{\mathrm L}=5/6\approx0.833$。对$y_2$，ROUGE-1为
$4/6\approx0.667$，ROUGE-2为$2/5$；LCS为“the on april 18”，$L=4$，
$F_{\mathrm L}=2L/(m+n)=8/13\approx0.615$。与BLEU一样，错日期候选在三项上都高于正确改写。

\subsubsection{chrF与METEOR}
\label{sec:nlp-generation-chrf-meteor}

词级匹配依赖分词：形态丰富的语言里同一个词有许多变体，中文还要先确定词边界。chrF改看
字符$n$元组\scite{popovic-2015-chrf}
对$n=1,\ldots,N$分别计算字符$n$元组的精确率$P_n$与召回率$R_n$，先取算术平均
$\bar P=N^{-1}\sum_nP_n$、$\bar R=N^{-1}\sum_nR_n$，再计算
\[
  \mathrm{chrF}_\beta=\frac{(1+\beta^2)\bar P\bar R}{\beta^2\bar P+\bar R},
\]
$\beta>1$使召回的权重更大。取参考“4月18日”、候选“4月17日”，按字符切分：五个一元字符
匹配四个，四个二元字符中只有“4月”“月1”匹配，所以$P_1=R_1=4/5$，$P_2=R_2=1/2$。教学取
$N=2,\beta=2$，得$\bar P=\bar R=0.65$，$\mathrm{chrF}=0.65$；工具包的默认设置用更高的阶数。
字符匹配减轻了分词问题，同样不理解数字的含义。

前面的指标都把starts与begins当作完全不同的词。METEOR在词级上逐步放宽匹配
\scite{banerjee-lavie-2005-meteor}
对齐分三个阶段进行：先匹配完全相同的词，再对剩余词匹配词干相同者，最后用同义词表匹配
同义词；每个词至多参与一次对齐，若有多种对齐方式，选交叉最少的一种。设匹配词数为$m_a$，
一元精确率$P=m_a/n$、召回率$R=m_a/m$。METEOR用召回权重为精确率九倍的调和平均，并按
匹配词在两句中被切成多少个连续块施加碎片惩罚：
\begin{equation}
  F_{\mathrm{mean}}=\frac{10PR}{R+9P},\qquad
  \mathrm{Pen}=0.5\Bigl(\frac{\mathit{ch}}{m_a}\Bigr)^3,\qquad
  \mathrm{METEOR}=F_{\mathrm{mean}}(1-\mathrm{Pen}),
  \label{eq:nlp-generation-meteor}
\end{equation}
其中$\mathit{ch}$是块数：一个块是一串在候选与参考中都相邻、且顺序一致的匹配词。所有匹配
连成一块时惩罚最小，词序越乱块越多，惩罚越大，最多扣去一半。

取参考“the meeting begins today”、候选“the meeting starts today”，并假定同义词表把starts
与begins归入同一组。精确阶段匹配the、meeting、today，同义词阶段匹配starts与begins，
$m_a=4$，$P=R=1$，四个词连成一块，$\mathrm{Pen}=0.5\times(1/4)^3\approx0.0078$，
$\mathrm{METEOR}\approx0.992$；同一对句子的一元BLEU精确率只有$3/4$。若候选改为
“today the meeting starts”，匹配仍是四个词，但today单独成块，$\mathit{ch}=2$，
$\mathrm{Pen}=0.0625$，得分降为$0.9375$。回到$y_1$，它与$r$匹配五个词、连成一块，
$\mathrm{METEOR}=\tfrac56(1-0.5\times0.2^3)\approx0.830$，日期错误仍只扣一个词的分。
METEOR的词干器与同义词表依赖语言资源，没有这类资源的语言只能退化为精确匹配。

\subsection{编辑任务指标}
\label{sec:nlp-generation-edit-metrics}

纠错与简化的输出大部分照抄来源。取来源$x$=“the previously announced meeting starts today”、
参考$r$=“the meeting begins today”。原样输出$x$的ROUGE-1为$3/4$，删去背景的候选
“the meeting starts today”同样是$3/4$：只看输出与参考的重合，分辨不出“做了正确的修改”
和“什么也没改”。编辑指标因此同时读取来源、输出和参考，比较的是修改本身：要么分别统计
增加、保留与删除了什么，要么先把输出与来源对齐成编辑，再与参考编辑比对。

\subsubsection{SARI}
\label{sec:nlp-generation-sari}

SARI为文本简化设计，分别评价增加、保留和删除三种操作\scite{xu-etal-2016-optimizing}
在单参考、集合计数的教学口径下，记来源、候选、参考的$n$元组集合为$X_n,Y_n,R_n$。增加
比较$Y_n\setminus X_n$与$R_n\setminus X_n$，保留比较$Y_n\cap X_n$与$R_n\cap X_n$，删除比较
$X_n\setminus Y_n$与$X_n\setminus R_n$。增加与保留用F1，精确率的分母是候选的操作数，召回
的分母是参考的操作数；删除只用精确率，避免参考删得多时奖励不加选择的删除。一至四阶
平均：
\begin{equation}
  \mathrm{SARI}=\frac{1}{4}\sum_{n=1}^{4}
  \frac{F_{\mathrm{add},n}+F_{\mathrm{keep},n}+P_{\mathrm{del},n}}{3}.
  \label{eq:nlp-generation-sari}
\end{equation}
分母为零的项记零，这沿用原实现的约定；多参考时还按参考中的频次加权，不能把参考简单拼接。

取候选$y$=“the meeting starts today”，任务允许略去“此前已宣布”的背景。一阶上候选没有新增
词，参考新增begins，增加F值为0；候选保留四词，参考要求保留其中三词，
$P_{\mathrm{keep},1}=3/4$，$R_{\mathrm{keep},1}=1$，F值为$6/7$；候选删除previously与announced，
都属于参考的删除，删除精确率为1。二阶候选只新增the meeting，参考新增the meeting、meeting
begins、begins today，增加精确率1、召回$1/3$、F值$1/2$；候选保留的meeting starts与starts
today都不在参考保留集合中，保留F值为0。各阶结果为
\begin{center}
  \small
  \begin{tabular}{@{}crrrccc@{}}
    \toprule
    阶 & 增加数 & 保留数 & 删除数 & 增加F1 & 保留F1 & 删除P\\
    \midrule
    1 & 0/1/0 & 4/3/3 & 2/3/2 & 0 & $6/7$ & 1\\
    2 & 1/3/1 & 2/0/0 & 3/5/3 & $1/2$ & 0 & 1\\
    3 & 1/2/0 & 1/0/0 & 3/4/3 & 0 & 0 & 1\\
    4 & 1/1/0 & 0/0/0 & 3/3/3 & 0 & 0 & 1\\
    \bottomrule
  \end{tabular}
\end{center}
每个操作数单元按“候选操作数／参考操作数／匹配数”排列。最终
$\mathrm{SARI}=(13/21+1/2+1/3+1/3)/4=25/56\approx0.446$。

原样输出$x$时，候选没有增加也没有删除，两项各阶都记零；一阶保留六词，参考保留其中三词，
保留F值为$2/3$；二阶以上参考保留集合为空，保留F值为零。于是
$\mathrm{SARI}=(2/9)/4=1/18\approx0.056$，远低于做了删除的候选，这正是编辑指标要区分的
情形。starts与begins意思相近，参考的操作仍偏好后者，说明操作匹配不等于语义判断；背景是否
允许删除也要由任务决定，不能反过来由SARI高低决定。

\subsubsection{M2与ERRANT}
\label{sec:nlp-generation-m2-errant}

纠错要求的是最小必要修改，评价时关心每一处编辑是否正确。M2评分器把系统输出与来源对齐
成编辑集合，与人工标注的参考编辑集合比对\scite{dahlmeier-ng-2012-better}
同一变化可以切成“替换一个短语”或“几处替换加插入”，切法不同会让计数不同；M2在来源到
输出的编辑格中保留所有最短对齐，加入跨越多个基本操作的短语编辑边，再选出与参考标注
匹配最多的编辑序列。对句子$s$，记系统编辑集合为$e_s$、参考编辑集合为$g_s$，跨句累加：
\begin{equation}
  P=\frac{\sum_s|e_s\cap g_s|}{\sum_s|e_s|},\qquad
  R=\frac{\sum_s|e_s\cap g_s|}{\sum_s|g_s|},\qquad
  F_\beta=\frac{(1+\beta^2)PR}{\beta^2P+R}.
  \label{eq:nlp-generation-m2}
\end{equation}
有多名标注者时，每句选使该句得分最高的那份参考。纠错任务通常取$\beta=0.5$，精确率的
权重更大，因为把正确的句子改错比漏改一处更伤害用户。

取来源“He go to school yesterday with book .”，零起始词坐标下参考编辑为
$e_1:[1,2)\mapsto\text{went}$和$e_2:[6,6)\mapsto\text{a}$。候选甲执行两项，匹配2、预测2、
参考2，$P=R=1$；候选乙只把go改为went，$P=1,R=1/2$；候选丙在甲上还把He改成The boy，
新增了来源未支持的人物描述，$P=2/3,R=1$。按$F_{0.5}$，乙得$5/6$，丙得$5/7$；按F1，
乙得$2/3$，丙得$4/5$，排序恰好相反。原始M2论文报告F1，$F_{0.5}$是后来的纠错评测指定的
协议，比较时必须写明。

M2只告诉我们改对了多少，不告诉我们哪类错误没改。ERRANT从平行句子自动抽取编辑并分类：
它用结合词元与词性信息的Damerau--Levenshtein对齐找出编辑，按语言规则合并相邻编辑，再依据
来源与目标词的词性、词形等特征把编辑归入错误类别\scite{bryant-etal-2017-automatic}
这样，系统输出与参考都能自动转成带类别的编辑，评分不再依赖人工标注的编辑边界。在上例中，
go到went可归入动词时态替换，缺a归入限定词缺失；按类别计数，乙的时态召回为1、限定词召回
为0，比总分更直接地指出遗漏。ERRANT的对齐合并规则与类别体系是为英语设计的，中文没有
词形变化与冠词，需要另定分词、编辑边界与错误类别。

\subsection{学习型指标}
\label{sec:nlp-generation-learned-metrics}

重合指标对合法改写不宽容，编辑指标也只认参考中出现的操作。学习型指标用预训练模型比较
语义：BERTScore不需训练，直接匹配上下文向量；BLEURT与COMET在人工评分上训练回归器，
直接预测人会打多少分；无参考质量估计再去掉参考，只用来源与候选。它们的共同代价是分数
依赖模型版本、语言覆盖与训练标注的领域，换了模型或领域就要重新检查与人工判断的一致性。

\subsubsection{BERTScore语义匹配}
\label{sec:nlp-generation-bertscore}

BERTScore用上下文词元向量代替字面匹配\scite{nlp-generation-zhang2020bertscore}
用同一预训练编码器分别编码候选与参考，取某一层的词元向量并归一化为单位长度，记候选
向量为$\bm h_i$（$i=1,\ldots,|y|$），参考向量为$\bm r_j$（$j=1,\ldots,|r|$）。让候选中
每个词元找最相似的参考词元，得到精确率；再从参考向候选反向匹配，得到召回：
\begin{equation}
  P_{\mathrm{B}}=\frac{1}{|y|}\sum_i\max_j\bm h_i^\top\bm r_j,\qquad
  R_{\mathrm{B}}=\frac{1}{|r|}\sum_j\max_i\bm h_i^\top\bm r_j,
  \label{eq:nlp-generation-bertscore}
\end{equation}
二者的调和平均为$F_{\mathrm B}$。这是双向最大相似匹配，不是一对一对齐，多个候选词元可以
匹配同一个参考词元。常见虚词贡献的相似度高却信息少，可选的逆文档频率加权把平均改为
加权平均，例如
$R_{\mathrm B}=\sum_j\mathrm{idf}(r_j)\max_i\bm h_i^\top\bm r_j/\sum_j\mathrm{idf}(r_j)$。
上下文向量的余弦普遍偏高，不相关的句子也常得到0.8以上的分数，可读性差；基线重标度
用随机句对上的平均分$b$把分数线性变换为$(F_{\mathrm B}-b)/(1-b)$，只改变尺度，不改变排序。

令参考词元为“甲／帮助／乙”，候选为“甲／协助／乙”。用四维标准单位向量构造教学表示：参考
依次取$\bm e_1,\bm e_2,\bm e_3$，候选取$\bm e_1,0.8\bm e_2+0.6\bm e_4,\bm e_3$。中间向量的
平方范数为$0.8^2+0.6^2=1$，余弦矩阵的对角线为$(1,0.8,1)$，其余为零。双向最大值相同，
$P_{\mathrm B}=R_{\mathrm B}=F_{\mathrm B}=2.8/3=14/15\approx0.933$，而一元词重合只有$2/3$：
近义替换得到了部分分数。若教学基线$b=0.8$，重标度后为$(14/15-0.8)/0.2=2/3$。再把候选改成
“乙／帮助／甲”，若各词的表示仍分别靠近$\bm e_3,\bm e_2,\bm e_1$，矩阵只是交换了行，双向
最大值全部为1，帮助者与被帮助者却已经互换。真实的上下文编码器会让“乙”在不同位置上的
向量有所不同，可能降低分数，但匹配公式本身不检查角色，数字与否定同样要单独核对。

\subsubsection{BLEURT与COMET}
\label{sec:nlp-generation-comet}

BERTScore的匹配规则是人为设计的，并不知道人工评价者在意哪些差异。BLEURT与COMET改为
直接学习人工评分。

BLEURT把参考与候选拼成一个序列送入BERT，取$[\mathrm{CLS}]$位置的向量$\bm v$，用线性层
预测分数$\widehat q=\bm w^\top\bm v+b$，以均方误差在人工评分上训练\scite{sellam-etal-2020-bleurt}
人工评分数据规模有限，且集中在少数语言与系统，直接微调容易在质量分布变化时失效。
BLEURT因此在微调前加了一个预训练阶段：从维基百科句子出发，用掩码语言模型填空、回译和随机
删词构造数百万对“参考—扰动句”，再为每对句子计算多种自动信号作为回归或分类目标，包括
BLEU、ROUGE、BERTScore、回译似然、蕴含分类器给出的蕴含/中立/矛盾概率，以及该扰动是否来自
回译。多任务预训练目标是各信号损失的加权和
\begin{equation}
  \mathcal L_{\mathrm{pre}}=\frac{1}{K}\sum_{k=1}^{K}\gamma_k\,\ell_k\bigl(\tau_k,\widehat\tau_k\bigr),
  \label{eq:nlp-generation-bleurt}
\end{equation}
$\tau_k$是第$k$个信号的值，$\widehat\tau_k$是模型在该信号上的预测，$\gamma_k$是权重。
预训练让模型先学会“什么样的扰动改变了意思、改变了多少”，再用少量人工评分校准到人的
尺度。BLEURT只读参考与候选，不读源文。

COMET面向翻译，把源文也作为输入\scite{rei-etal-2020-comet}
它的估计器结构是：用跨语言编码器XLM-R分别编码源句$x$、候选$y$与参考$r$，对各层输出做
可学习的加权混合，再对词元做平均池化，得到句向量$\bm s,\bm h,\bm r$；把候选与另两者的
关系写成特征
\begin{equation}
  \bm f=[\bm h;\bm r;\bm h\odot\bm s;\bm h\odot\bm r;|\bm h-\bm s|;|\bm h-\bm r|],
  \label{eq:nlp-generation-comet}
\end{equation}
交给前馈网络回归人工质量分，以均方误差$n^{-1}\sum_i(\widehat q_i-q_i)^2$训练。逐元素乘积
衡量两句在各维上的一致程度，绝对差衡量偏离程度；源向量不单独作为一块，但通过乘积和
差值参与预测。与BLEURT的交叉编码不同，COMET分别编码三句，候选与参考之间没有逐词元
注意力，换来的是可以预先缓存源与参考的向量。若教学标签为$(1.0,0.2)$，预测为$(0.8,0.5)$，
这一批的损失为$(0.04+0.09)/2=0.065$，仅演示回归目标。

读取源文的好处在参考不完全时显现。源文补充记录为“资料袋共12份”，正确候选为“There are 12
information packs”，错误候选把12改成21；若参考恰好写成“Twelve packs are provided”，两条
候选与它的词重合几乎一样低，而COMET可以从源向量与候选向量的差别中察觉第二条的数字不符。
这只是有机会察觉，句向量并不保证对数字交换敏感，关键数字仍要按上一节的方法直接核验。
系统级分数是固定测试集上句级预测的平均，整体均值高不说明每个关键句都合格。

\subsubsection{无参考质量估计}
\label{sec:nlp-generation-qe}

上线后的翻译系统没有参考，BLEU、BLEURT和参考型COMET都无从计算。\term{无参考质量估计}
（reference-free quality estimation, QE）只输入源文与候选，预测候选的质量。它仍需人工质量
标注来训练，只是推断时不再需要参考；把参考型模型的参考字段删掉并不能得到同样有效的
QE模型。QE的输出有两种粒度：句级回归一个质量分；词级或跨度级为每个词元预测OK或BAD标签，
或直接预测错误跨度及其严重程度，后者能指出该改哪里。COMET框架下的CometKiwi把源文与候选
拼接后联合编码，同时输出句级分数与词级标签\scite{rei-etal-2022-cometkiwi}
官方发布的模型中，\code{wmt22-comet-da}需要源文、候选与参考，\code{wmt22-cometkiwi-da}只需
源文与候选，两者的输入、训练标签与分数尺度不同，需分别记录。

QE有三种主要用途。其一是候选重排：对同一源句生成多个候选，按QE分数选优；也可以把参考型COMET当作效用函数，
以其他候选充当伪参考做最小贝叶斯风险解码，两种做法及其代价见第~\ref{sec:nlp-generation-qe-mbr}小节。
其二是过滤：把低分译文送人工复核，或从回译数据中剔除低质量句对。其三是无参考监控：
跟踪线上译文质量随时间的变化。用同一个指标既选候选又报告质量，会把选择偏差算作质量
提升，报告时应换用另一个指标或人工评价。QE分数不是“译文正确的概率”，在训练语向之外、
陌生领域、术语密集或需要篇章上下文的输入上，校准与错误分布都要另行检查；数字与角色
绑定仍需独立核验。

\subsection{人工评价与指标元评价}
\label{sec:nlp-generation-human-meta}

前面的指标都声称在近似人工判断，人工判断本身也需要规定：评价者看什么材料、按什么类别
记错、错误多重算多少分。元评价则反过来，以人工判断为参照衡量一个自动指标可信到什么
程度、可以在什么粒度上使用。

\subsubsection{MQM错误标注与人工比较}
\label{sec:nlp-generation-mqm}

只给译文打一个总分，既难以复核，也不告诉译者怎样修订。\term{多维质量指标}（Multidimensional
Quality Metrics, MQM）是按错误类别与严重程度组织人工标注的框架：评价者在译文中标出每处
错误的位置，归入准确性（误译、遗漏、增译）、流畅性、术语、风格等类别，并标注严重程度。
片段的罚分是各错误按严重程度加权后的和：
\begin{equation}
  \mathrm{MQM}(y)=\sum_{e\in\mathcal E(y)}w_{\mathrm{sev}(e)},
  \label{eq:nlp-generation-mqm}
\end{equation}
$\mathcal E(y)$是标出的错误集合，罚分越高质量越差。专家MQM标注的研究使用主要错误记5分、
次要错误记1分的权重，并发现非专家众包评分可能把人工译文排在机器译文之后，而专家MQM能
纠正这一排序\scite{nlp-intro-freitag2021}
具体研究还要规定类别子集、标注单位、评价者能否看到篇章上下文、复评方式与汇总规则，
这些设置不是一个缩写就能确定的。

固定源句“校外来宾须在4月16日18:00前提交报名信息”。候选甲为“Guests are welcome to join”，
候选乙为“External guests must, before 18:00 on April 16, submit registration information”。
甲简短流畅，却没有告诉校外来宾必须完成什么；乙较生硬，但保留了角色、动作与截止条件。
只问“更喜欢读哪句”与问“能否据此正确参加活动”，可能得到相反的结论。评价界面应提供
同一份源文与用途说明，随机交换候选的左右位置，分别收集偏好、错误位置与用途是否满足。
按上述权重，甲的关键条件遗漏记为一处主要错误，罚分为5；乙无已识别错误，罚分为0。
若用途规定任何关键条件遗漏都不合格，还应同时记录甲未通过硬验收，不能让语料平均罚分
掩盖这次失败。

局部错误标注与总体差异的证据也不同。一个教学的配对文档罚分差为$d=(5,0,5,0)$，正值表示
甲更差。若暂把四份文档看作独立样本，平均差为2.5，样本标准差约2.887，标准误约1.443，
自由度3的$t_{0.975}=3.182$给出区间$2.5\pm3.182\times1.443\approx[-2.09,7.09]$，包含零。
样本这么少，不能断言总体上甲更差；但这不妨碍对已核实的那一处遗漏作出明确修订。区间
估计与重采样的一般做法见第~\ref{sec:nlp-eval-uncertainty}小节。

\subsubsection{指标与人工判断的相关性}
\label{sec:nlp-generation-meta-eval}

一个指标好不好，要看它的排序与人工判断有多一致，这类研究称为\term{元评价}
（meta-evaluation）。一致性分两个粒度衡量。系统级比较各系统在整个测试集上的平均指标分
与平均人工分，回答“能否用它选出更好的系统”；段级比较同一来源句的不同译文，回答“能否
用它判断单句好坏或做候选重排”。两种粒度上的结论可以不同。

系统级常用Pearson相关系数。设四个系统的教学人工分依次为72、75、78、80，指标分依次为
30、33、32、35，两组均值为76.25与32.5，离差乘积之和为18.5，离差平方和分别为36.75与13，
相关系数$18.5/\sqrt{36.75\times13}\approx0.846$。相关不低，但指标把第二、三个系统的顺序
排反了。若再加入一个明显很差的系统，人工分40、指标分10，相关系数升到约0.995：一个离群
系统就能让相关显得近乎完美，却没有让指标更会区分前四个接近的系统
\scite{mathur-etal-2020-tangled}
系统级相关因此应在去掉离群系统后复核，或改报两两系统排序的一致比例。

段级常用\term{Kendall秩相关系数}（Kendall rank correlation coefficient），它只看成对排序
是否一致。对$n$个对象的全部$\binom n2$对，若人工与指标对某对的排序相同记为一致，相反
记为不一致，则
\begin{equation}
  \tau=\frac{\#\text{一致}-\#\text{不一致}}{\#\text{一致}+\#\text{不一致}} .
  \label{eq:nlp-generation-kendall}
\end{equation}
无平局时分母就是$\binom n2$。设同一来源句的四个系统译文，人工分为$(60,75,40,90)$，指标分为
$(0.52,0.48,0.30,0.81)$。六对中只有第1、2条译文的顺序相反，$\tau=(5-1)/6=2/3$。人工直接
打分本身有噪声，分差很小的两条译文谁好谁坏并不可靠；WMT曾只保留人工分差超过25分的译文
对来计算$\tau$\scite{ma-etal-2019-results}
本例按此规则只剩第1与4、第2与3、第3与4三对，全部一致，$\tau=1$。阈值与平局处理都会
改变$\tau$，报告时必须写明。

WMT指标评测任务把这些做法组织成年度比较：参评者提交自己的指标在当年翻译任务全部系统
输出上的分数，组织者收集人工判断，按语向分别计算系统级与段级相关，并用重采样或成对
显著性检验判断指标之间的差异是否可靠。人工判断起初主要来自众包的直接评估打分，自2021年
起改以专家MQM标注为主要参照\scite{freitag-etal-2021-results}
这些评测反复显示，学习型指标与人工判断的相关普遍高于BLEU，但在不同语向、领域与高质量
系统之间的排序上仍不稳定；在一个语向或领域上的元评价结论，不能直接搬到另一个。

\subsection{大语言模型评判}
\label{sec:nlp-generation-llm-judge}

参考与人工评分都难以覆盖的维度，例如摘要是否连贯、对话回复是否有趣，可以写成文字准则，
让大语言模型按准则打分。第~\ref{sec:nlp-eval-judge-protocol}小节已给出\term{大模型评判}的
三种形式、位置与长度偏差的修正，以及用一致率和$\kappa$校准评判模型的方法；这里只讨论
生成任务特有的做法。

G-Eval把评判拆成三步\scite{liu-etal-2023-geval}
第一步写出任务说明与一个维度的评价准则，例如摘要的一致性定义为“摘要中的事实都能在原文中
找到支持，按1至5分打分”。第二步让模型根据准则自行生成一串评价步骤，例如“先通读原文，再
逐条找出摘要中的事实陈述，检查每条是否与原文一致”，即思维链形式的评价程序。第三步把
原文、摘要与这些步骤填进评分表单，让模型输出分数。直接取模型输出的整数分，大量样本会
挤在同一个分数上，难以区分；G-Eval改用分数词元的概率加权求期望：
\begin{equation}
  \mathrm{score}=\sum_{s=1}^{5}s\cdot p(s\mid\text{准则},\text{步骤},x,y).
  \label{eq:nlp-generation-geval}
\end{equation}
接口不提供词元概率时，原文以多次采样的频率近似$p(s)$。设两条摘要的教学分数分布为
$(0.05,0.10,0.20,0.45,0.20)$与$(0,0.05,0.10,0.50,0.35)$，取最可能的整数分都是4，期望分则为
3.65与4.15，后者更好的信号得以保留。

翻译评价上也有类似做法：直接让模型按0至100分给译文打分，或让模型按MQM的类别与严重程度
标出错误跨度，再按式\eqref{eq:nlp-generation-mqm}换算成分数
\scite{kocmi-federmann-2023-large,kocmi-federmann-2023-gemba}
后者的输出可以逐条复核，与人工MQM标注的接口一致。按维度分别评判也是生成任务的常见
做法：摘要按一致性、连贯性、流畅性与相关性四个维度分别打分，与SummEval的人工标注维度对应，
便于逐维做元评价。

评判模型在生成任务上有一种额外偏差：它倾向于给大语言模型写出的文本更高的分数，即使人工
评价者认为人写的摘要更好；G-Eval的作者也报告了这一倾向。评判与被评系统同源时，这种偏好
会放大。因此，评判的分数与上一小节的指标一样，要在本任务、本语言的人工标注子集上做
元评价，评判提示、模型版本与解码设置任一变化都要重新校准。忠实性等有来源可查的维度，
仍应优先用第~\ref{sec:nlp-generation-faithfulness}小节的核验程序，评判模型适合作补充。

表~\ref{tab:nlp-generation-metrics}汇总本节各指标的测量对象、所需输入与主要失效情形。

\begin{table*}[htbp]
  \centering
  \small
  \renewcommand{\tabularxcolumn}[1]{>{\raggedright\arraybackslash}p{#1}}
  \begin{tabularx}{\linewidth}{@{}>{\raggedright\arraybackslash}p{24mm}>{\raggedright\arraybackslash}p{30mm}XX@{}}
    \toprule
    指标 & 所需输入 & 测量对象 & 主要失效情形\\
    \midrule
    SummaC & 来源、输出 & 每个输出句能否由来源推出 & 跨句证据被分句切断；平均冲淡单个严重错误\\
    QAFactEval & 来源、输出 & 选出的信息点在来源中能否答出相同答案 & 只覆盖被选中的信息点；三模型误差串联\\
    BLEU & 输出、参考 & 截断$n$元精确率与长度 & 合法改写低分；句级不稳定；关键数字错误扣分少\\
    ROUGE & 输出、参考 & 参考$n$元或LCS的覆盖 & 同上；多参考与汇总方式需写明\\
    chrF、METEOR & 输出、参考 & 字符重合；含词干与同义词的对齐与词序 & 依赖语言资源；仍不理解数字与角色\\
    SARI & 来源、输出、参考 & 增加、保留、删除操作 & 近义操作不计；删除是否允许由任务决定\\
    M2、ERRANT & 来源、输出、参考编辑 & 编辑的精确率与召回，按类别细分 & 编辑边界与类别依赖语言\\
    BERTScore & 输出、参考 & 上下文向量的双向最大匹配 & 角色互换可能高分；依赖编码器与层\\
    BLEURT、COMET & 输出、参考；COMET另需来源 & 预测的人工质量分 & 领域与语向外的校准；对数字交换不保证敏感\\
    质量估计 & 来源、输出 & 无参考的质量分或错误跨度 & 用于选择后再报告会高估；校准随领域变化\\
    MQM & 来源、输出、人工 & 带类别与严重程度的错误 & 成本高；类别与权重需按用途规定\\
    大模型评判 & 准则、来源、输出 & 按文字准则的打分或错误标注 & 偏好模型生成文本；须在本任务上元评价\\
    \bottomrule
  \end{tabularx}
  \caption{生成评价指标的输入与测量对象。各指标测量不同的事件，分数不能相互换算；关键事实与
  硬约束应由来源核验或值绑定检查单独判定，再用这些指标观察表达质量。}
  \label{tab:nlp-generation-metrics}
\end{table*}

\Needspace{8\baselineskip}


\section*{文献说明}
\label{sec:nlp-generation-literature}

下面按本章各节集中给出方法的原始文献和数据集入口。正文中每个方法至多保留一个支撑其机制的对照数字，均注明出自原文；其余数值是本书的教学例子，不表示复现了论文实验。

\begin{fullwidth}
\small
\textbf{条件生成的共同问题。}
注意力编码器—解码器：\fullcite{bahdanau2015}
\par\smallskip
指针生成网络与覆盖机制：\fullcite{nlp-intro-see2017}
\par\smallskip
序列级训练与暴露偏差：\fullcite{nlp-generation-ranzato2016mixer}
\par\smallskip
最小风险训练：\fullcite{nlp-generation-shen2016mrt}
\par\medskip
\textbf{数据到文本与开放写作。}
经典生成流水线：\fullcite{nlp-generation-reiter2000building}
\par\smallskip
RotoWire数据集与抽取式评价：\fullcite{nlp-generation-wiseman2017}
\par\smallskip
WebNLG提供RDF三元组到文本的任务入口：\fullcite{nlp-intro-gardent2017}
\par\smallskip
ToTTo提供带高亮单元格的受控表格到文本任务：\fullcite{nlp-generation-parikh2020totto}
\par\smallskip
内容选择门与指针规划：\fullcite{nlp-generation-puduppully2019}
\par\smallskip
WritingPrompts数据集与分层故事生成：\fullcite{nlp-inference-fan2018story}
\par\smallskip
即插即用语言模型：\fullcite{nlp-generation-dathathri2020pplm}
\par\smallskip
前缀分类器的软属性控制：\fullcite{yang-klein-2021-fudge}
\par\smallskip
长篇故事的递归重提示与修订：\fullcite{nlp-generation-yang2022re3}
\par\medskip
\textbf{文本编辑。}
Delete--Retrieve--Generate的属性删除、检索与生成设置：\fullcite{li-etal-2018-delete}
\par\smallskip
BEA-2019纠错共享任务及其数据：\fullcite{nlp-generation-bryant2019bea}
\par\smallskip
JFLEG提供英文流畅性纠错材料：\fullcite{nlp-intro-napoles2017}
\par\smallskip
Newsela简化语料及其对简化研究数据问题的分析：\fullcite{nlp-generation-xu2015newsela}
\par\smallskip
WikiLarge简化语料随DRESS一起发布：\fullcite{nlp-generation-zhang2017dress}
\par\smallskip
ASSET提供含多种改写操作的简化参考：\fullcite{nlp-intro-alvamanchego2020}
\par\medskip
\textbf{摘要与信息综合。}
BART的去噪预训练：\fullcite{lewis-etal-2020-bart}
\par\smallskip
CNN/DailyMail新闻摘要语料的来源：\fullcite{nlp-generation-hermann2015cnndm}
\par\smallskip
XSum对应高度压缩的新闻摘要任务：\fullcite{nlp-intro-narayan2018}
\par\smallskip
SummEval连接自动分数与多维人工评价：\fullcite{fabbri-etal-2021-summeval}
\par\smallskip
FActScore评估原子事实的来源支持：\fullcite{min-etal-2023-factscore}
\par\medskip
\textbf{机器翻译。}
短语统计翻译的模型与短语抽取：\fullcite{nlp-intro-koehn2003}
\par\smallskip
STACL提供前缀可见条件与等待策略：\fullcite{nlp-intro-ma2019stacl}
\par\smallskip
NLLB连接多语数据、模型与语言覆盖评价：\fullcite{nlp-intro-nllb2022}
\par\smallskip
WMT24是固定任务条件下的系统比较入口：\fullcite{nlp-intro-kocmi2024}
\par\smallskip
DiscoX提供中英专业篇章评价材料：\fullcite{nlp-intro-zhao2026discox}
\par\medskip
\textbf{生成结果的评价。}
BLEU的截断计数与长度惩罚：\fullcite{papineni-etal-2002-bleu}
\par\smallskip
BERTScore的词元表示与匹配定义：\fullcite{nlp-generation-zhang2020bertscore}
\par\smallskip
MQM错误标注、专家评价及上下文影响：\fullcite{nlp-intro-freitag2021}
\par\smallskip
COMET各版本的输入格式（参考型与无参考型）以官方文档为准：\url{https://github.com/Unbabel/COMET}。
\end{fullwidth}

\section{本章小结}
\label{sec:nlp-generation-summary}

章首列出的七项共同挑战，在本章中各由一组机制承担。读取与保留条件靠交叉注意力和复制：注意力让解码器每步按需读取任意长度的条件，指针生成网络让名称和数字可以从输入原样取出；训练与生成的不一致则由计划采样、序列级训练和最小风险训练逐步缓解。内容选择与顺序靠规划：经典流水线用规则完成，内容选择门与指针规划把它变成可学习、可检查的中间结果，分层故事生成和递归规划把同一思路推广到开放写作与长篇。输入与输出高度重合时，编辑标签只预测改动的位置；缺少平行数据时，删除—检索—生成、回译与挖掘复述从单语语料中构造训练信号。长输入由长文档摘要结构、篇章翻译和分层规划处理。术语、格式和值绑定这类硬要求由约束解码与核验负责，不能指望软属性控制或采样参数替代。

这些机制没有一项单独保证忠实。复制能写出正确的名称，却不阻止它被绑定到错误的对象；规划能检查选了哪些内容，却不检查句子是否忠实实现了计划；属性控制改变措辞，可能同时提高错误内容的概率。参考重合、编辑匹配、学习型指标、人工评价和大语言模型评判各自只测量质量的一个方面，因此结果仍需按第~\ref{sec:nlp-generation-recovery}小节的流程恢复成可检查的对象，逐项核对实体绑定、否定、日期、数量和来源支持，无解或预算不足时明确返回失败。

本章的条件都在生成之前给定。下一章转向问答与对话：回答所需的证据往往要先检索，交流目标要在多轮中维护和澄清。本章的条件生成、复制与核验机制在那里继续使用，证据获取与交互状态则需要另外建模。
